\documentclass[11pt]{article}

\usepackage[utf8]{inputenc}
\usepackage[T1]{fontenc}
\usepackage[margin=1in]{geometry}
\usepackage{amsmath,amssymb}
\usepackage{graphicx}
\usepackage{longtable,booktabs}
\usepackage{microtype}
% hyperref goes last; it makes \cite clickable through to the bibliography.
\usepackage[colorlinks=true,linkcolor=blue,citecolor=teal,urlcolor=magenta]{hyperref}

% pandoc emits \tightlist but expects the preamble to define it.
\providecommand{\tightlist}{%
  \setlength{\itemsep}{0pt}\setlength{\parskip}{0pt}}

% If the compile times out on Overleaf, comment out \tableofcontents below first.
\title{A Comprehensive Survey on In-Context Learning}
\author{Generated by AutoSurvey}
\date{}

\begin{document}
\maketitle
\tableofcontents
\newpage

\section{Introduction to In-Context Learning}

\subsection{What is In-Context Learning?}

In-context learning (ICL) refers to the remarkable ability of large language models (LLMs) to adapt to and perform a wide range of tasks using only a few demonstration examples provided as part of the input, without any model parameter updates \cite{ref1,ref2}. This emerging paradigm has gained significant attention in the field of natural language processing (NLP) and beyond, as it offers a promising alternative to the traditional fine-tuning approach, which requires extensive task-specific training data and computational resources \cite{ref3}.

The emergence of in-context learning in LLMs can be traced back to the scaling up of language models, such as GPT-3 \cite{ref4}. As these models were trained on vast amounts of diverse text data, they were able to acquire a deep understanding of language and the ability to leverage contextual information to adapt to new tasks \cite{ref5}. This property is particularly remarkable given that LLMs are not explicitly trained on in-context learning; rather, it emerges as a result of the models' ability to capture and recombine the patterns and structures present in the pretraining data \cite{ref6}.

The significance of in-context learning lies in its potential to revolutionize the way we approach natural language tasks. Unlike traditional machine learning approaches, which often require large amounts of labeled data and extensive fine-tuning, in-context learning enables LLMs to quickly adapt to new tasks with minimal effort \cite{ref7}. This can have far-reaching implications for a wide range of applications, from language understanding and generation to task-oriented dialogue systems and knowledge-intensive applications \cite{ref8}. Furthermore, in-context learning has shown promise in domains beyond natural language processing, such as computer vision and scientific computing \cite{ref9,ref10}, opening up new avenues for integrating language-based reasoning with other forms of intelligence \cite{ref11,ref12}.

Despite the remarkable success of in-context learning, the underlying mechanisms that enable this capability in LLMs are not yet fully understood \cite{ref6}. Researchers have proposed various theoretical frameworks and empirical analyses to shed light on the factors that contribute to the emergence of in-context learning, such as the properties of the pretraining data \cite{ref5}, the architectural design of the language models \cite{ref13}, and the role of prompt engineering \cite{ref14}. Furthermore, the limitations and challenges associated with in-context learning have also been the subject of investigation, highlighting the sensitivity to the choice and design of the demonstration examples \cite{ref15}, the susceptibility to biases \cite{ref16}, and the need for scalable and efficient algorithms \cite{ref17,ref18}.

As the field of in-context learning continues to evolve, researchers are exploring various future directions, such as the integration of in-context learning with other learning paradigms \cite{ref2,ref19}, the development of novel architectural designs \cite{ref13}, and the potential for cross-modal and cross-lingual transfer \cite{ref11,ref20}. These advancements promise to further enhance the capabilities and adaptability of LLMs, ultimately leading to more versatile and human-centric artificial intelligence \cite{ref21}.

\subsection{Foundations of In-Context Learning}

The foundations of in-context learning are built upon several key techniques and methodologies that enable language models to effectively adapt to new tasks and contexts. These include prompt engineering, meta-learning, few-shot learning, and contextual representation learning.

Prompt engineering refers to the process of designing informative prompts that can elicit the desired model behaviors \cite{ref22,ref23}. Prompts can be viewed as a form of task specification, guiding the model to understand the desired output and perform the appropriate reasoning steps. Effective prompt engineering is crucial for in-context learning, as the quality and relevance of the prompts can significantly impact the model's ability to adapt to new tasks \cite{ref24}.

Meta-learning, also known as learning-to-learn, is another fundamental technique that enables in-context learning \cite{ref25}. Meta-learning approaches, such as Model-Agnostic Meta-Learning (MAML) and prototypical networks, aim to learn task-agnostic representations and optimization strategies that can be quickly adapted to new tasks with limited data \cite{ref26,ref27}. By learning to learn, meta-learning models can rapidly adapt to new tasks, even with small amounts of training data, which is a key requirement for in-context learning \cite{ref28}.

Few-shot learning is closely related to in-context learning, as it focuses on the ability to learn new concepts from a small number of examples \cite{ref29,ref30}. Few-shot learning techniques leverage contextual information, data augmentation, and transfer learning to enhance the model's performance on new tasks with limited training data \cite{ref31}. The ability to learn effectively from a few examples is a crucial aspect of in-context learning, as it allows the model to quickly adapt to new tasks without requiring extensive fine-tuning \cite{ref32}.

Contextual representation learning is another key component of in-context learning, as it enables models to effectively capture and leverage various types of contextual cues, such as multi-modal context, visual-semantic context, and temporal context \cite{ref33,ref34}. By learning contextual representations that can adapt to the current task and input, models can better utilize the available information to improve their in-context learning performance \cite{ref35}.

The combination of these techniques, namely prompt engineering, meta-learning, few-shot learning, and contextual representation learning, forms the foundation of in-context learning. By leveraging these approaches, language models can effectively adapt to new tasks and contexts, even with limited training data, making in-context learning a powerful and versatile paradigm for a wide range of applications.

\subsection{The Role of Contextual Information}

The role of contextual information is paramount in enabling effective in-context learning, as it allows models to leverage a wealth of cues beyond the primary input to enhance their understanding and performance. In-context learning, by its very nature, relies on the model's ability to extract and utilize relevant contextual information from the provided examples and the current input. Various types of contextual information, including multi-modal, visual-semantic, and temporal context, have been explored and shown to play a crucial role in boosting the capabilities of in-context learning models.

Multi-modal context, which encompasses information from different modalities such as text, images, and audio, has been a key focus in the in-context learning domain. Recent studies have demonstrated the benefits of incorporating multi-modal context to enhance the performance of in-context learning models \cite{ref36}. For instance, the MERLOT model \cite{ref36} learns to jointly understand events in the visual world by watching and listening to millions of video clips with transcribed speech, enabling it to perform strongly on a variety of video-based tasks through in-context learning.

Visual-semantic context, which captures the relationships between visual and textual information, has also been extensively studied in the context of in-context learning. The integration of visual features and semantic representations has been shown to be highly beneficial for improving the model's ability to understand and reason about visual scenes, as well as to adapt to new tasks through in-context learning \cite{ref37}. For example, the work by \cite{ref37} demonstrates how the incorporation of visual-semantic context, through a multi-knowledge representation approach, can significantly enhance the in-context learning capabilities of vision-language models.

Temporal context, which captures the dynamic and sequential nature of information, has also been recognized as an important factor in enabling effective in-context learning, particularly in tasks involving time-series data or sequential decision-making. The \cite{ref38} study shows how incorporating temporal context, through the use of a transformer-based model that leverages the action sequence information, can improve the in-context learning performance for egocentric action recognition tasks.

The integration of these various types of contextual information has been a subject of extensive research in the in-context learning domain. The \cite{ref11} work, for instance, presents a novel framework that jointly models text and visual prompts in a unified representation space, enabling the model to effectively handle in-context learning tasks with multimodal output. Similarly, the \cite{ref39} approach demonstrates how the strategic incorporation of multimodal in-context examples can significantly enhance the performance of large multimodal models in various downstream tasks.

The importance of contextual information in in-context learning extends beyond just the input modalities. Researchers have also explored the role of prompts, which can be considered a form of contextual information, in shaping the in-context learning capabilities of models. The \cite{ref24} study, for instance, highlights the crucial impact of prompt selection and prompt fusion on the performance of in-context learning in computer vision tasks.

Furthermore, the \cite{ref40} work showcases the significance of effectively modeling the intra- and inter-modality interactions, which can be seen as a form of contextual information, for enhancing the in-context learning abilities of multimodal models.

In summary, the role of contextual information, including multi-modal, visual-semantic, and temporal context, has been extensively explored in the in-context learning domain. The incorporation of these various types of contextual cues has been shown to be instrumental in improving the performance and adaptability of in-context learning models across a wide range of tasks and domains. As the field of in-context learning continues to evolve, the strategic leveraging of contextual information will undoubtedly remain a key focus for researchers and practitioners alike.

\subsection{Evaluation and Benchmarking of In-Context Learning}

Evaluating the performance of in-context learning models is a critical aspect of understanding and advancing this emerging paradigm. Researchers have developed a variety of benchmarks, datasets, and evaluation protocols to assess the capabilities of in-context learning systems and identify their strengths and limitations.

One key area of evaluation is the development of benchmark datasets specifically designed to assess in-context learning abilities. These datasets often include a diverse set of tasks, ranging from natural language understanding to generation and reasoning, and are constructed to capture the multi-faceted nature of in-context learning performance \cite{ref41,ref42}.

In addition to task-specific datasets, researchers have also proposed more general evaluation frameworks that aim to capture the multi-dimensional aspects of in-context learning. For example, the InstructEval suite \cite{ref43} includes a diverse set of language models, tasks, and evaluation metrics to enable a comprehensive assessment of instruction selection methods for in-context learning. Similarly, the Ludwig Benchmarking Toolkit \cite{ref44} provides a configurable and extensible platform for running end-to-end benchmark studies, allowing users to customize the evaluation based on their specific needs and objectives.

When it comes to the actual evaluation metrics used to assess in-context learning, researchers have explored a wide range of approaches. Traditional metrics, such as accuracy, F1-score, and BLEU, have been widely adopted to gauge the performance of in-context learning models on various tasks \cite{ref45,ref46}. However, these metrics may not fully capture the nuances of in-context learning, as they often focus on the final output rather than the underlying learning process.

To address this, researchers have proposed more specialized metrics that aim to better reflect the unique characteristics of in-context learning. For example, the Normalized Invariability to Choice of Examples (NICE) metric \cite{ref45} quantifies the degree to which a task's learnability is dependent on the choice of in-context examples, providing insights into the stability and robustness of in-context learning. Similarly, the pointwise \(\mathcal{V}\)-usable information (PVI) metric \cite{ref46} has been adapted to the in-context setting, allowing researchers to assess the difficulty of learning a task from a given set of examples.

Alongside the development of specialized metrics, researchers have also recognized the importance of standardized evaluation protocols and leaderboards to facilitate fair comparisons between different in-context learning models \cite{ref47}. This has led to the creation of initiatives such as the GEM benchmark \cite{ref42}, which provides a modular infrastructure for evaluating natural language generation models, including those that leverage in-context learning.

However, the evaluation of in-context learning is not without its challenges. One major issue is the sensitivity of in-context learning performance to the choice and formatting of the prompt, which can significantly impact the model's performance \cite{ref48}. This has led to the development of approaches that aim to mitigate the impact of prompt design, such as the use of prompt ensembles \cite{ref48}.

Another challenge is the potential for biases and inconsistencies in the evaluation datasets themselves. As in-context learning models are expected to perform well on a wide range of tasks and domains, it is important to ensure that the evaluation datasets capture a diverse range of linguistic and semantic phenomena \cite{ref49}. Researchers have proposed the use of synthetic tasks and datasets \cite{ref49,ref50} as well as the integration of human-in-the-loop evaluation \cite{ref49} to address these concerns.

Looking to the future, the continued development of comprehensive and robust evaluation frameworks for in-context learning will be crucial for driving progress in this field. Researchers will need to explore novel evaluation metrics, datasets, and protocols that can capture the nuances of in-context learning, while also ensuring the reliability, fairness, and generalizability of the assessment process. By addressing these challenges, the research community can unlock a deeper understanding of the capabilities and limitations of in-context learning, paving the way for more effective and impactful applications of this emerging paradigm.

\subsection{Limitations, Challenges, and Ethical Considerations}

In-context learning (ICL) has emerged as a powerful paradigm that enables large language models (LLMs) to adapt to new tasks and perform them effectively without updating their parameters. However, this promising capability comes with its own set of limitations, challenges, and ethical considerations that need to be addressed.

One of the primary limitations of in-context learning is its sensitivity to prompt design. The performance of ICL models is critically dependent on the choice and ordering of the examples provided in the prompt. Small variations in the prompt can lead to significant changes in the model's behavior and output, making it challenging to ensure reliable and consistent performance across different contexts \cite{ref51}. This sensitivity to prompt design can also amplify the susceptibility of ICL models to biases, which is another significant challenge that plagues this learning paradigm \cite{ref52,ref53}.

LLMs, which form the backbone of ICL systems, are known to exhibit various types of biases, including social biases, stereotypes, and other undesirable patterns that can be perpetuated and even amplified during the in-context learning process \cite{ref53,ref54}. These biases can lead to unfair and discriminatory decisions, particularly in high-stakes applications such as hiring, lending, and healthcare, where the consequences of such biases can be severe.

Moreover, the need for scalable and efficient algorithms in in-context learning is another critical challenge. As the number of in-context examples and the complexity of the tasks increase, the computational and memory requirements of ICL models can become prohibitively high, making it difficult to deploy them in real-world scenarios \cite{ref55}. Addressing this challenge requires the development of novel algorithms and architectures that can effectively leverage contextual information while maintaining a reasonable computational footprint.

The ethical implications of in-context learning also deserve careful consideration. The ability of ICL models to adapt to new tasks and contexts raises concerns about transparency, accountability, and the potential for unintended consequences \cite{ref56,ref57}. As these models are increasingly integrated into decision-making systems, it is crucial to ensure that they adhere to ethical principles and consider the broader societal impact of their decisions.

One key ethical concern is the potential for ICL models to perpetuate or amplify existing biases and inequalities, as discussed earlier. This can lead to the disproportionate impact of algorithmic decisions on marginalized or vulnerable groups, further exacerbating social injustice \cite{ref58,ref59}. Addressing these biases requires a comprehensive approach that goes beyond simply optimizing for fairness metrics, but also considers the underlying societal and historical factors that shape the data used to train these models \cite{ref60}.

Another ethical challenge is the need for transparency and accountability in the deployment of ICL systems. As these models become more complex and opaque, it becomes increasingly difficult to understand the reasoning behind their decisions and hold them accountable for the consequences \cite{ref61,ref62}. This raises concerns about the trustworthiness and reliability of ICL systems, particularly in high-stakes applications where the stakes are high.

To mitigate these limitations, challenges, and ethical concerns, researchers and practitioners in the field of in-context learning must adopt a multifaceted approach. This may involve developing novel prompt engineering techniques that are less sensitive to variations, incorporating robust debiasing mechanisms into the learning process, and designing scalable and efficient algorithms that can maintain performance while reducing computational requirements \cite{ref63,ref64}.

Additionally, there is a pressing need to integrate ethical considerations and human oversight into the development and deployment of ICL systems. This may involve the creation of transparent and auditable decision-making processes, the establishment of ethical guidelines and frameworks, and the involvement of diverse stakeholders in the design and evaluation of these systems \cite{ref65,ref66}.

By addressing these limitations, challenges, and ethical concerns, the field of in-context learning can continue to progress and unlock the full potential of this powerful paradigm, while ensuring that it is deployed in a responsible and trustworthy manner that benefits society as a whole.

\subsection{Future Directions and Conclusion}

The field of in-context learning (ICL) has seen rapid advancements in recent years, driven by the emergence of large language models (LLMs) that have demonstrated remarkable capabilities in leveraging contextual information to quickly adapt to new tasks \cite{ref6,ref67}. As ICL continues to evolve, several promising future directions have emerged that hold the potential to further enhance its capabilities and versatility.

One key future direction is the integration of ICL with other learning paradigms, such as transfer learning, meta-learning, and multi-task learning \cite{ref68,ref69}. By leveraging the synergies between these complementary approaches, researchers can explore ways to develop more adaptable, sample-efficient, and generalized ICL models. For instance, meta-learning techniques can enable ICL models to rapidly adapt to new tasks by learning task-agnostic representations and optimization strategies \cite{ref2}. Similarly, the integration of multi-task learning can equip ICL models with the ability to simultaneously learn and transfer knowledge across a diverse range of tasks, further enhancing their versatility and real-world applicability \cite{ref70}.

Another promising direction involves the exploration of novel architectural designs and attention mechanisms that can better capture and integrate various types of contextual cues, such as multi-modal, visual-semantic, and temporal context \cite{ref71,ref72}. These innovations can further enhance the in-context learning capabilities of LLMs and vision-language models, enabling them to more effectively leverage contextual information \cite{ref11,ref73}.

The potential for cross-modal and cross-lingual transfer also represents a crucial future direction for in-context learning \cite{ref74,ref75}. Enabling ICL models to effectively adapt and generalize across different modalities and diverse languages can significantly broaden the reach and impact of these technologies. Addressing the challenges associated with cross-modal and cross-lingual transfer, such as the development of robust multimodal representations and effective alignment strategies, can unlock new application domains and foster more inclusive and accessible ICL-based systems \cite{ref76,ref77}.

As ICL models become more capable of rapidly adapting to new tasks and contexts, they can play a crucial role in shaping the development of more adaptive, efficient, and human-centric AI systems \cite{ref78,ref21}. By enabling AI agents to learn and reason in a more contextual and interactive manner, ICL can contribute to the realization of artificial general intelligence (AGI) and the creation of AI systems that can better understand, communicate, and collaborate with humans \cite{ref73,ref79}.

\section{Fundamental Techniques and Approaches for In-Context Learning}

\subsection{Prompt Engineering}

Prompt engineering has emerged as a critical technique for harnessing the immense potential of large language models (LLMs) and enabling them to effectively adapt to new tasks and contexts \cite{ref80,ref81}. At its core, prompt engineering involves designing informative prompts that can elicit the desired behaviors from LLMs, guiding them to perform a wide range of natural language processing tasks, from text generation to question answering, commonsense reasoning, and beyond.

The paramount importance of prompt engineering lies in the fact that LLMs are trained on vast, diverse corpora, which endows them with a rich knowledge base and the ability to generalize to novel situations. However, this general-purpose nature also means that LLMs may not be optimally tuned for specific tasks or applications. Prompt engineering bridges this gap by providing a flexible and efficient means of customizing the LLM's behavior to the task at hand \cite{ref82}.

One of the key aspects of prompt engineering is the design of effective prompts. Prompts can take various forms, such as natural language instructions, templates, or even sequences of tokens that do not necessarily adhere to grammatical conventions \cite{ref83}. The choice of prompt can significantly impact the LLM's performance, as it serves as a guiding mechanism that shapes the model's internal representations and reasoning processes \cite{ref84,ref85}.

Researchers have explored a wide range of techniques for designing effective prompts, including manual prompt engineering by domain experts and automated prompt engineering approaches that leverage optimization algorithms and meta-learning strategies \cite{ref86,ref87,ref88}. One prominent class of prompt engineering techniques is prompt tuning, which involves learning a set of prompt parameters that can be efficiently fine-tuned on a specific task \cite{ref89,ref90}.

Another important aspect of prompt engineering is the role of contextual information. Prompts can be designed to leverage various types of contextual cues, such as multi-modal information, visual-semantic relationships, or temporal dynamics, to enhance the LLM's understanding and reasoning capabilities \cite{ref91,ref92,ref93}. By incorporating these contextual elements into the prompt, researchers have shown that LLMs can better adapt to the task at hand and produce more accurate and coherent outputs.

However, prompt engineering is not without its challenges. LLMs can be highly sensitive to the specific wording and structure of the prompt, and prompts may be susceptible to biases present in the training data or the LLM itself \cite{ref51,ref94}. To address these challenges, researchers have explored techniques such as prompt ensemble methods and strategies for prompt calibration and optimization \cite{ref81,ref95}.

Overall, prompt engineering has emerged as a powerful and versatile technique for harnessing the capabilities of large language models. By carefully designing prompts that can effectively guide the LLM's behavior, researchers and practitioners can unlock the full potential of these powerful AI systems and apply them to a wide range of natural language processing tasks and applications.

\subsection{Meta-Learning}

Meta-learning, also known as ``learning to learn,'' aims to enable models to rapidly adapt to new tasks using limited data. This is particularly relevant in the context of in-context learning, where models need to quickly adapt to new contexts or tasks presented to them. Two prominent meta-learning approaches that are closely tied to in-context learning are Model-Agnostic Meta-Learning (MAML) and prototypical networks.

MAML is a gradient-based meta-learning algorithm that learns a good initialization for model parameters, which can then be quickly adapted to new tasks through a few gradient updates \cite{ref96,ref97}. The key idea behind MAML is to learn a shared initialization of the model parameters that can be efficiently fine-tuned on new tasks, rather than learning task-specific models from scratch. During meta-training, MAML optimizes for an initialization that can be rapidly adapted to new tasks through a few gradient descent steps, as opposed to learning a single model for all tasks \cite{ref98}.

A critical aspect of MAML is its ability to learn task-agnostic representations and optimization strategies. By optimizing for an initialization that can be quickly adapted, MAML encourages the model to learn features and representations that are broadly useful across the distribution of tasks, rather than specialized for any particular task \cite{ref99}. This task-agnostic nature of the learned representations helps in the context of in-context learning, where models need to be able to quickly adapt to new and diverse contexts.

In addition to MAML, prototypical networks are another prominent meta-learning approach that has shown promise for in-context learning \cite{ref100}. Prototypical networks learn a metric space in which classification can be performed by computing distances to prototype representations of each class. During meta-training, the network learns to embed examples from new tasks into this metric space, such that the prototypes can be quickly computed from a small number of examples and used for classification. Similar to MAML, prototypical networks aim to learn task-agnostic representations that can be effectively leveraged for few-shot adaptation to new tasks \cite{ref100}.

While MAML and prototypical networks have shown impressive results in various in-context learning tasks, there are still several challenges and open questions that need to be addressed. For example, MAML is known to be sensitive to hyperparameter tuning and can be computationally expensive due to its nested optimization structure \cite{ref101}. Researchers have proposed various extensions and modifications to MAML to address these limitations, such as Alpha MAML \cite{ref102}, which incorporates an online hyperparameter adaptation scheme, and Contextual Gradient Scaling (CxGrad) \cite{ref103}, which scales the gradients of the backbone layers to facilitate task-specific adaptation.

Another important consideration in meta-learning for in-context learning is the role of task similarity and the distribution of tasks encountered during meta-training. It has been shown that the performance of MAML can be highly dependent on the diversity and similarity of the tasks used during meta-training \cite{ref104}. Researchers have proposed approaches like Task-Agnostic Meta-Learning (TAML) \cite{ref105} and Multimodal MAML (MMAML) \cite{ref106} to better handle diverse task distributions and improve the generalization of meta-learned representations.

Additionally, recent work has explored the importance of representation change and feature learning in meta-learning, rather than just feature reuse \cite{ref107,ref108}. These studies suggest that the ability to learn new features and representations during the adaptation phase can be crucial for effective in-context learning, especially when the new context or task is significantly different from the meta-training tasks.

In summary, meta-learning approaches like MAML and prototypical networks have shown great promise for in-context learning, as they enable models to quickly adapt to new contexts or tasks using limited data. These methods learn task-agnostic representations and optimization strategies that can be effectively leveraged for rapid adaptation. However, there are still several challenges and open questions, such as addressing the computational and hyperparameter sensitivity of MAML, handling diverse task distributions, and understanding the role of representation change versus feature reuse in meta-learning for in-context learning. Ongoing research in these areas will continue to advance the capabilities of meta-learning for in-context learning and other related applications.

\subsection{Few-Shot Learning}

Few-shot learning (FSL) has emerged as a promising approach to tackling the challenge of learning new concepts from limited labeled data \cite{ref109}. The core idea behind FSL is to leverage prior knowledge, often in the form of representations learned from a large corpus of data, to enable rapid adaptation to novel tasks with few examples.

One of the fundamental techniques underlying FSL is meta-learning, which aims to learn a learning algorithm that can quickly adapt to new tasks \cite{ref110}. The key idea is to train the meta-learner on a diverse set of tasks, such that it can acquire the necessary inductive biases and optimizers to effectively learn new tasks from limited data. Prominent examples of meta-learning approaches include Model-Agnostic Meta-Learning (MAML) \cite{ref111} and Prototypical Networks \cite{ref112}, which learn task-agnostic representations and optimization strategies, respectively.

Beyond meta-learning, another important technique for FSL is few-shot data augmentation, which aims to generate new labeled examples to enrich the limited training data \cite{ref113}. These data augmentation techniques often leverage the inherent structure and properties of the data, such as semantic relationships and visual-spatial cues, to synthesize plausible examples that can improve the generalization of the learned models. For instance, \cite{ref113} proposed a curriculum-based data augmentation strategy that gradually introduces augmented samples to the training process, leading to faster convergence and better performance.

Contextual information has also been shown to play a crucial role in enhancing few-shot learning performance \cite{ref33,ref34}. By incorporating relevant contextual cues, such as object co-occurrences, semantic relationships, and spatiotemporal dependencies, FSL models can better leverage the available information to learn more generalizable representations. \cite{ref33} proposed a plug-and-play module that can identify task-relevant features by exploiting the contextual information in the support set, leading to significant performance improvements on few-shot object classification tasks.

Transfer learning is another key component in many successful FSL approaches, where pre-trained models are fine-tuned or adapted to the target few-shot tasks \cite{ref114,ref115}. By leveraging the knowledge acquired from large-scale datasets, FSL models can gain a strong initial understanding of the underlying patterns, which can then be fine-tuned with the limited target data. However, the effectiveness of transfer learning in FSL is highly dependent on the similarity between the source and target domains, as well as the specific fine-tuning strategies employed \cite{ref116}.

Recent advancements in self-supervised learning have also demonstrated the potential to boost the performance of FSL models \cite{ref117,ref118}. By learning rich and transferable representations from unlabeled data, self-supervised pre-training can provide a strong foundation for few-shot learning tasks, often outperforming supervised pre-training approaches \cite{ref119}.

Despite the significant progress in FSL, several challenges remain. One key issue is the sensitivity of FSL models to the design of prompts or task descriptions, which can dramatically impact their performance \cite{ref51}. Addressing this challenge requires developing more robust and generalizable FSL approaches that are less dependent on specific prompt formulations. Another challenge is the susceptibility of FSL models to various biases, such as data biases and social biases, which can get amplified during the few-shot learning process \cite{ref94}. Developing debiasing techniques and more transparent FSL models is crucial for ensuring the fairness and trustworthiness of these systems.

Furthermore, the computational and memory efficiency of FSL algorithms remains an important challenge, particularly as the number of in-context examples and the complexity of the tasks increase \cite{ref120}. Exploring novel architectures and optimization strategies that can effectively leverage contextual information while maintaining high efficiency is a promising direction for future research.

\subsection{Contextual Representation Learning}

Contextual information plays a crucial role in enabling effective in-context learning, as it allows models to adapt to different tasks and situations by leveraging relevant background knowledge and relationships. Researchers have explored various methods for learning contextual representations that can capture and utilize diverse types of contextual cues, such as multi-modal context, visual-semantic context, and temporal context.

One key approach to learning contextual representations is to leverage multi-modal information \cite{ref121,ref122,ref123}. By incorporating data from multiple modalities, such as text, images, and audio, models can learn richer and more comprehensive representations that better capture the contextual relationships between different aspects of a given situation or task. For example, \cite{ref124} demonstrates how incorporating complementary information from RGB images, optical flow, and depth maps can enhance the performance of video temporal grounding tasks.

Similarly, researchers have explored the use of visual-semantic context to improve in-context learning \cite{ref125,ref126,ref127}. By modeling the relationships between visual features and their semantic interpretations, models can learn more meaningful and coherent representations that better capture the underlying context. For instance, \cite{ref125} shows how generating contextual descriptions of images can help improve the performance of visual question answering tasks.

Temporal context is another important aspect that has been explored in the context of in-context learning \cite{ref128,ref39,ref11}. By taking into account the temporal dynamics and dependencies within a sequence of data, models can learn more robust and contextually-aware representations that better capture the evolving nature of the task or situation. For example, \cite{ref128} demonstrates how incorporating temporal context can improve the performance of a multimodal dialogue system for learning visually grounded word meanings.

In addition to these specific types of contextual information, researchers have also explored more general techniques for learning contextual representations that can be broadly applied to in-context learning tasks. One such approach is the use of prompt engineering, which involves designing informative prompts that can effectively elicit the desired model behaviors and capture the relevant contextual information \cite{ref39,ref129,ref130}. By carefully crafting the prompts, models can be guided to focus on the most relevant contextual cues and adapt their representations accordingly.

Another promising direction is the use of meta-learning and few-shot learning techniques \cite{ref39,ref35}, which enable models to rapidly adapt to new tasks and contexts by leveraging a small amount of labeled data. These approaches often involve learning task-agnostic representations and optimization strategies that can be effectively applied to a wide range of in-context learning scenarios.

Furthermore, researchers have explored the integration of contextual information with other learning paradigms, such as contrastive learning \cite{ref131,ref131}, which can help models learn more discriminative and contextually-aware representations by explicitly capturing the relationships between different modalities or data points.

Overall, the development of effective contextual representation learning techniques has been a crucial area of research in enabling robust and adaptable in-context learning. By leveraging multi-modal information, visual-semantic context, temporal dynamics, and other contextual cues, models can learn more comprehensive and contextually-aware representations that can be effectively applied to a wide range of in-context learning tasks and scenarios.

\section{Contextual Information Modeling and Utilization}

\subsection{Modeling and Leveraging Multi-modal Context}

The role of multi-modal context, including visual-semantic context and temporal context, is crucial in enabling effective in-context learning. In-context learning (ICL) has gained significant attention in recent years, particularly with the emergence of large language models (LLMs) \cite{ref36}. These LLMs have demonstrated impressive capabilities in adapting to new tasks and contexts by leveraging the information provided in the input, often referred to as the ``context.'' However, the context in these models is typically limited to textual information, overlooking the wealth of information available in other modalities, such as visual and temporal data.

Recognizing the importance of multi-modal context, researchers have started exploring ways to incorporate different types of contextual information to enhance the performance of in-context learning models \cite{ref132,ref39}. Visual-semantic context, which captures the relationships between visual features and their semantic meanings, has been shown to be particularly beneficial for tasks involving visual reasoning and understanding \cite{ref125}. By leveraging the information contained in visual data, in-context learning models can better ground their language-based understanding in the physical world, leading to more accurate and coherent responses \cite{ref133}.

Temporal context, on the other hand, is crucial for in-context learning models operating on time-series data, such as video understanding tasks \cite{ref124}. By capturing the temporal dynamics and dependencies within the input data, these models can better reason about events, predict future outcomes, and provide more comprehensive and holistic responses \cite{ref134}.

Researchers have explored various approaches to effectively model and leverage multi-modal context in in-context learning. One common approach is to use multimodal fusion techniques, where information from different modalities is combined to create a unified representation \cite{ref40,ref135}. These fusion methods can range from simple concatenation to more sophisticated attention-based mechanisms that dynamically weigh the importance of different modalities based on the task at hand \cite{ref72}.

Another key aspect in modeling multi-modal context is the effective extraction and integration of contextual features. Techniques such as cross-modal feature interaction and co-attention mechanisms have been used to capture the interplay between different modalities, enabling the in-context learning model to better understand the relationships and dependencies between the various sources of information \cite{ref136,ref137}.

Prompt engineering, a crucial component of in-context learning, has also been explored in the context of multi-modal data \cite{ref11,ref138}. Researchers have investigated how to design effective prompts that can leverage contextual information from various modalities to enhance the performance of in-context learning models \cite{ref139,ref140}.

It is worth noting that the integration of multi-modal context is not limited to the field of in-context learning. Broader applications in computer vision and multimodal settings have also benefited from the incorporation of contextual information \cite{ref141,ref142}. By leveraging the contextual cues from different modalities, these models have demonstrated improved performance in tasks such as image classification, object detection, and multimodal dialogue generation \cite{ref143,ref144}.

However, the effective modeling and utilization of multi-modal context in in-context learning is not without its challenges. Researchers have identified issues such as the sensitivity of in-context learning models to prompt design, the susceptibility to biases, and the need for scalable and efficient algorithms \cite{ref139,ref145}. Addressing these challenges is crucial for the widespread adoption and practical deployment of in-context learning models in real-world applications.

\subsection{Modeling and Leveraging Visual-Semantic Context}

Modeling and leveraging visual-semantic context is an essential component of in-context learning, as it enables models to effectively integrate and utilize the relationships between visual and textual information. Recent advancements in vision-language models have demonstrated the power of capturing and exploiting visual-semantic context to enhance the performance of various tasks, including in-context learning.

One key aspect of modeling visual-semantic context is the integration of visual features and semantic representations. Several studies have explored different approaches to achieve this integration. For instance, the work on ``Probing Multimodal Embeddings for Linguistic Properties: the Visual-Semantic Case'' \cite{ref146} investigates the complementary nature of visual and textual information in visual-semantic embeddings. The authors propose a generalization of probing tasks to the visual-semantic domain and demonstrate that visual-semantic embeddings can outperform their unimodal counterparts in capturing linguistic properties, suggesting the benefits of integrating visual and semantic information.

Similarly, the ``Compositional Mixture Representations for Vision and Text'' \cite{ref147} paper presents a model that learns a shared Gaussian mixture representation, imposing the compositionality of text onto the visual domain without explicit location supervision. By combining a spatial transformer with a representation learning approach, the model learns to split images into separately encoded patches and associate visual and textual representations in an interpretable manner. This approach showcases the potential of modeling the relationships between visual and semantic information to enhance in-context learning.

Another relevant work, ``IMProv: Inpainting-based Multimodal Prompting for Computer Vision Tasks'' \cite{ref148}, explores a generative model that can perform in-context learning for various computer vision tasks by leveraging both textual and visual prompts. The model, trained on a dataset of figures from computer vision papers and their associated captions, is able to generate visual outputs conditioned on multimodal prompts, demonstrating the effectiveness of integrating visual and semantic information for in-context learning.

Furthermore, the ``Towards More Unified In-context Visual Understanding'' \cite{ref11} paper presents a novel in-context learning framework for visual understanding tasks that enables multimodal output. The authors propose a unified representation space for both text and visual prompts, which are then processed by a decoder-only transformer architecture to perform generative modeling. This approach highlights the importance of modeling the interplay between visual and textual information to facilitate more comprehensive in-context learning in the visual domain.

Researchers have also investigated the role of visual-semantic context in enhancing the performance of specific in-context learning tasks. For instance, the ``Investigating the Role of Attribute Context in Vision-Language Models for Object Recognition and Detection'' \cite{ref149} paper examines the influence of attribute context in vision-language models, such as its impact on object embeddings and the utility of attribute context in training for open-vocabulary object detection and fine-grained text-region retrieval. The findings suggest that attribute context can provide valuable information for improving in-context learning in these tasks, but its effective utilization remains an ongoing challenge.

Another example is the ``Efficient Object-Level Visual Context Modeling for Multimodal Machine Translation: Masking Irrelevant Objects Helps Grounding'' \cite{ref150} paper, which explores the application of in-context learning in the domain of multimodal machine translation. The authors propose an object-level visual context modeling framework that encourages the model to ground the translation on relevant visual objects by masking irrelevant objects. This approach demonstrates the importance of effectively modeling and leveraging visual-semantic context to enhance the performance of in-context learning in specialized domains.

Overall, the research on modeling and leveraging visual-semantic context has shown promising results in improving the performance of in-context learning. By integrating visual features and semantic representations, models can better capture the relationships between visual and textual information, leading to more effective in-context adaptation and generalization. However, the field still faces challenges, such as the sensitivity to prompt design, the need for more comprehensive evaluation frameworks, and the development of scalable and efficient algorithms to handle the complexity of visual-semantic context. Continued research in this area is crucial for advancing the capabilities of in-context learning, particularly in the context of multimodal tasks and applications.

\subsection{Modeling and Leveraging Temporal Context}

The ability to effectively model and leverage temporal information is critical for achieving robust in-context learning, particularly in domains involving temporal dynamics, such as video understanding. Existing research has explored various approaches to address this challenge.

One key aspect of modeling temporal context is the recognition that different tasks may require reasoning over different temporal scales. For example, in video understanding, some tasks may focus on short-term dynamics, such as recognizing individual actions, while others may require understanding longer-term dependencies, such as recognizing complex events unfolding over time \cite{ref151,ref151}. To address this, researchers have proposed multi-granular temporal aggregation frameworks that can effectively capture temporal information at multiple scales \cite{ref152}.

Another important consideration is the role of temporal ordering and dynamics in in-context learning. Several studies have found that simply shuffling the frames in a video can significantly degrade the performance of in-context learning models, suggesting that these models are heavily reliant on capturing the correct temporal relationships \cite{ref153,ref154}. This highlights the need for in-context learning models to develop robust temporal reasoning capabilities that go beyond mere pattern matching.

Recent work has explored various approaches to modeling and leveraging temporal context for in-context learning. One line of research has focused on incorporating explicit temporal representations, such as temporal positional encodings or temporal attention mechanisms, into neural network architectures \cite{ref155,ref156}. These approaches aim to enable the models to effectively capture and reason about the temporal dynamics of the input data.

Another approach has been to design self-supervised learning objectives that encourage the models to learn temporal relationships, such as predicting the relative order of frames or reconstructing the temporal evolution of visual features \cite{ref157,ref158}. These techniques can help the models develop a stronger understanding of temporal patterns, which can then be leveraged for in-context learning tasks.

Additionally, researchers have explored the integration of temporal context with other modalities, such as language, to further enhance in-context learning capabilities. For example, studies have shown that the combination of visual and language features can lead to improved temporal reasoning and understanding \cite{ref159,ref160}.

Beyond simply modeling temporal context, some studies have also investigated the effective utilization of temporal information for in-context learning. Approaches such as temporal reasoning over graph-structured representations \cite{ref161,ref162} or the integration of temporal context with causal reasoning \cite{ref163} have demonstrated promising results in various video understanding tasks.

It is also worth noting that the modeling and utilization of temporal context is not limited to the video domain. Researchers have explored the role of temporal information in other specialized domains, such as scientific computing \cite{ref164}, reinforcement learning \cite{ref165}, and speech processing \cite{ref166}. These studies highlight the broader applicability of temporal context modeling and the need for domain-specific considerations in the design of in-context learning approaches.

In summary, the modeling and leveraging of temporal context is a crucial aspect of achieving robust in-context learning, particularly in domains involving temporal dynamics. Existing research has explored various approaches, including multi-granular temporal aggregation, the incorporation of explicit temporal representations, self-supervised learning objectives for temporal understanding, and the integration of temporal context with other modalities and reasoning techniques. As the field of in-context learning continues to evolve, the effective modeling and utilization of temporal information will likely remain a key area of focus for researchers and practitioners alike.

\subsection{Contextual Feature Extraction and Fusion}

Extracting and fusing contextual features from multiple modalities is crucial for enabling seamless integration within in-context learning frameworks. Given the diverse nature of contextual information, such as visual, semantic, and temporal cues, effective techniques are needed to capture the interplay between different modalities and leverage their complementary strengths \cite{ref167,ref168}.

One key aspect of contextual feature extraction and fusion is cross-modal feature interaction, where the representations from different modalities are combined to enhance the learning of task-relevant information \cite{ref40,ref169}. Cross-attention mechanisms have proven to be effective in modeling the dependencies between modalities, allowing the model to selectively focus on the most informative features from each modality \cite{ref170,ref171}.

These cross-attention mechanisms can be extended to capture the fine-grained interactions between modalities, going beyond just aligning high-level features. For example, \cite{ref172} proposed a context-based fusion approach that uses attention to model the interplay between modality-specific features at different levels of granularity, from low-level visual and audio cues to higher-level semantic representations. Similarly, \cite{ref173} introduced a fusion matrix that simultaneously captures intra-modal and inter-modal relationships, enabling the model to effectively leverage the complementary information across modalities.

Another important aspect of contextual feature extraction and fusion is the ability to handle varying degrees of modality availability and alignment. In real-world scenarios, the input data may not always be fully aligned or contain all the expected modalities \cite{ref174,ref169}. To address this challenge, researchers have explored techniques like modality-specific feature extraction and dynamic fusion strategies \cite{ref175,ref176}.

\cite{ref169} proposed a fusion module that splits each modality into channel-wise equal feature blocks and creates a joint representation that is used to generate soft attention for each channel across the feature blocks. This approach allows the model to effectively fuse features of various spatial dimensions and sequence lengths, making it suitable for both convolutional and recurrent neural networks.

\cite{ref175} introduced two adaptive fusion networks, Auto-Fusion and GAN-Fusion, which learn to compress and align information from different modalities in a context-aware manner, without relying on a predefined fusion operation. These adaptive fusion techniques can better model the complex relationships between modalities and handle missing or noisy data.

Furthermore, \cite{ref176} proposed a graph-based fusion method that leverages self-attention and multi-head attention to capture the interactions between modalities and their sentiment-related features. This approach allows the model to adaptively fuse features while preserving the unique characteristics of each modality.

Beyond feature-level fusion, researchers have also explored techniques to integrate contextual information at the model architecture level. \cite{ref177} proposed a hierarchical fusion approach that first applies cross-task attention to facilitate the exchange of relevant information between multiple task features, and then applies cross-scale attention to aggregate useful information from feature maps at different resolutions.

\cite{ref127} introduced a transformer-based model that explicitly disentangles the text and video into semantic roles, such as objects, spatial contexts, and temporal contexts, and learns the intra- and inter-role correlations to discover discriminative features for cross-modal matching.

These architectural-level fusion techniques demonstrate the importance of carefully designing the model structure to effectively leverage the contextual information across modalities, going beyond simple feature concatenation or attention-based fusion.

In summary, the extraction and fusion of multi-modal contextual features is a crucial component of in-context learning frameworks. Researchers have explored a variety of techniques, including cross-attention mechanisms, dynamic fusion strategies, and hierarchical fusion architectures, to capture the interplay between different modalities and seamlessly integrate contextual information within the learning process. As the field of in-context learning continues to evolve, further advancements in contextual feature extraction and fusion will be key to unlocking the full potential of these powerful learning models.

\subsection{Contextual Prompt Engineering}

Prompt engineering has emerged as a crucial technique for enhancing the performance of in-context learning, particularly in multimodal settings. The design of effective prompts that can leverage contextual information from various modalities is essential for improving the model's ability to adapt to new tasks and contexts.

In multimodal in-context learning, the prompt plays a crucial role in guiding the model to effectively integrate and utilize the available contextual information. Unlike text-only prompts, multimodal prompts need to consider the interplay between different modalities, such as vision and language, to ensure that the model can fully leverage the complementary nature of the contextual information.

One key aspect of contextual prompt engineering is the design of prompts that can effectively capture and represent the relationships between the various modalities \cite{ref178}. Traditional prompt engineering techniques that focus solely on textual prompts may fall short in capturing the nuanced interactions between different modalities. To address this, researchers have explored strategies for designing prompts that can seamlessly integrate visual, audio, and other contextual cues.

For example, \cite{ref179} proposes a modality-missing-aware prompt design that can handle the presence or absence of different modalities during both training and inference. By incorporating prompts that can adapt to missing modalities, the model can maintain robust performance in real-world scenarios where certain contextual information may be unavailable.

Furthermore, \cite{ref180} introduces the concept of ``domain-invariant'' prompts, which are designed to enhance the generalization capabilities of vision-language models across different domains and classes. By learning prompts that can capture the underlying relationships and patterns in the data, the model can better adapt to unseen scenarios without requiring extensive retraining.

Another important aspect of contextual prompt engineering is the incorporation of prior knowledge and domain-specific information into the prompt design. \cite{ref181} explores the use of multimodal prompts, including both text and image-based prompts, to guide the fine-tuning of vision transformer models for image classification tasks. By leveraging domain-specific information encoded in the prompts, the model can better adapt to the target task and improve its performance.

In addition to the design of the prompts themselves, researchers have also investigated the importance of prompt positioning and formatting in the context of in-context learning. \cite{ref84} highlights the significant impact that the placement of prompts can have on the model's performance, emphasizing the need for a systematic exploration of prompt formatting strategies.

Moreover, \cite{ref138} discusses the role of contextual information in prompt engineering, particularly in terms of its impact on the model's ability to generalize and adapt to new tasks and domains. The paper suggests that the effective modeling and utilization of contextual cues, such as multi-modal, visual-semantic, and temporal context, can greatly enhance the performance of in-context learning models.

Beyond the technical aspects of prompt engineering, researchers have also explored the ethical implications of in-context learning and the potential challenges associated with it. \cite{ref81} highlights the susceptibility of in-context learning to biases, the need for scalable and efficient algorithms, and the importance of developing ethical frameworks to ensure the responsible deployment of these systems.

In summary, contextual prompt engineering is a critical component of in-context learning, particularly in multimodal settings. The design of effective prompts that can leverage the available contextual information from various modalities is crucial for enhancing the model's adaptability, generalization, and performance in real-world applications. Ongoing research in this area explores advanced prompt engineering techniques, the impact of prompt formatting, the incorporation of prior knowledge, and the consideration of ethical implications, all of which contribute to the advancement of in-context learning capabilities.

\subsection{Evaluation of Contextual In-Context Learning}

As the field of in-context learning continues to evolve, particularly with the incorporation of contextual information from multiple modalities, the design of appropriate evaluation protocols and benchmarks has become increasingly crucial. Evaluating the performance of in-context learning models that leverage contextual cues poses several unique challenges that must be addressed to ensure a comprehensive and meaningful assessment.

One of the primary challenges in evaluating contextual in-context learning is the inherent complexity and diversity of the contextual information involved. In-context learning models may utilize a wide range of contextual data, including visual-semantic representations, temporal dynamics, and multi-modal interactions \cite{ref182}. Designing evaluation protocols that can effectively capture the nuances and interdependencies of these different modalities requires a deep understanding of the underlying task requirements and the specific capabilities of the in-context learning models.

Furthermore, the evaluation of contextual in-context learning must consider the scalability and generalizability of the models. As the complexity of the contextual information increases, the evaluation process must be able to handle the growing volume of data and the potential for higher-dimensional feature spaces \cite{ref183}. This necessitates the development of comprehensive test suites and datasets that can systematically probe the models' abilities to leverage contextual information in diverse scenarios, ranging from simple to more complex and challenging tasks.

One approach to addressing these challenges is the creation of specialized benchmarks and datasets that focus on the evaluation of contextual in-context learning. These benchmarks can incorporate a wide range of contextual modalities, including visual, temporal, and multimodal information, and can be designed to assess the models' performance on tasks that require the effective integration and utilization of these contextual cues \cite{ref184}. For instance, the ConTextual benchmark \cite{ref184} has been specifically developed to evaluate the context-sensitive text-rich visual reasoning capabilities of large multimodal models, providing valuable insights into their ability to effectively leverage contextual information.

Furthermore, the design of comprehensive test suites can also contribute to the systematic evaluation of contextual in-context learning. These test suites can include a wide range of tasks and scenarios, each designed to probe specific aspects of the models' contextual understanding and reasoning abilities \cite{ref185,ref186}. By incorporating a diverse set of tasks and varying the complexity and characteristics of the contextual information, these test suites can help researchers and practitioners identify the strengths, weaknesses, and limitations of in-context learning models in a more holistic manner.

In addition to the development of specialized benchmarks and test suites, the evaluation of contextual in-context learning must also consider the use of appropriate metrics and evaluation protocols. The choice of metrics can significantly impact the assessment of the models' performance, and it is crucial to select metrics that can accurately capture the relevant aspects of contextual understanding and reasoning \cite{ref187,ref188}. For instance, in the case of visual-semantic in-context learning, metrics that focus on the overall accuracy or F1-score may not be sufficient to capture the nuances of the models' performance, and more specialized metrics may be required.

Moreover, the evaluation of contextual in-context learning must also consider the potential biases and limitations of the datasets and test suites used. As with any evaluation framework, there is a risk of biases being introduced, which can lead to an incomplete or skewed assessment of the models' capabilities \cite{ref189,ref190}. Addressing these biases and ensuring the representativeness and diversity of the evaluation datasets is crucial for obtaining reliable and meaningful insights into the performance of contextual in-context learning models.

In conclusion, the evaluation of contextual in-context learning poses unique challenges that require a comprehensive and systematic approach. The development of specialized benchmarks, test suites, and evaluation protocols that can effectively capture the nuances of contextual information and reasoning is crucial for advancing the field of in-context learning and ensuring the reliable deployment of these models in real-world applications. By addressing these challenges, researchers and practitioners can gain a deeper understanding of the capabilities and limitations of in-context learning models, ultimately leading to the creation of more robust and versatile artificial intelligence systems.

\section{In-Context Learning in Language Models}

\subsection{LLMs' In-context Learning Capabilities}

Large language models (LLMs) have demonstrated remarkable capabilities in performing in-context learning across a wide range of language-based tasks, including natural language understanding, generation, and reasoning. The emergence of LLMs, such as GPT-3 \cite{ref191}, has revolutionized the field of natural language processing, showcasing their ability to effectively adapt to new tasks and contexts with minimal training data.

One of the key factors contributing to the success of LLMs in in-context learning is their ability to effectively leverage the provided context to inform their predictions and generate coherent and relevant responses. LLMs are trained on vast amounts of text data, enabling them to build rich representations of language and the world \cite{ref192,ref193}. When presented with a new task or context, LLMs can draw upon this accumulated knowledge and adapt their behavior accordingly, often outperforming specialized models trained on the task-specific data.

The in-context learning capabilities of LLMs have been demonstrated across a wide range of language-based tasks, such as question answering, text summarization, and language translation. For example, in the task of question answering, LLMs can effectively utilize the provided context, including the question and any relevant information, to generate accurate and informative responses \cite{ref8}. This is particularly valuable in scenarios where task-specific training data may be scarce or difficult to obtain, as LLMs can leverage their general language understanding to adapt to the new task.

Furthermore, LLMs have shown promising capabilities in reasoning-based tasks, such as logical inference and mathematical problem-solving \cite{ref194}. By analyzing the in-context learning dynamics of LLMs, researchers have observed the models' ability to learn and apply formal languages and basic logical reasoning skills, showcasing their potential to tackle more complex reasoning tasks.

One of the key factors contributing to the in-context learning capabilities of LLMs is their ability to effectively leverage and generate natural language explanations. Recent research has demonstrated that incorporating human-annotated rationales or chain-of-thought prompting can significantly enhance the performance of LLMs on tasks that require reasoning \cite{ref195}. Moreover, studies have shown that LLMs can learn to generate their own natural language explanations, which can further improve their in-context learning abilities and provide valuable insights into their decision-making processes \cite{ref196}.

While large language models (LLMs) have demonstrated impressive in-context learning capabilities, they can still be prone to generating unreliable or factually incorrect outputs, a phenomenon known as hallucination \cite{ref197}. This limitation of LLMs can be particularly problematic in real-world applications where the accuracy and factuality of the generated outputs are crucial.

\subsection{Supervised Knowledge for In-context Learning}

While large language models (LLMs) have demonstrated impressive in-context learning capabilities, they can still be prone to generating unreliable or factually incorrect outputs, a phenomenon known as hallucination \cite{ref197}. This limitation can be particularly problematic in real-world applications where the accuracy and factuality of the generated outputs are crucial. To address this challenge, researchers have explored the use of supervised fine-tuned language models to enhance the reliability and generalizability of LLMs during in-context learning \cite{ref8}.

One key approach is to leverage the strengths of supervised fine-tuned models to mitigate the hallucinations and improve the factuality of LLM outputs. Supervised fine-tuning, where language models are trained on task-specific datasets, can help instill domain-specific knowledge and task-oriented capabilities in the models. By incorporating this supervised knowledge into the in-context learning process, researchers have demonstrated that LLMs can become more reliable and accurate in their outputs \cite{ref8}.

For example, \cite{ref8} proposed a framework that combines a pre-trained LLM with a task-specific fine-tuned language model (SLM) to enhance the in-context learning performance. The SLM, trained on a curated dataset for a specific task, is used to provide additional guidance and supervision to the LLM during the inference stage. This approach has been shown to improve the generalizability of the LLM, reduce hallucinations, and improve the factuality of the generated outputs across a range of knowledge-intensive tasks.

Similarly, \cite{ref198} explored a method for fine-tuning LLMs to be more factual, without relying on expensive human labeling. The authors leveraged recent advancements in automatically assessing the factuality of open-ended text, as well as direct preference optimization algorithms, to fine-tune LLMs to generate more factual outputs. Their experiments demonstrated that this approach can significantly reduce the factual error rate of LLMs on various tasks.

The benefits of incorporating supervised knowledge into LLMs' in-context learning go beyond just reducing hallucinations and improving factuality. \cite{ref199} investigated the impact of fine-tuning on the general in-context learning performance of LLMs, finding that integrating the in-context learning strategy during the fine-tuning process can help mitigate the potential decrease in the models' generalization capabilities.

Furthermore, \cite{ref200} explored the potential biases introduced during the reinforcement learning fine-tuning process, suggesting that a careful balance must be struck between task-specific fine-tuning and preserving the broader in-context learning capabilities of LLMs.

Overall, the integration of supervised fine-tuned language models into the in-context learning process of LLMs has shown promising results in enhancing the reliability, factuality, and generalizability of these powerful models. By leveraging the strengths of both pre-trained LLMs and task-specific fine-tuned models, researchers have been able to mitigate the hallucination issue and improve the overall trustworthiness of LLM-generated outputs. However, as the field continues to evolve, there is an ongoing need to carefully balance the trade-offs between task-specific fine-tuning and maintaining the broad in-context learning capabilities of LLMs, to ensure their effective and responsible deployment in real-world applications.

\subsection{In-context Learning Dynamics with Random Sequences}

Recent studies have shown that large language models (LLMs) can exhibit intriguing in-context learning (ICL) capabilities, where the models are able to adapt to new tasks by conditioning on a few input-output examples provided in the context, without any explicit fine-tuning \cite{ref194}. This emergent behavior of LLMs has sparked significant interest in the research community, as it suggests that these models can leverage the information contained in the context to rapidly acquire new skills and perform complex tasks.

One approach to better understand the mechanisms underlying ICL in LLMs is to examine their behavior when presented with seemingly random input sequences, such as binary sequences. By using random sequences as context, researchers can gain insights into how LLMs transition between seemingly random and deterministic behaviors, and how they utilize the contextual information to make predictions \cite{ref194}.

The key idea behind this line of investigation is that by analyzing the in-context learning dynamics of LLMs on random sequences, researchers can uncover the latent concepts and the underlying reasoning processes that enable these models to perform ICL. Random binary sequences are particularly useful for this purpose, as they can be easily generated and controlled, and they lack the inherent structure and semantics present in natural language data.

In a study by \cite{ref194}, the researchers investigated the in-context learning dynamics of LLMs, specifically GPT-3.5+, when presented with random binary sequences as context. The authors observed that the LLM exhibited the ability to generate seemingly random numbers and learn basic formal languages, with striking in-context learning dynamics where the model's outputs transitioned sharply from seemingly random behaviors to deterministic repetition.

The authors suggest that these findings provide a more nuanced understanding of the LLM's capabilities compared to traditional success-or-failure evaluation benchmarks, as they reveal the model's ability to adapt and reason about the underlying structure of the input sequences, even when they appear to be random.

Furthermore, the study draws inspiration from the cognitive science of human randomness perception, indicating that the ability to perceive and reason about randomness may be a fundamental characteristic of human cognition that LLMs are also able to exhibit to some extent.

Another study by \cite{ref201} investigated the in-context learning dynamics of LLMs using tasks adapted from cognitive psychology. The researchers found that LLMs update their beliefs in an asymmetric manner, learning more from better-than-expected outcomes than from worse-than-expected ones. Interestingly, this effect was reversed when the models were learning about counterfactual feedback, and disappeared when no agency was implied.

The authors corroborated these findings by investigating idealized in-context learning agents derived through meta-reinforcement learning, where they observed similar patterns. These results contribute to the understanding of how in-context learning works in LLMs, highlighting that the framing of a problem can significantly influence the learning dynamics, a phenomenon also observed in human cognition.

These studies suggest that the in-context learning capabilities of LLMs are not solely driven by their ability to memorize and regurgitate information, but rather involve more complex reasoning and adaptation processes. By probing the models' responses to seemingly random inputs and analyzing their learning dynamics, researchers can gain valuable insights into the underlying mechanisms that enable LLMs to perform ICL.

Furthermore, these investigations into the in-context learning dynamics of LLMs can have broader implications for our understanding of the nature of intelligence and cognition. The similarities between the behavior of LLMs and human cognitive processes, as observed in these studies, suggest that there may be fundamental principles underlying intelligence that are shared across both artificial and biological systems.

Overall, the exploration of in-context learning dynamics with random sequences is a promising avenue of research that can deepen our understanding of the remarkable capabilities of large language models and their potential to shed light on the nature of intelligence and cognition. As the field of AI continues to rapidly evolve, these types of controlled and targeted investigations will be crucial in unlocking the full potential of these powerful models.

\subsection{Leveraging Post-hoc Explanations for In-context Learning}

In-context learning has emerged as a powerful capability of large language models (LLMs), enabling them to quickly adapt to new tasks and contexts with minimal training. A key aspect of this success is the ability of LLMs to generate natural language explanations, also known as post-hoc explanations, that can provide insights into the models' reasoning process \cite{ref202,ref203}. These post-hoc explanations can not only help users understand the model's decision-making, but also serve as a valuable signal to further improve the model's in-context learning performance.

Recent research has investigated the use of post-hoc explanations to enhance the in-context learning capabilities of LLMs. One approach, known as AMPLIFY \cite{ref195}, leverages post-hoc explanation methods to automatically generate natural language rationales that can provide corrective signals to the LLM during in-context learning. The key idea behind AMPLIFY is to construct these rationales by embedding insights from post-hoc explanations, such as attribution scores that capture the influence of each input feature on the model's prediction. By incorporating these rationales into the in-context learning process, AMPLIFY has been shown to significantly improve the prediction accuracy of LLMs across a variety of tasks, outperforming methods that rely on human-annotated rationales.

The potential of post-hoc explanations to enhance in-context learning is further explored in the Self-AMPLIFY framework \cite{ref204}. In this approach, the researchers leverage post-hoc explanation methods applied to small language models (SLMs) to automatically generate rationales and use them to improve the SLMs' own performance through in-context learning. This method is particularly valuable as it addresses the challenge of relying on human-annotated rationales, which can be costly and labor-intensive, especially for smaller models that may not have the same in-context learning capabilities as their larger counterparts.

The potential of post-hoc explanations to improve in-context learning is not limited to text-based tasks. Researchers have also explored the use of post-hoc explanations in the context of computer vision and multimodal learning \cite{ref205}. For example, the ExaRanker model \cite{ref206} uses LLMs to generate explanations for neural rankers, which are then used to improve the rankers' performance during in-context learning.

While the use of post-hoc explanations to enhance in-context learning has shown promising results, there are still several challenges and limitations that need to be addressed. One key challenge is the reliability and faithfulness of the generated explanations \cite{ref207,ref208}. Researchers have found that LLMs can sometimes generate plausible-sounding explanations that may not accurately reflect the model's true decision-making process \cite{ref203}. This lack of faithfulness can undermine the effectiveness of the post-hoc explanations as a signal for in-context learning.

To address this challenge, researchers have proposed approaches that aim to improve the faithfulness of the generated explanations. For example, the xLLM framework \cite{ref209} introduces a generative explanation model that is optimized to maximize the faithfulness of the natural language explanations, ensuring that they accurately reflect the LLM's decision-making process.

Another limitation of the current approaches is that they primarily focus on the generation of explanations, without explicitly considering the influence of the explanations on the model's learning process. Researchers have argued that for post-hoc explanations to be truly effective in enhancing in-context learning, it is necessary to understand how the model utilizes the provided explanations and integrate them into the learning process in a more principled manner \cite{ref210}.

To address this, recent studies have explored the use of counterfactual reasoning and causal analysis to better understand the relationship between the generated explanations and the model's in-context learning performance \cite{ref196}. These approaches aim to identify the specific characteristics of the explanations that are most effective in improving the model's learning and generalization capabilities.

Looking ahead, the integration of post-hoc explanations and in-context learning is a promising area of research that has the potential to significantly enhance the capabilities of LLMs. By leveraging the insights provided by post-hoc explanations, researchers can develop more robust and reliable in-context learning algorithms that can adapt to a wider range of tasks and contexts. Additionally, the continued exploration of faithfulness and the relationship between explanations and learning can lead to the development of more transparent and trustworthy LLMs, which is crucial for their widespread adoption in real-world applications.

\subsection{Understanding LLMs' Logical Reasoning Capabilities}

The extent to which large language models (LLMs) truly understand logical rules versus learning a probabilistic mapping through context has been a subject of intense scrutiny and debate in the field of artificial intelligence. Recent studies have employed counterfactual methods to closely examine the logical reasoning capabilities of these models.

One key aspect of in-context learning that has been explored in the literature is the use of post-hoc explanations, or natural language explanations, generated by LLMs to provide insights into their decision-making process. These post-hoc explanations have been investigated as a means to enhance the in-context learning capabilities of LLMs. For example, the AMPLIFY framework leverages post-hoc explanation methods to automatically generate natural language rationales that can provide corrective signals to the LLM during in-context learning, leading to significant improvements in prediction accuracy \cite{ref195}. Similarly, the Self-AMPLIFY framework explores the use of post-hoc explanations generated by small language models to improve their own in-context learning performance, addressing the challenge of relying on human-annotated rationales \cite{ref204}.

However, the reliability and faithfulness of the generated explanations remain a key challenge, as LLMs can sometimes produce plausible-sounding explanations that do not accurately reflect their true decision-making process \cite{ref203}. To address this, researchers have proposed approaches like the xLLM framework, which aims to optimize the faithfulness of the generated natural language explanations \cite{ref209}.

Furthermore, recent studies have explored the use of counterfactual reasoning and causal analysis to better understand the relationship between the generated explanations and the model's in-context learning performance \cite{ref196}. These approaches aim to identify the specific characteristics of the explanations that are most effective in improving the model's learning and generalization capabilities.

One study, ``Do Large Language Models Understand Logic or Just Mimic Context'' \cite{ref211}, investigates the underlying mechanisms driving the success of LLMs in logical reasoning tasks. The authors found that LLMs do not truly understand logical rules; instead, they have learned a probabilistic mapping through the context provided in the training data. By modifying the context text or changing the concepts of logical terms, the authors were able to significantly disrupt the outputs of LLMs, leading to counterintuitive responses. This suggests that LLMs are not relying on a deep understanding of logical rules but rather on surface-level patterns in the data.

Building on this, the ``CLOMO: Counterfactual Logical Modification with Large Language Models'' \cite{ref212} paper introduces a novel task, Counterfactual Logical Modification (CLOMO), to further investigate the counterfactual reasoning capabilities of LLMs. The authors developed a human-annotated benchmark to assess the ability of LLMs to effectively alter a given argumentative text to uphold a predetermined logical relationship. Their experiments revealed that while LLMs demonstrate a notable capacity for logical counterfactual thinking, there remains a discernible gap between their current abilities and human performance.

These studies underscore the limitations of LLMs in truly understanding logical rules and causal relationships, and highlight the importance of incorporating more explicit reasoning mechanisms and structured knowledge to enhance their logical and causal reasoning capabilities.

\subsection{Efficiency Considerations for In-context Learning}

The efficiency of deploying large language models (LLMs) for in-context learning tasks has emerged as a critical challenge in the field of natural language processing. LLMs, which have demonstrated remarkable capabilities in a wide range of tasks, often require substantial computational resources and memory, making their practical deployment a significant hurdle.

One of the primary efficiency considerations in in-context learning with LLMs is the computational complexity associated with the models themselves. The attention mechanism, which is a core component of transformers and is widely used in LLMs, has a quadratic time and space complexity with respect to the input sequence length \cite{ref213}. This complexity can lead to significant performance bottlenecks, particularly when dealing with long input contexts, which are often necessary for effective in-context learning.

To address this challenge, researchers have explored various techniques for improving the efficiency of LLMs in the context of in-context learning. One such approach is model compression, which aims to reduce the size and complexity of LLMs while preserving their performance. Several model compression methods have been proposed, including quantization \cite{ref214} and pruning \cite{ref215}.

Quantization techniques, which involve reducing the bit-width of model parameters, have shown promising results in reducing the memory footprint and computational requirements of LLMs. By representing weights with fewer bits, quantized models can achieve significant memory savings and faster inference times, making them more suitable for deployment on resource-constrained devices \cite{ref214}.

Pruning, on the other hand, focuses on removing redundant connections or neurons within the model architecture, thereby reducing the overall model size and complexity. Recent advancements in structured pruning methods, such as the work presented in \cite{ref215}, have shown that LLMs can be significantly compressed while maintaining their performance on in-context learning tasks.

In addition to model compression techniques, researchers have also explored data-centric approaches to improve the efficiency of in-context learning. One such approach is the use of synthetic data generation, where fine-tuned LLMs are employed as ``teachers'' to generate high-quality synthetic training data for smaller downstream models \cite{ref216}. By leveraging the in-context learning capabilities of LLMs, this approach can help mitigate the need for large amounts of manually annotated data, thereby reducing the overall cost and effort required for in-context learning.

Furthermore, techniques like retrieval-based knowledge transfer \cite{ref217} have been proposed, where knowledge from large-scale LLMs is distilled into smaller models, enabling efficient in-context learning without the need for resource-intensive fine-tuning of the entire model.

Another promising direction for improving the efficiency of in-context learning with LLMs is the development of novel architectural designs that are specifically tailored for this task. For example, the work presented in \cite{ref218} introduces a structured attention mechanism that can effectively scale in-context demonstrations, achieving comparable or better performance than full attention while obtaining significant inference speed-ups.

Additionally, techniques like \cite{ref219} have explored methods for compressing the intermediate activation of a specified span of tokens, thereby reducing both memory and computational costs during in-context learning.

Overall, the efficiency challenges associated with deploying LLMs for in-context learning tasks are multifaceted, encompassing computational complexity, memory requirements, and the need for data-efficient approaches. Researchers have made significant strides in addressing these challenges through model compression techniques, data-centric methods, and novel architectural designs. As the field of in-context learning with LLMs continues to evolve, further advancements in these areas will be crucial for enabling the widespread and practical deployment of these powerful models in real-world applications.

\subsection{Complementary Explanations for In-context Learning}

Recent research has shown that incorporating natural language explanations into the prompting process can significantly enhance the performance of large language models (LLMs) in in-context learning tasks \cite{ref196,ref220}. These studies have provided valuable insights into the mechanisms by which explanations can be effectively leveraged to improve the learning capabilities of LLMs.

One key factor that contributes to the effectiveness of explanations in in-context learning is the computation trace, or the way in which the solution is decomposed and presented to the model \cite{ref196}. The natural language used to express the prompt has also been shown to play a crucial role in the effectiveness of explanations \cite{ref220}.

While the importance of explanations in in-context learning has been well-established, the field still lacks a comprehensive understanding of how to construct the most effective sets of explanations to optimize the performance of LLMs. To address this gap, researchers have proposed approaches for generating complementary explanations that can further enhance the in-context learning capabilities of LLMs \cite{ref196}.

The key idea behind the concept of complementary explanations is that diverse reasoning skills demonstrated by different exemplars can lead to better performance. By constructing a set of explanations that are both relevant and complementary, the LLM can be exposed to a wider range of reasoning strategies and approaches, allowing it to develop a more robust and versatile understanding of the task.

One such approach, proposed in \cite{ref196}, utilizes a maximal marginal relevance-based exemplar selection strategy to identify the optimal set of explanations. This method selects explanations that are both relevant to the task at hand and diverse in their reasoning patterns, ensuring that the LLM is presented with a comprehensive and complementary set of explanations.

The effectiveness of this approach has been demonstrated across a variety of in-context learning tasks, including question answering, natural language inference, and symbolic reasoning \cite{ref196}. The results show that the use of complementary explanations can lead to significant performance improvements compared to using a single, or even a set of non-complementary, explanations.

Moreover, the impact of complementary explanations has been found to be particularly pronounced in the context of larger LLMs, such as GPT-3 and InstructGPT \cite{ref196}. This suggests that as the complexity and capabilities of LLMs continue to grow, the need for effective explanation-based approaches to in-context learning will become increasingly crucial.

It is important to note that the effectiveness of explanations in in-context learning is not limited to the computation trace and natural language used to express the prompt. Other factors, such as the quality and relevance of the explanations themselves, as well as the way in which they are integrated into the prompting process, can also play a significant role \cite{ref221,ref222}.

Researchers have explored various strategies for optimizing the selection and integration of explanations, including the use of proxy metrics to evaluate the quality of explanations \cite{ref222}, and the development of frameworks for automatically generating and updating explanations based on the performance of the LLM \cite{ref223}.

By continuing to investigate the role of explanations in in-context learning and exploring new approaches for constructing and leveraging complementary explansets, the field can make significant strides towards developing more effective and reliable LLMs that can truly harness the power of natural language understanding and reasoning.

\subsection{Synthetic Data Generation for Fine-tuning}

The emergence of large language models (LLMs) has significantly advanced the field of natural language processing, with their remarkable abilities to perform a wide range of tasks. However, the training of these large models often requires vast amounts of data, which can be a significant challenge, especially in low-resource settings.

To address this challenge, researchers have explored the use of synthetic data generation as a means to improve the performance of downstream models in data-scarce scenarios. One promising approach is to leverage the capabilities of fine-tuned LLMs as teachers to generate high-quality synthetic training data for smaller models.

The underlying premise of this approach is that fine-tuned LLMs, with their deep understanding of the task at hand, can serve as effective generators of realistic and informative training examples. By tapping into the knowledge and generation abilities of these powerful models, researchers aim to augment the limited real-world data available, thereby enhancing the performance of smaller, more resource-efficient models.

One of the key advantages of this approach is its ability to bridge the gap between the vast knowledge captured by LLMs and the specific task requirements of downstream applications. \cite{ref224} found that simply fine-tuning LLMs on the available data often falls short in terms of generating high-quality synthetic data, as the models may struggle to capture the nuances and idiosyncrasies of the target domain.

By leveraging the fine-tuned LLMs as teachers, researchers have developed more sophisticated data generation techniques that can effectively mimic the distribution of the real-world data, while also introducing novel and diverse examples that can help expand the training capabilities of the downstream models.

For instance, \cite{ref216} explored the use of fine-tuned LLMs to generate synthetic training data for smaller models in low-resource settings. The authors found that the synthetic data, combined with the limited real-world data, can significantly improve the performance of the downstream models, sometimes even matching or exceeding the performance achieved with the full real-world dataset.

Similarly, \cite{ref225} proposed a method that uses a teacher LLM to enhance a small seed dataset by generating additional data points based on the model's mistakes during training. This approach effectively amplifies the signal from incorrectly predicted data points and reintegrates them into the training data, allowing the student model to focus on more challenging examples.

The success of these approaches highlights the potential of leveraging the capabilities of fine-tuned LLMs to address the data scarcity challenges faced by smaller models. By generating high-quality synthetic data, researchers can effectively expand the training data available to these models, enabling them to achieve better performance without the need for extensive data collection and annotation efforts.

Moreover, the ability to generate task-specific synthetic data can also have broader implications for the development of more robust and adaptable language models. By fine-tuning LLMs on diverse datasets and tasks, and then using them to generate synthetic data, researchers can explore the transfer of knowledge and skills across different domains, potentially leading to the creation of more versatile and generalized language models.

\cite{ref226} provides a comprehensive overview of the advancements in this field, outlining the various methodologies, evaluation techniques, and practical applications of using generative AI, particularly LLMs, for synthetic data generation. The authors also discuss the current limitations and potential future research directions, highlighting the promising trajectory of this approach.

Overall, the use of fine-tuned LLMs as teachers to generate synthetic training data for smaller downstream models has emerged as a powerful strategy for addressing the challenges of data scarcity in the field of natural language processing. By leveraging the vast knowledge and generation capabilities of these large models, researchers can significantly improve the performance of smaller, more resource-efficient models, paving the way for more accessible and widespread deployment of advanced language technologies.

\subsection{Reducing Structured Data for In-context Learning}

The rapid advancements in large language models (LLMs) have enabled their widespread adoption across diverse applications, including natural language understanding, generation, and reasoning. A key capability that has emerged in LLMs is their ability to perform in-context learning, where the model can adapt to new tasks or contexts with minimal examples or fine-tuning \cite{ref227,ref228}. This ability has been particularly useful in settings where access to labeled data is limited, as it allows LLMs to leverage the knowledge they have acquired during pre-training to quickly adapt to new tasks.

However, a significant challenge that arises in the context of in-context learning is the model's ability to effectively process and utilize the input context. As the input context grows in size and complexity, the computational and memory requirements of LLMs can quickly become prohibitive, limiting their practical deployment in real-world scenarios \cite{ref229}. To address this challenge, researchers have explored techniques for fine-tuning LLMs to generate reduced versions of input contexts, with the aim of improving the reasoning performance of a fixed LLM on downstream tasks.

One approach to this problem is presented in the paper ``Unveiling the Generalization Power of Fine-Tuned Large Language Models''. The authors investigate the impact of fine-tuning on the generalization ability of LLMs, and observe that fine-tuning on generation tasks can enhance the model's ability to generalize to different domains and tasks. Building on this insight, they suggest that fine-tuning LLMs to generate reduced versions of input contexts could potentially improve their reasoning performance on downstream tasks.

Similarly, the paper ``Self-Selected Attention Span for Accelerating Large Language Model Inference'' explores a related approach, where the authors fine-tune an LLM to identify the minimal attention spans required for each step of a task, and then use this information to optimize the inference process. By reducing the context that the model needs to attend to, the authors were able to achieve significant improvements in the throughput of LLM inference, without compromising the model's performance.

Another relevant work is ``Compressing Context to Enhance Inference Efficiency of Large Language Models'', which proposes a method called Selective Context that enhances the inference efficiency of LLMs by identifying and pruning redundancy in the input context. The authors demonstrate that their approach can significantly reduce memory cost and inference time while maintaining comparable performance on downstream tasks, such as summarization, question answering, and response generation.

These approaches share a common underlying principle: by fine-tuning LLMs to generate reduced versions of input contexts, the models can more effectively process and utilize the available information, leading to improved reasoning performance on downstream tasks. This is particularly important in settings where the input context is large and complex, as it can help alleviate the computational and memory constraints that often limit the deployment of LLMs in real-world applications.

One key advantage of these techniques is their potential for broader applicability. By focusing on the input context, rather than task-specific architectures or training procedures, these methods can be leveraged across a wide range of applications and domains, where the challenges of processing large and complex input contexts are prevalent. This flexibility could be especially valuable in scenarios where the deployment of LLMs is constrained by resource limitations, such as on edge devices or in low-power computing environments.

Moreover, the ability to generate reduced versions of input contexts could also have implications for the interpretability and transparency of LLM-based systems. By identifying the most relevant and informative subsets of the input context, these techniques could potentially aid in the development of more interpretable and explainable models, which is a critical requirement in many high-stakes applications, such as healthcare, finance, and legal decision-making \cite{ref230}.

\section{In-Context Learning in Computer Vision and Multimodal Settings}

\subsection{Leveraging Contextual Information in Computer Vision}

The importance of modeling and leveraging multimodal contextual information for in-context learning in computer vision has also been highlighted in recent research. Beyond the integration of textual and visual modalities, audio cues have also been explored to enhance in-context learning for computer vision tasks.

One key factor in enhancing the performance of in-context learning in computer vision is the effective modeling and utilization of multimodal contextual information. In contrast to traditional language-based ICL, visual ICL needs to go beyond just textual prompts and leverage various types of contextual cues, including visual, textual, and audio information, to enable more robust and versatile task adaptation.

A growing body of research has explored different techniques for incorporating multimodal contextual information to boost the in-context learning capabilities of computer vision models. \cite{ref148} proposed a generative transformer-based model that can in-context learn visual tasks from multimodal prompts, including textual descriptions and input-output visual examples. Similarly, \cite{ref132} investigated the significance of both visual and textual information in the demonstrations for in-context learning on vision-language models (VLMs), and proposed the Mixed Modality In-Context Example Selection (MMICES) approach to boost the ICL performance of VLMs.

The importance of modeling and leveraging multimodal contextual information for in-context learning in computer vision has also been highlighted in \cite{ref184}. Beyond the integration of textual and visual modalities, recent research has also explored the role of audio cues in enhancing in-context learning for computer vision tasks, as demonstrated in \cite{ref231}.

The effective modeling and utilization of multimodal contextual information for in-context learning in computer vision is a rapidly evolving research area. Researchers are exploring various techniques, including prompt engineering, multimodal feature fusion, and cross-modal attention mechanisms, to enable models to seamlessly integrate and leverage the complementary information from different modalities. One promising direction is the development of specialized architectural designs and learning paradigms, such as \cite{ref9}, which introduced the Instruct Me More (InMeMo) method to enhance the performance of in-context learning on visual tasks.

The development of comprehensive evaluation frameworks and benchmarks, like \cite{ref184}, is also crucial to assess the performance of in-context learning models in computer vision and multimodal settings.

\subsection{Multimodal Context Modeling}

The emergence of large language models (LLMs) \cite{ref232} has enabled in-context learning (ICL) to become a prominent approach in computer vision, allowing models to adapt to new tasks by leveraging a few demonstration examples. While ICL has been extensively explored in the language domain, its application in the visual domain poses unique challenges due to the inherent complexity and multimodal nature of visual information.

One of the key factors in enhancing the performance of in-context learning in computer vision is the effective modeling and utilization of multimodal contextual information. In contrast to traditional language-based ICL, visual ICL needs to go beyond just textual prompts and leverage various types of contextual cues, including visual, textual, and audio information, to enable more robust and versatile task adaptation.

A growing body of research has explored different techniques for incorporating multimodal contextual information to boost the in-context learning capabilities of computer vision models. \cite{ref148} proposed a generative transformer-based model that can in-context learn visual tasks from multimodal prompts, including textual descriptions and input-output visual examples. Similarly, \cite{ref132} investigated the significance of both visual and textual information in the demonstrations for in-context learning on vision-language models (VLMs), and proposed the Mixed Modality In-Context Example Selection (MMICES) approach to boost the ICL performance of VLMs.

The importance of modeling and leveraging multimodal contextual information for in-context learning in computer vision has also been highlighted in \cite{ref184}. Beyond the integration of textual and visual modalities, recent research has also explored the role of audio cues in enhancing in-context learning for computer vision tasks, as demonstrated in \cite{ref231}.

The effective modeling and utilization of multimodal contextual information for in-context learning in computer vision is a rapidly evolving research area. Researchers are exploring various techniques, including prompt engineering, multimodal feature fusion, and cross-modal attention mechanisms, to enable models to seamlessly integrate and leverage the complementary information from different modalities. One promising direction is the development of specialized architectural designs and learning paradigms, such as \cite{ref9}, which introduced the Instruct Me More (InMeMo) method to enhance the performance of in-context learning on visual tasks.

The development of comprehensive evaluation frameworks and benchmarks, like \cite{ref184}, is also crucial to assess the performance of in-context learning models in computer vision and multimodal settings.

\subsection{Prompt Engineering for Visual In-Context Learning}

Prompt engineering has emerged as a crucial technique in enabling effective in-context learning for computer vision tasks \cite{ref24}. The choice and design of the prompts can have a significant impact on the performance of in-context learning models, as they serve as the primary means of conveying task-specific information and guiding the model's adaptation to new scenarios.

One key aspect of prompt engineering for visual in-context learning is prompt selection \cite{ref233}. Given a set of potential prompts, the model needs to be able to select the most appropriate ones to provide the necessary context and guidance for a particular task or input. This selection process can be based on various factors, such as the similarity between the prompt and the task, the informativeness of the prompt, or the model's internal representation of the prompt.

Another important aspect is prompt composition \cite{ref233}. In-context learning models often require a combination of multiple prompts, either in the form of text-based instructions or visual examples, to effectively convey the task requirements. The way these prompts are composed and integrated can have a significant impact on the model's ability to understand and execute the task.

Recent studies have explored various approaches to prompt composition, including the use of soft prompts \cite{ref234}, prototype-based prompts \cite{ref235}, and hierarchical prompt structures \cite{ref236}. These methods aim to leverage the complementary information from different prompts and enhance the model's understanding of the task requirements.

Prompt fusion is another crucial aspect of prompt engineering for visual in-context learning \cite{ref148}. In-context learning models often need to integrate information from multiple prompts, such as text-based instructions and visual examples, to effectively solve a given task. The way these prompts are fused and the strategies used to combine them can significantly impact the model's performance.

Recent studies have explored various prompt fusion techniques, including attention-based fusion \cite{ref148}, retrieval-enhanced fusion \cite{ref237}, and hierarchical fusion \cite{ref238}. These methods aim to effectively leverage the complementary information from different prompts and enhance the model's ability to understand and execute the task.

Additionally, researchers have investigated the use of synthetic data and data augmentation techniques to improve prompt engineering for visual in-context learning \cite{ref233}. By generating or augmenting prompts, these approaches aim to expand the diversity of the prompt set and enhance the model's ability to generalize to new tasks and contexts.

Overall, prompt engineering plays a crucial role in enabling effective in-context learning for computer vision tasks. The selection, composition, and fusion of prompts, as well as the use of synthetic data and data augmentation, are key aspects that researchers have explored to enhance the performance and generalization of in-context learning models in the visual domain \cite{ref233}. As the field of visual in-context learning continues to evolve, further advancements in prompt engineering are expected to play a crucial role in unlocking the full potential of these models.

\subsection{Architectural Advances for Visual In-Context Learning}

The emergence of large language models (LLMs) \cite{ref239} and their impressive in-context learning (ICL) abilities have inspired researchers to explore similar capabilities in the visual domain. While LLMs can effectively leverage contextual information and few-shot demonstrations to adapt to new tasks, the same level of ICL performance has not been easily replicated in vision-based models. This has motivated the exploration of novel architectural designs and advancements to enable more effective in-context learning for visual tasks.

One key architectural development has been the incorporation of attention mechanisms to enhance the models' ability to capture and leverage relevant contextual information. Attention-based models, such as Vision Transformers (ViTs) \cite{ref240}, have demonstrated superior performance in various visual tasks compared to traditional convolutional neural networks (CNNs). The attention mechanisms in ViTs enable the models to dynamically focus on relevant regions of the input image, which can be particularly beneficial for in-context learning scenarios where the model needs to adapt to new tasks or environments based on the provided demonstrations.

Researchers have further explored ways to improve the attention mechanisms in ViTs to better suit the needs of visual in-context learning. For example, \cite{ref11} proposed a unified model that encodes both text-based prompts and visual inputs into a shared representation space, allowing the attention mechanism to seamlessly integrate and leverage the contextual information from both modalities. Similarly, \cite{ref241} investigated the transfer of in-context learning capabilities from large language models to vision-language models, highlighting the importance of cross-modal attention mechanisms in enabling effective in-context learning in the visual domain.

Beyond attention-based architectures, the incorporation of meta-learning techniques has also shown promise for enhancing visual in-context learning. Meta-learning approaches, such as Model-Agnostic Meta-Learning (MAML) \cite{ref242}, aim to learn a model initialization that can be quickly adapted to new tasks or environments with limited data. In the context of visual in-context learning, meta-learning techniques can enable the model to rapidly adapt to new tasks based on the provided demonstrations, without requiring extensive fine-tuning.

Furthermore, researchers have explored the integration of cross-modal interactions to improve in-context learning in the visual domain. \cite{ref243} investigated the effectiveness of cross-attention mechanisms, which capture the interactions between different modalities (e.g., vision and language), compared to self-attention mechanisms that operate within a single modality. The findings suggest that cross-attention can be a valuable tool for enhancing the models' ability to leverage contextual information from multiple modalities during in-context learning.

Additionally, \cite{ref244} proposed a method that tightly couples the vision and language representations during prompt tuning, enabling the model to effectively leverage the interplay between the two modalities for improved few-shot generalization and domain adaptation. This approach highlights the importance of explicitly modeling the cross-modal interactions to enable more effective in-context learning in the visual domain.

Another important aspect of architectural advancements for visual in-context learning is the ability to handle varying levels of contextual information and task complexity. \cite{ref130} explored techniques to enhance the models' capability to adapt to complex multi-modal prompts, including the retrieval of relevant images, summarization of visual information, and composition of language-based demonstrations. By addressing these challenges, the proposed methods enabled large vision-language models to more effectively leverage in-context learning for a wide range of visual tasks.

Overall, the architectural advancements in the field of visual in-context learning reflect the growing recognition of the importance of contextual information, cross-modal interactions, and meta-learning capabilities. These developments have enabled more effective adaptation and generalization of visual models to new tasks and environments, paving the way for more robust and versatile in-context learning in the visual domain.

\subsection{Benchmarking and Evaluation of Visual In-Context Learning}

Developing robust and comprehensive evaluation protocols for assessing the performance of in-context learning models in computer vision is crucial for advancing the field. Existing benchmarks and evaluation approaches often suffer from various limitations, highlighting the need for more rigorous and standardized evaluation frameworks.

One of the primary challenges in benchmarking visual in-context learning is the lack of standardized datasets and evaluation protocols. While several benchmarks, such as IGLUE \cite{ref245} and MMBench \cite{ref246}, have been proposed to evaluate multimodal language models, their focus is often on vision-language tasks, which may not fully capture the unique challenges of in-context learning in the visual domain.

To address this gap, researchers have proposed specialized benchmarks for evaluating in-context learning in computer vision. For instance, \cite{ref9} introduced a benchmark that focuses on foreground segmentation and single object detection tasks, using in-context image pairs as prompts. Similarly, \cite{ref247} presented a benchmark covering a range of vision tasks, including high-level visual understanding and low-level image processing, to assess the in-context learning capabilities of their proposed Painter model.

While these benchmarks have made notable contributions, they often target specific tasks or use limited datasets, which may not be representative of the diverse range of visual in-context learning challenges. To overcome this limitation, researchers have called for the development of more comprehensive and cross-domain benchmarks that can capture the varying complexities and requirements of in-context learning in computer vision.

One such effort is the BLINK benchmark \cite{ref248}, which focuses on core visual perception abilities, such as relative depth estimation, visual correspondence, and multi-view reasoning. BLINK reformats classic computer vision tasks into a multiple-choice question format, allowing for a more rigorous and standardized evaluation of multimodal language models' in-context learning capabilities in the visual domain. The benchmark's focus on perception-demanding tasks highlights the limitations of current multimodal language models, which often rely on language-mediated understanding rather than genuine visual reasoning.

In addition to task-specific benchmarks, researchers have also explored the development of more general-purpose evaluation frameworks for visual in-context learning. For example, \cite{ref249} introduces a dataset based on the game of chess, with a curated rule set that constrains the set of allowable predictions. This dataset aims to probe the semantic and logical reasoning abilities of vision models, providing insights into their in-context learning performance beyond standard accuracy-based metrics.

Another important aspect of benchmarking visual in-context learning is the consideration of contextual information and its role in model performance. The \cite{ref184} benchmark specifically focuses on evaluating the ability of multimodal language models to perform context-sensitive text-rich visual reasoning, highlighting the importance of incorporating contextual cues in the evaluation process.

Alongside dataset-specific benchmarks, researchers have also explored the development of more general evaluation protocols and metrics for assessing in-context learning performance. \cite{ref250} introduces the ContextRef benchmark, which examines the alignment between referenceless metrics (such as CLIPScore) and human preference judgments in the context of image description generation. This work emphasizes the importance of considering context dependence in the evaluation of in-context learning models.

Despite these advancements, the field of benchmarking and evaluating visual in-context learning remains an active area of research, with several ongoing challenges and limitations. One key issue is the sensitivity of in-context learning performance to prompt design, which can significantly impact model behavior and lead to inconsistent evaluation results \cite{ref251}. Addressing this challenge requires the development of standardized prompt engineering guidelines and evaluation protocols that can account for the influence of prompt design on model performance.

Another limitation is the potential for biases and inconsistencies in the evaluation datasets themselves. As highlighted in \cite{ref251}, the definitions and representations of visual concepts can vary across datasets and annotators, leading to ambiguity and inconsistency in the evaluation process. Addressing these issues requires a more rigorous and systematic approach to dataset curation and annotation, as well as the development of evaluation metrics that can account for these nuances.

Furthermore, the scalability and efficiency of in-context learning evaluation protocols remain a concern. Many existing benchmarks rely on human-in-the-loop evaluation or costly fine-tuning procedures, which can limit their practical applicability and hinder the rapid development and iteration of in-context learning models. Addressing these challenges may require the exploration of more automated and efficient evaluation approaches, such as those proposed in \cite{ref252}.

In conclusion, the benchmarking and evaluation of visual in-context learning is a critical area of research that requires continued attention and innovation. Developing comprehensive, standardized, and scalable evaluation frameworks will be crucial for driving progress in this field, enabling researchers to accurately assess the capabilities of in-context learning models and identify areas for further improvement. By overcoming the current limitations and addressing the emerging challenges, the computer vision community can build a more robust and reliable foundation for advancing the state-of-the-art in visual in-context learning.

\subsection{Specialized Applications of Visual In-Context Learning}

Medical imaging and autonomous driving have emerged as specialized domains where in-context learning can provide significant benefits. In the field of medical imaging, the application of in-context learning has demonstrated great potential \cite{ref253,ref254}.

The scarcity of annotated medical imaging data poses a significant challenge for traditional supervised learning approaches. In-context learning can help overcome this limitation by enabling rapid adaptation to new tasks and datasets with minimal additional training. For example, \cite{ref253} has shown that the Generative Pretrained Transformer 4 with Vision capabilities (GPT-4V) can effectively classify cancer histopathology images by leveraging in-context learning, without the need for extensive task-specific fine-tuning.

One unique aspect of in-context learning in medical imaging is the potential to leverage pre-existing domain knowledge and privileged information to enhance generalization. \cite{ref254} demonstrates that using privileged information, such as tumor shape or location, can lead to stronger domain generalization ability compared to conventional techniques. This is particularly relevant in medical imaging, where domain shifts can occur due to differences in imaging modalities, hospital settings, or patient populations.

In the field of autonomous driving, in-context learning can play a crucial role in enabling robots and vehicles to adapt to diverse and dynamic environments \cite{ref255,ref256}. Traditional computer vision approaches often struggle with the complexity and variability of real-world driving scenarios, which can lead to performance degradation and safety concerns.

In-context learning offers a promising solution by allowing autonomous driving systems to rapidly adapt to new situations and tasks, such as object detection, semantic segmentation, or lane detection, based on a few illustrative examples \cite{ref257,ref258}. This flexibility is crucial for ensuring the robustness and generalization of autonomous driving systems, as they need to operate in a wide range of conditions, including varying weather, lighting, and traffic patterns.

Moreover, the integration of in-context learning with other techniques, such as multimodal sensing and diffusion models, can further enhance the capabilities of autonomous driving systems. \cite{ref256} presents a framework that leverages multimodal inputs, including depth maps and text captions, to adaptively enhance low-light images, improving the performance of object detection in challenging nighttime conditions.

In the domain of robotic perception, in-context learning can enable robots to rapidly learn new tasks and adapt to diverse environments through interactive human-robot collaboration \cite{ref258,ref259}. By incorporating human guidance and feedback, robots can expand their knowledge and capabilities beyond the initial training data, allowing them to tackle novel challenges and operate in unstructured settings.

The ability to learn from a few examples through in-context learning is particularly valuable in robotic applications, where data collection and annotation can be time-consuming and labor-intensive. By leveraging human interaction and task demonstrations, robots can efficiently acquire new skills and adapt to changing environments, ultimately enhancing their overall versatility and performance.

In summary, the application of in-context learning in specialized domains such as medical imaging and autonomous driving has demonstrated significant potential. The unique challenges and opportunities in these domains, such as data scarcity, domain shifts, and the need for rapid adaptation, make in-context learning a promising approach for enhancing the robustness, generalization, and adaptability of vision-based systems. As the field continues to evolve, further advancements in in-context learning for specialized applications will likely play a crucial role in driving the progress of these critical domains.

\subsection{Limitations and Future Directions of Visual In-Context Learning}

While the recent advancements in in-context learning (ICL) for computer vision have shown promising results, there are several key limitations and challenges that need to be addressed to fully realize the potential of this approach.

One of the primary limitations of visual in-context learning is its sensitivity to prompt design \cite{ref148,ref24}. The performance of ICL models can vary significantly depending on the choice and structure of the prompts, which often require extensive manual tuning and experimentation. This sensitivity makes it difficult to deploy these models in real-world scenarios where the input data and user requirements may be highly diverse and unpredictable. Researchers need to explore more robust and generalizable prompt engineering techniques to reduce the reliance on manual prompt design.

Another critical challenge is the susceptibility of visual in-context learning to biases \cite{ref260}. The performance of these models can be heavily influenced by biases present in the training data, such as dataset biases, social biases, and algorithmic biases. These biases can get amplified through the in-context learning process, leading to unfair and unreliable predictions, especially for underrepresented or marginalized groups. Developing debiasing strategies, including data curation, model design, and evaluation protocols, is crucial to ensuring the fairness and trustworthiness of visual in-context learning systems.

The need for scalable and efficient algorithms is another key limitation of current visual in-context learning approaches \cite{ref148,ref24}. As the complexity of tasks and the size of input data increase, the computational and memory requirements of in-context learning can become prohibitive, limiting their practical deployment. Researchers need to explore novel architectural designs, efficient prompt generation techniques, and distributed computing strategies to enable the scalable and efficient application of visual in-context learning in real-world scenarios.

In terms of future research directions, one promising area is the integration of visual in-context learning with other learning paradigms, such as transfer learning, meta-learning, and multi-task learning \cite{ref232}. By leveraging the synergies between these approaches, researchers can develop more adaptable, sample-efficient, and generalized vision systems that can rapidly adapt to new tasks and domains.

Another emerging trend is the exploration of novel architectural designs and attention mechanisms that are tailored specifically for visual in-context learning \cite{ref247,ref261}. These architectural innovations can enhance the models' ability to effectively leverage contextual information, including multi-modal, visual-semantic, and temporal cues, to improve their in-context learning performance.

The potential for cross-modal and cross-lingual transfer in visual in-context learning is another exciting research direction \cite{ref232}. Developing models that can seamlessly transfer their in-context learning capabilities across different modalities, such as from language to vision, and across different languages, can greatly expand the accessibility and versatility of these systems, enabling their deployment in diverse real-world applications.

\section{In-Context Learning in Specialized Domains}

\subsection{In-Context Learning in Scientific Computing}

In-context learning (ICL) has shown great promise in enhancing the performance and sample efficiency of various machine learning tasks, including natural language processing and computer vision. Recently, researchers have also explored the application of ICL in the domain of scientific computing, which encompasses a wide range of tasks such as numerical optimization, differential equation solving, and physics-based simulations.

One of the key advantages of ICL in scientific computing is its potential to reduce the computational burden of traditional simulation-based approaches. Many scientific computing tasks, such as solving partial differential equations (PDEs) or optimizing complex engineering systems, require extensive numerical simulations, which can be computationally expensive and time-consuming. By leveraging ICL, researchers aim to develop data-driven surrogate models that can quickly approximate the solutions to these complex problems, without the need for costly simulations \cite{ref262}.

In the context of numerical optimization, ICL has been explored as a means of accelerating the optimization process by conditioning the model on a few examples of optimal solutions or promising search directions. For instance, \cite{ref32} have demonstrated the feasibility of using sequence-to-sequence (Seq2Seq) models for few-shot learning in optimization tasks, outperforming traditional optimization algorithms in certain settings.

Similarly, in the domain of differential equation solving, ICL has been investigated as a way to enhance the performance and generalization capabilities of data-driven PDE solvers. \cite{ref262} and \cite{ref263} have proposed novel neural network architectures and pretraining strategies that leverage ICL to effectively solve a wide range of PDE systems, including those with complex boundary conditions and parametric dependencies.

One of the unique challenges in applying ICL to scientific computing tasks is the need to ensure that the model's predictions adhere to the underlying physical laws and constraints. Unlike natural language or image-based tasks, where the model can learn from a diverse set of examples, scientific computing tasks often require the model to respect the fundamental principles of physics, chemistry, or engineering \cite{ref264}. To address this challenge, researchers have explored the integration of physics-informed neural networks (PINNs) with ICL, where the model's architecture and training process are designed to enforce the relevant physical constraints \cite{ref262}, \cite{ref263}.

Another key challenge in scientific computing is the need to handle high-dimensional and complex data, such as multi-dimensional fields or unstructured grids. To address this, researchers have explored the use of advanced neural network architectures, such as convolutional neural networks and graph neural networks, in combination with ICL \cite{ref265}, \cite{ref266}. These approaches aim to leverage the model's ability to extract and leverage relevant features from the input data, while also benefiting from the adaptability and generalization capabilities of ICL.

Recent advancements in hardware accelerators, such as Tensor Processing Units (TPUs), have also played a crucial role in enabling the efficient deployment of ICL-based models in scientific computing \cite{ref267}. By harnessing the parallel processing capabilities of these specialized hardware, researchers have demonstrated significant speedups in scientific simulations and optimization tasks compared to traditional CPU-based approaches.

Looking ahead, the integration of ICL with scientific computing holds great promise for revolutionizing the way complex scientific and engineering problems are approached. By leveraging the adaptability and generalization capabilities of ICL, researchers can develop more efficient and versatile tools for tasks such as parameter optimization, inverse problem solving, and real-time decision-making in various domains, including fluid dynamics, material science, and climate modeling. As the field continues to evolve, we can expect to see further advancements in the theoretical understanding of ICL, the development of novel architectures and training strategies, and the successful deployment of ICL-based models in a wide range of scientific computing applications.

\subsection{In-Context Learning in Reinforcement Learning}

In-context learning (ICL) has shown great promise in enhancing the performance and sample efficiency of various machine learning tasks, including natural language processing and computer vision. Recently, researchers have also explored the application of ICL in the domain of scientific computing, which encompasses a wide range of tasks such as numerical optimization, differential equation solving, and physics-based simulations.

One of the key advantages of ICL in scientific computing is its potential to reduce the computational burden of traditional simulation-based approaches. Many scientific computing tasks, such as solving partial differential equations (PDEs) or optimizing complex engineering systems, require extensive numerical simulations, which can be computationally expensive and time-consuming. By leveraging ICL, researchers aim to develop data-driven surrogate models that can quickly approximate the solutions to these complex problems, without the need for costly simulations \cite{ref262}.

In the context of numerical optimization, ICL has been explored as a means of accelerating the optimization process by conditioning the model on a few examples of optimal solutions or promising search directions. For instance, \cite{ref32} have demonstrated the feasibility of using sequence-to-sequence (Seq2Seq) models for few-shot learning in optimization tasks, outperforming traditional optimization algorithms in certain settings.

Similarly, in the domain of differential equation solving, ICL has been investigated as a way to enhance the performance and generalization capabilities of data-driven PDE solvers. \cite{ref262} and \cite{ref263} have proposed novel neural network architectures and pretraining strategies that leverage ICL to effectively solve a wide range of PDE systems, including those with complex boundary conditions and parametric dependencies.

One of the unique challenges in applying ICL to scientific computing tasks is the need to ensure that the model's predictions adhere to the underlying physical laws and constraints. Unlike natural language or image-based tasks, where the model can learn from a diverse set of examples, scientific computing tasks often require the model to respect the fundamental principles of physics, chemistry, or engineering \cite{ref264}. To address this challenge, researchers have explored the integration of physics-informed neural networks (PINNs) with ICL, where the model's architecture and training process are designed to enforce the relevant physical constraints \cite{ref262}, \cite{ref263}.

Another key challenge in scientific computing is the need to handle high-dimensional and complex data, such as multi-dimensional fields or unstructured grids. To address this, researchers have explored the use of advanced neural network architectures, such as convolutional neural networks and graph neural networks, in combination with ICL \cite{ref265}, \cite{ref266}. These approaches aim to leverage the model's ability to extract and leverage relevant features from the input data, while also benefiting from the adaptability and generalization capabilities of ICL.

Recent advancements in hardware accelerators, such as Tensor Processing Units (TPUs), have also played a crucial role in enabling the efficient deployment of ICL-based models in scientific computing \cite{ref267}. By harnessing the parallel processing capabilities of these specialized hardware, researchers have demonstrated significant speedups in scientific simulations and optimization tasks compared to traditional CPU-based approaches.

Looking ahead, the integration of ICL with scientific computing holds great promise for revolutionizing the way complex scientific and engineering problems are approached. By leveraging the adaptability and generalization capabilities of ICL, researchers can develop more efficient and versatile tools for tasks such as parameter optimization, inverse problem solving, and real-time decision-making in various domains, including fluid dynamics, material science, and climate modeling. As the field continues to evolve, we can expect to see further advancements in the theoretical understanding of ICL, the development of novel architectures and training strategies, and the successful deployment of ICL-based models in a wide range of scientific computing applications.

\subsection{In-Context Learning in Speech Processing}

The use of in-context learning in speech processing tasks, such as speech recognition, speech synthesis, and spoken language understanding, has gained significant attention in recent years. The multimodal nature of speech data and the temporal dependencies inherent in speech signals pose unique challenges in adapting in-context learning techniques to these domains.

One of the key opportunities in leveraging in-context learning for speech processing lies in the potential to improve the performance and robustness of speech recognition systems. Large language models (LLMs) have exhibited impressive in-context learning capabilities in natural language tasks \cite{ref268}, and these capabilities can potentially be extended to speech recognition. By providing the LLM with a few examples of speech-text pairs as context, the model may be able to adapt its language model to perform better on the recognition of unseen speech \cite{ref269}. This could be particularly beneficial in scenarios where the speech data exhibits significant variation, such as due to accents, background noise, or speaker characteristics.

Moreover, in-context learning can also be explored in the context of speech synthesis, where a language model could be conditioned on a few examples of the desired speaking style or voice characteristics to generate high-quality synthetic speech \cite{ref270}. This could enable more personalized and expressive text-to-speech systems, catering to the specific needs and preferences of users.

Another promising application of in-context learning in speech processing is in the area of spoken language understanding (SLU). SLU tasks, such as intent recognition and slot filling, require the model to extract semantic information from speech input. By incorporating in-context learning, the SLU model could be quickly adapted to new domains or tasks using a few examples, without the need for extensive fine-tuning \cite{ref271}. This could greatly enhance the flexibility and deployment of SLU systems in real-world conversational interfaces.

However, adapting in-context learning techniques to speech processing tasks also presents unique challenges. The multimodal nature of speech, which combines acoustic, linguistic, and paralinguistic information, requires the model to effectively capture and leverage these different modalities \cite{ref272}. Additionally, the temporal dependencies inherent in speech signals, such as the sequential nature of phonemes and the prosodic features, need to be accounted for in the in-context learning approach \cite{ref273}.

Recent studies have explored various strategies to address these challenges. For instance, \cite{ref274} proposed a speech-text dialog pre-training model that explicitly aligns the speech and text modalities, enabling the model to better leverage contextual information for spoken language understanding tasks. Similarly, \cite{ref275} introduced a contextual audio-visual switching component that dynamically utilizes both audio and visual cues to enhance speech recognition performance in varying noise conditions.

Another key aspect to consider in the context of in-context learning for speech processing is the role of unsupervised speech representation learning. Techniques such as contrastive learning \cite{ref270} and self-supervised pre-training \cite{ref276} have shown promising results in extracting robust and task-agnostic speech representations. These representations can potentially serve as a strong foundation for in-context learning, enabling the model to quickly adapt to new speech-based tasks with limited data.

Furthermore, the integration of in-context learning with established speech processing techniques, such as acoustic modeling, language modeling, and speaker adaptation, could lead to synergistic advancements. For example, \cite{ref277} demonstrated the benefits of incorporating contextual information, such as time and location, to improve the performance of an RNN-T-based automatic speech recognition system.

In conclusion, the application of in-context learning in speech processing tasks presents both opportunities and challenges. By effectively leveraging the multimodal and temporal nature of speech data, and seamlessly integrating in-context learning with existing speech processing techniques, researchers and practitioners can unlock new possibilities in improving the flexibility, robustness, and personalization of speech-based systems. As the field continues to evolve, further advancements in this direction could significantly enhance the user experience and accessibility of spoken language technologies.

\section{Evaluation and Benchmarking of In-Context Learning}

\subsection{Benchmarks and Datasets for In-Context Learning Evaluation}

The rapid development of large language models (LLMs) has led to a surge of interest in the field of in-context learning, where these models demonstrate the ability to adapt to new tasks by leveraging a small number of demonstration examples within the input prompt \cite{ref278,ref4}. To facilitate the systematic evaluation of in-context learning capabilities, researchers have introduced a variety of benchmark datasets.

These benchmark datasets often exhibit distinct characteristics that can influence the performance of in-context learning models. One key aspect is the complexity of the tasks, which can range from relatively simple classification problems to more sophisticated reasoning and multi-step tasks \cite{ref279,ref280}. The linguistic diversity of the datasets is another important factor, as models may struggle with tasks involving low-resource or typologically distant languages \cite{ref281}.

Furthermore, the domain coverage of the benchmark datasets is crucial, as in-context learning performance can be heavily influenced by the relevance of the demonstration examples to the target task \cite{ref282}. Some benchmarks, such as the IGLUE (Image-Grounded Language Understanding Evaluation) \cite{ref245}, incorporate multimodal tasks that require models to jointly process visual and textual information, adding another layer of complexity to the in-context learning challenge.

The evaluation of in-context learning also needs to consider the format and composition of the prompts used to present the demonstration examples to the models \cite{ref48,ref15}. Factors such as the number of examples, their ordering, and the way they are formatted can significantly impact the performance of in-context learning models, as these aspects can influence the models' ability to effectively extract and apply the relevant information from the demonstrations.

To address these complexities, researchers have proposed more comprehensive evaluation frameworks, such as the ICL Consistency Test \cite{ref283}, which assesses the stability and robustness of in-context learning across multiple setups and prompts. Additionally, the CausalBench \cite{ref284} benchmark specifically focuses on evaluating the causal reasoning capabilities of LLMs in the in-context learning setting, providing insights into their ability to understand and leverage causal relationships within the task prompts.

\subsection{Evaluation Metrics for In-Context Learning}

Evaluating the performance of in-context learning models is a crucial aspect of understanding their capabilities and limitations. Various metrics have been proposed to assess the effectiveness of these models, each with its own strengths and weaknesses. In this subsection, we explore the different evaluation metrics that have been used in the context of in-context learning and discuss their suitability for different scenarios.

One of the most commonly used metrics in the field of in-context learning is few-shot accuracy \cite{ref285}. This metric measures the model's ability to correctly classify or predict the output for a small number of examples (e.g., 1-shot, 5-shot) provided in the context. Few-shot accuracy is a straightforward and intuitive metric that aligns well with the goal of in-context learning, which is to leverage a limited number of examples to adapt to a new task or domain. However, this metric may not capture the full range of a model's capabilities, as it focuses solely on the binary correctness of the predictions and does not provide insights into the model's ability to handle more nuanced or complex scenarios.

Another metric that has been employed in the context of in-context learning is average precision \cite{ref286}. This metric measures the area under the precision-recall curve, providing a more comprehensive assessment of the model's performance across different thresholds. Average precision can be particularly useful in scenarios where the task involves ranking or scoring multiple options, as it captures the model's ability to correctly order the predictions. However, the interpretation of average precision can be less intuitive than that of few-shot accuracy, and it may be more sensitive to the specific task and data distribution.

Normalized discounted cumulative gain (nDCG) \cite{ref287} is another metric that has been explored in the context of in-context learning. nDCG measures the quality of a model's ranking of relevant items, considering both the relevance of the items and their position in the ranking. This metric can be particularly useful in scenarios where the model is required to generate a ranked list of outputs, such as in question-answering or information retrieval tasks. However, the normalization process in nDCG can introduce some challenges, as it can lead to inconsistencies in the relative ranking of models \cite{ref287}.

In addition to these general-purpose metrics, task-specific evaluation metrics have also been used to assess the performance of in-context learning models \cite{ref288}. For example, in the context of text summarization, metrics such as ROUGE, BERTScore, and factual consistency have been employed to evaluate the quality of the generated summaries. These task-specific metrics can provide more nuanced and informative assessments of a model's performance, as they are tailored to the specific characteristics and requirements of the task at hand. However, the use of task-specific metrics may also limit the comparability of results across different tasks or domains.

When selecting appropriate evaluation metrics for in-context learning, it is important to consider the specific characteristics of the task, the desired performance criteria, and the intended use case of the model. For example, in scenarios where the model is required to provide a ranked list of predictions, metrics like nDCG may be more suitable than few-shot accuracy. Conversely, in binary classification tasks, few-shot accuracy may be a more appropriate choice.

Additionally, it is crucial to address the potential limitations and biases of the evaluation metrics. As highlighted in \cite{ref286}, the reuse of validation sets and the use of small datasets can lead to unreliable performance estimates, even with seemingly robust metrics like few-shot accuracy. To address these challenges, researchers have proposed stronger random baselines \cite{ref50} and have emphasized the importance of comprehensive and diverse evaluation frameworks \cite{ref289}.

In conclusion, the evaluation of in-context learning models is a multifaceted challenge that requires careful consideration of the appropriate metrics, task characteristics, and potential biases. The choice of evaluation metrics should be guided by the specific goals and requirements of the in-context learning application, and researchers should strive to develop more comprehensive and robust evaluation frameworks to gain a deeper understanding of the capabilities and limitations of these models.

\subsection{Standardized Evaluation Protocols and Leaderboards}

The development of standardized evaluation protocols and leaderboards has been a crucial step in advancing the field of in-context learning (ICL). Researchers have recognized the importance of having a consistent and reproducible evaluation framework to accurately and reliably assess the performance of ICL models.

One of the key challenges in achieving standardized evaluation protocols for ICL is the sensitivity of model performance to the design of the prompt \cite{ref51}. The prompt, which includes the input examples and the task description, can have a significant impact on the model's ability to learn and perform well in the in-context setting. This sensitivity to prompt design can lead to inconsistencies in model evaluation, as different prompts may result in vastly different performance outcomes.

To address this challenge, researchers have proposed the development of standardized evaluation protocols and leaderboards that aim to reduce the impact of prompt design on model evaluation. \cite{ref290} highlights the need for leaderboard systems that go beyond a single, static test set and instead focus on evaluating models in a real-world setting. The authors argue that current leaderboard systems may not accurately reflect a model's performance in practical applications, as they often fail to capture the discrepancy between testing and real-world conditions.

In response to these concerns, researchers have started to explore the development of more comprehensive and robust evaluation frameworks for ICL. \cite{ref291} proposes a framework that incorporates a metric to evaluate the predictive bias of a fixed prompt against labels or a given attribute. By identifying prompts with higher bias, this approach aims to improve the performance of ICL models in an effective and interpretable manner.

Another key challenge in developing standardized evaluation protocols for ICL is the issue of dataset biases. \cite{ref292} highlights how causal language models (CausalLMs) can be more sensitive to the order of in-context demonstration examples compared to prefix language models (PrefixLMs). This sensitivity can be attributed to the auto-regressive attention masks within CausalLMs, which restrict each token from accessing information from subsequent tokens. The authors propose an unsupervised fine-tuning method, termed the Information-Augmented and Consistency-Enhanced approach, to address this challenge and reduce the sensitivity to the order of in-context examples.

In addition to prompt design and dataset biases, the stability and discriminative power of evaluation metrics have also been a topic of discussion in the ICL research community. \cite{ref293} emphasizes the importance of conducting significance testing when evaluating the performance of ICL models, as the rapid progress in the field can lead to frequent replacement of the ``state of the art'' (SOTA) model on leaderboards. The authors propose an evaluation framework that explicitly treats certain outcomes as distinct and avoids aggregating them into a single-point metric, which can provide more nuanced insights into the performance of ICL models.

The development of standardized evaluation protocols and leaderboards for ICL is not limited to natural language processing tasks. \cite{ref294} highlights the need for robust evaluation frameworks in the context of conversational recommender systems, where the quality and diversity of the generated training data are critical for the performance of the model. The authors propose a multi-faceted approach to evaluating the quality of synthetic datasets produced by generative models, addressing the limitations of current evaluation methods.

Furthermore, the research community has also recognized the importance of considering the utility of ICL models from the perspective of end-users, rather than solely focusing on performance metrics. \cite{ref295} argues that current leaderboard systems may not accurately capture the utility of ICL models for practitioners, as they often prioritize performance-based evaluation over other qualities, such as compactness, fairness, and energy efficiency. The authors advocate for more transparency on leaderboards, including the reporting of statistics that are of practical concern, to better estimate a model's utility for end-users.

In summary, the development of standardized evaluation protocols and leaderboards for ICL is an active area of research, with researchers addressing various challenges, such as the sensitivity to prompt design, dataset biases, and the need for more comprehensive and robust evaluation frameworks. The goal is to create evaluation protocols that can reliably assess the performance of ICL models, while also considering the practical utility of these models for end-users. As the field of ICL continues to evolve, the research community will need to remain vigilant in developing and refining these evaluation frameworks to ensure the reliable and responsible development of ICL technologies.

\subsection{Comparative Analyses of In-Context Learning Models}

Comparative analyses of in-context learning (ICL) capabilities across different language models have been an active area of research, shedding light on the factors that contribute to the emergence and performance of this remarkable ability. These studies have examined the sensitivity of ICL to model architecture, size, and pre-training data, offering valuable insights for future model development.

One key finding from comparative analyses is that the emergence of ICL is not strictly tied to model size. Several studies have demonstrated that smaller models, when trained on appropriately designed pre-training data, can also exhibit strong ICL abilities \cite{ref296}. This suggests that the specific characteristics of the pre-training data, rather than just the model scale, play a crucial role in enabling ICL.

For instance, \cite{ref297} found that the presence of ``parallel structures'' in the pre-training data, where pairs of phrases follow similar templates in the same context window, is a key factor that contributes to the emergence of ICL. The authors demonstrate that removing these parallel structures from the pre-training data significantly reduces the ICL accuracy of the resulting models, even when excluding common patterns like n-gram repetitions and long-range dependencies. This highlights the importance of the underlying structure and diversity of the pre-training data in shaping the ICL capabilities of language models.

Similarly, \cite{ref3} proposes a theoretical framework that attributes the emergence of ICL to the pre-training distribution having long-range coherence. In this setting, the language model must infer a latent document-level concept during pre-training to generate coherent next tokens. The authors show that at test time, ICL occurs when the model can also infer a shared latent concept between the examples in the prompt, despite a distribution mismatch between the prompt and pre-training data.

Comparative analyses have also shed light on the interplay between model architecture and ICL performance. \cite{ref298} evaluates thirteen different model architectures, including recurrent, convolutional, and attention-based models, on a suite of synthetic ICL tasks. The authors find that all the considered architectures can perform ICL under a wider range of conditions than previously documented, but with stark differences in statistical efficiency and consistency. Interestingly, they also observe that several attention-alternative architectures are sometimes competitive with or better in-context learners than transformer-based models, suggesting that the attention mechanism may not be the sole driver of ICL abilities.

Furthermore, comparative studies have explored the relationship between language modeling performance, measured by perplexity, and ICL capabilities. \cite{ref299} examines how a small subset of the pre-training data that is particularly supportive of ICL can be identified and used to significantly improve the model's ICL ability, by up to 18\%. Interestingly, the authors find that this supportive subset does not necessarily have a higher domain relevance to the downstream tasks, but rather contains a higher mass of rarely occurring, long-tail tokens, indicating that learning to incorporate difficult long-range context may be a key factor in enabling effective ICL.

Similarly, \cite{ref300} delves into the interplay between model architecture, layer depth, and multilingual processing capabilities of language models. The authors uncover the phenomenon of ``over-layerization,'' where increasing layer depth without corresponding adjustments to other parameters can degrade model performance, suggesting that the balance between different architectural components is crucial for leveraging contextual information, including for ICL.

Taken together, the comparative analyses of ICL capabilities across different language models have revealed that the emergence and performance of this ability is a complex interplay between model architecture, pre-training data characteristics, and the balance of different architectural components. These insights have important implications for future model development, underscoring the need to move beyond a one-size-fits-all approach and to carefully consider the design of pre-training data and model architectures to optimize for ICL and other desirable capabilities.

\subsection{Evaluating In-Context Learning in Specialized Domains}

The evaluation of in-context learning models in specialized domains, such as scientific computing, medical diagnosis, and legal reasoning, presents unique challenges and considerations that go beyond the general evaluation frameworks discussed in the previous sections. These specialized domains often require domain-specific datasets and evaluation metrics to accurately assess the performance and capabilities of in-context learning models.

In the domain of scientific computing, in-context learning models have shown promise in tasks such as numerical optimization, differential equation solving, and physics-based simulations \cite{ref228}. However, evaluating the performance of these models in such specialized tasks requires careful consideration of the domain-specific requirements and constraints. For example, the accuracy and sample efficiency of in-context learning models in solving complex scientific problems may be more relevant than their performance on standard natural language processing benchmarks. Additionally, the ability of these models to handle numerical inputs and outputs, as well as their capability to capture and leverage the underlying physical and mathematical principles, should be evaluated \cite{ref228}.

Similarly, in the medical domain, the evaluation of in-context learning models for tasks such as diagnosis, treatment recommendations, and patient education should go beyond traditional natural language understanding benchmarks \cite{ref301}. These models should be assessed on their ability to generate medically accurate and contextually relevant responses, their robustness to handle diverse clinical scenarios, and their potential to improve patient outcomes \cite{ref302}. Specialized datasets and evaluation metrics that capture the nuances of medical decision-making and patient-provider interactions are essential for a comprehensive assessment of in-context learning models in the healthcare domain.

In the legal domain, the evaluation of in-context learning models for tasks such as legal reasoning, case law analysis, and legal document generation should consider the unique challenges of this specialized field \cite{ref303}. These models should be assessed on their ability to understand and reason about legal concepts, interpret case law and statutes, and generate legally sound and contextually appropriate responses. Domain-specific datasets that capture the complexities of legal language, reasoning, and decision-making are crucial for evaluating the performance of in-context learning models in the legal domain.

Addressing the unique challenges of evaluating in-context learning models in specialized domains requires the development of domain-specific datasets, evaluation metrics, and benchmarks. These specialized evaluation frameworks should capture the unique characteristics, constraints, and requirements of each domain, while also allowing for meaningful comparisons across different domains and models.

For example, in the medical domain, the evaluation of in-context learning models could include metrics that assess the medical accuracy, safety, and ethical considerations of the generated responses \cite{ref304}. Similarly, in the legal domain, the evaluation could focus on the models' ability to interpret and reason about legal precedents, statutes, and case law, as well as their compliance with legal principles and ethical standards.

Furthermore, the evaluation of in-context learning models in specialized domains should not be limited to just the models' performance on specific tasks. It should also consider the models' robustness, adaptability, and generalizability across a wide range of scenarios within the domain \cite{ref49}. This comprehensive evaluation approach can help ensure that in-context learning models are truly capable of providing reliable and trustworthy support in specialized, high-stakes domains.

In conclusion, the evaluation of in-context learning models in specialized domains, such as scientific computing, medical diagnosis, and legal reasoning, requires a tailored approach that takes into account the unique characteristics, requirements, and constraints of each domain. The development of domain-specific datasets, evaluation metrics, and benchmarks is crucial for assessing the performance and capabilities of these models in a meaningful and comprehensive way. By addressing the unique challenges of specialized domains, the research community can ensure that in-context learning models are effectively deployed and integrated into real-world applications, ultimately enhancing their impact and trust in these critical domains.

\subsection{Benchmarking In-Context Learning with Synthetic Tasks}

The use of synthetic tasks and datasets has emerged as a promising approach for benchmarking in-context learning (ICL) models. Synthetic tasks offer several advantages compared to real-world datasets, primarily the ability to control for confounding factors and systematically probe the capabilities of ICL models \cite{ref305}.

One of the key advantages of using synthetic tasks is the ability to create a controlled environment where the underlying data-generating process is known. This allows researchers to isolate specific aspects of in-context learning, such as the model's ability to perform logical reasoning, arithmetic operations, or understand spatial relationships \cite{ref305}. By manipulating the complexity and structure of the synthetic tasks, researchers can better understand the factors that contribute to a model's in-context learning performance and identify the model's strengths and weaknesses.

Moreover, synthetic tasks can be designed to be scalable, allowing researchers to systematically increase the difficulty or complexity of the tasks to push the limits of a model's capabilities. This is particularly important in the context of in-context learning, where models are expected to perform well on a wide range of tasks with varying levels of complexity \cite{ref216}.

Another advantage of using synthetic tasks is the ability to create datasets that are free from the biases and idiosyncrasies often present in real-world datasets. Real-world datasets can be influenced by a variety of factors, such as the way the data was collected, the demographics of the participants, or the specific domains and contexts represented in the data \cite{ref189}. By using synthetic tasks, researchers can ensure that the evaluation is not skewed by these types of biases, allowing for a more accurate assessment of the model's in-context learning abilities.

However, the use of synthetic tasks for benchmarking in-context learning is not without its limitations. One of the main concerns is the potential for models to overfit to the specific structure and patterns present in the synthetic tasks, leading to inflated performance that may not generalize to real-world scenarios \cite{ref306}. This is particularly true if the synthetic tasks are not designed to capture the complexity and diversity of real-world problems.

To address this concern, researchers have proposed the use of more diverse and realistic synthetic datasets that better mimic the characteristics of real-world data \cite{ref307}. By incorporating elements of real-world data, such as multimodal information or temporal dependencies, these synthetic datasets can provide a more comprehensive and challenging benchmark for in-context learning models.

Another potential limitation of synthetic tasks is the difficulty in ensuring that the tasks accurately capture the nuances and subtleties of real-world problems. While synthetic tasks can be designed to test specific capabilities, they may not fully capture the contextual and pragmatic aspects of natural language understanding that are crucial for in-context learning \cite{ref278}. This is an ongoing challenge in the field, and researchers are exploring ways to bridge the gap between synthetic and real-world tasks, such as by incorporating more realistic language modeling objectives or using meta-learning techniques to adapt the synthetic tasks to better match the target domains \cite{ref308}.

Despite these limitations, the use of synthetic tasks for benchmarking in-context learning has proven to be a valuable tool in the field. By providing a controlled and scalable environment, synthetic tasks have enabled researchers to gain deeper insights into the factors that drive in-context learning performance and to develop more robust and generalizable models \cite{ref183}.

Moreover, the lessons learned from synthetic task benchmarking can inform the design and evaluation of real-world in-context learning applications. By identifying the strengths and weaknesses of models on synthetic tasks, researchers can better understand the limitations of current approaches and develop strategies to address them, leading to more effective and reliable in-context learning systems \cite{ref309}.

As the field of in-context learning continues to evolve, the use of synthetic tasks for benchmarking will likely remain an important tool in the researcher's toolkit. By combining the insights gained from synthetic task benchmarking with real-world evaluations, the research community can work towards developing more robust, versatile, and trustworthy in-context learning models that can effectively tackle a wide range of practical challenges.

\subsection{Towards More Comprehensive and Robust Evaluation}

As the field of in-context learning continues to advance, the need for more comprehensive and robust evaluation frameworks becomes increasingly crucial. The existing evaluation approaches, while valuable, often fall short in capturing the multifaceted nature of in-context learning and the potential challenges that may arise in real-world deployment scenarios.

One key direction for developing more robust evaluation frameworks is the incorporation of multidimensional performance assessment \cite{ref310,ref46}. Current evaluation metrics, such as few-shot accuracy or average precision, provide a limited view of in-context learning capabilities. A more holistic approach would involve assessing in-context learning models across a broader spectrum of performance dimensions, including but not limited to task-specific metrics, interpretability, sample efficiency, and robustness to various distribution shifts.

Evaluating model robustness to distributional shifts is particularly crucial, as in-context learning models are often deployed in dynamic and unpredictable environments \cite{ref311,ref312}. Existing benchmarks and datasets may not capture the full range of potential distribution shifts that these models may encounter in real-world applications. Developing more comprehensive evaluation protocols that systematically test in-context learning models' performance under a diverse set of distribution shifts, including but not limited to domain shifts, label shifts, and spurious correlation shifts, would provide a more realistic assessment of their capabilities and limitations.

Furthermore, the integration of human-in-the-loop evaluation \cite{ref313,ref314} can significantly enhance the robustness and reliability of in-context learning evaluation frameworks. By incorporating human feedback, perspectives, and judgments into the evaluation process, researchers can gain valuable insights into how in-context learning models perform in tasks that require human-level understanding, nuance, and contextual awareness. This can help identify areas where in-context learning models excel or fall short, and guide the development of more user-centric and trustworthy systems.

One approach to incorporating human-in-the-loop evaluation could involve the development of interactive evaluation platforms, where users can engage with in-context learning models, provide feedback, and assess the models' performance in real-time \cite{ref315}. These platforms could also facilitate the collection of diverse user preferences and perspectives, which can then be used to refine the evaluation criteria and identify potential biases or limitations in the models.

Another promising direction is the exploration of synthetic task generation for comprehensive in-context learning evaluation \cite{ref24}. By creating a diverse set of synthetic tasks with varying levels of complexity, contextual information, and domain-specific knowledge requirements, researchers can assess the generalization capabilities of in-context learning models and identify their strengths and weaknesses across a broader range of scenarios.

Additionally, the development of standardized evaluation protocols and benchmarks \cite{ref316,ref24} can greatly enhance the comparability and reproducibility of in-context learning research. These protocols should strive to incorporate the multidimensional performance assessment, distributional robustness evaluation, and human-in-the-loop feedback mentioned earlier, providing a comprehensive and rigorous framework for assessing the capabilities of in-context learning models.

Finally, the exploration of novel evaluation metrics and analysis techniques \cite{ref46} can offer deeper insights into the mechanisms underlying in-context learning. By developing new approaches to quantify the models' ability to effectively leverage contextual information, understand task-relevant concepts, and generalize to unseen scenarios, researchers can gain a more nuanced understanding of the strengths and limitations of in-context learning, ultimately guiding the development of more robust and reliable systems.

In summary, the pursuit of more comprehensive and robust evaluation frameworks for in-context learning is a multifaceted challenge that requires the integration of various approaches, including multidimensional performance assessment, distributional robustness evaluation, human-in-the-loop feedback, synthetic task generation, standardized benchmarking, and novel evaluation metrics. By addressing these gaps, the research community can contribute to the development of in-context learning models that are more reliable, trustworthy, and adaptable to real-world deployments.

\section{Limitations, Challenges, and Ethical Considerations}

\subsection{Sensitivity to Prompt Design}

In-context learning (ICL) has emerged as a powerful capability of large language models (LLMs), allowing them to adapt to new tasks by conditioning on a prompt that includes a few task examples \cite{ref48}. However, the performance of in-context learning has been shown to be highly sensitive to the design of the prompt, including the choice of examples, ordering, and formatting \cite{ref317}.

The choice of examples included in the prompt is a key factor contributing to the sensitivity of in-context learning. Studies have shown that the performance of in-context learning can vary significantly depending on the specific examples used, even when the overall set of examples is held constant \cite{ref318}. This is because the examples included in the prompt can have a significant impact on the model's understanding of the task, as well as its ability to generalize to new inputs \cite{ref291}.

The ordering of the examples within the prompt has also been found to be a critical factor in in-context learning. \cite{ref51} demonstrates that the order of the examples can have a significant impact on the model's performance, with certain orderings leading to much better results than others. This is likely due to the fact that the model's attention and reasoning processes can be heavily influenced by the order in which the examples are presented \cite{ref292}.

In addition to the choice and ordering of examples, the formatting of the prompt itself can also have a significant impact on in-context learning performance \cite{ref48}. Subtle changes in the way the prompt is structured, such as the use of different formatting or the inclusion of additional instructions or context, can lead to large variations in the model's behavior \cite{ref317}.

The sensitivity of in-context learning to prompt design poses a significant challenge for the reliability and generalization of these models. If the performance of a model is highly dependent on the specific prompt used, it becomes difficult to trust the model's outputs or to deploy it in real-world applications \cite{ref319}. Additionally, the need to carefully engineer prompts to achieve good performance can be time-consuming and resource-intensive, limiting the scalability and practicality of in-context learning \cite{ref48}.

\subsection{Susceptibility to Biases}

In-context learning (ICL) has shown remarkable performance in a wide range of tasks, from natural language processing to computer vision. However, this paradigm is not immune to the susceptibility of various biases that can severely undermine the fairness and trustworthiness of the resulting models. The primary sources of bias in ICL can be broadly categorized into data biases, social biases, and algorithmic biases.

Data biases are inherent in the training datasets used to pre-train the large language models (LLMs) that form the backbone of many ICL systems. These datasets often reflect societal inequities, historical prejudices, and sampling biases, which can then be propagated and amplified through the ICL process \cite{ref53}. For example, \cite{ref320,ref321} demonstrates how biased training data can lead to unfair outcomes in machine learning models, particularly in high-stakes decision-making domains.

Social biases, such as those related to gender, race, age, or socioeconomic status, can also creep into ICL systems through the prompts and examples used during the in-context learning stage. The choice of prompts, the framing of the task, and the selection of in-context examples can all inadvertently reinforce harmful stereotypes and prejudices \cite{ref322}. For instance, \cite{ref323} finds that text-only models exhibit less bias compared to multimodal models, highlighting the potential for social biases to be amplified when incorporating additional modalities.

Algorithmic biases can also arise from the very nature of the in-context learning process. The sensitivity of ICL to prompt design and the potential for feedback loops between the model and the user can lead to the amplification of biases over successive iterations \cite{ref324}. \cite{ref325} further demonstrates how the interaction between the predictive model and the data can result in the model-induced exacerbation of biases, even when the initial data is unbiased.

The implications of these biases for the fairness and trustworthiness of ICL systems are significant. Biased models can make decisions that disproportionately harm specific demographic groups, leading to unfair outcomes and a lack of trust in the system \cite{ref326}. This is particularly concerning in high-stakes applications, such as healthcare, finance, or criminal justice, where the consequences of biased decisions can be severe.

Moreover, the opacity of the in-context learning process and the potential for hidden biases can make it challenging to detect and mitigate these issues. \cite{ref46} proposes a framework for evaluating the hardness and contextual information in ICL, but more research is needed to develop comprehensive methods for auditing and addressing bias in these systems.

To address the susceptibility of ICL to biases, a multifaceted approach is required. First, researchers and practitioners must be vigilant about the potential sources of bias in the pre-trained LLMs, the training datasets used, and the prompts and examples employed during the in-context learning phase \cite{ref53,ref320,ref321}. Techniques such as debiasing the training data, carefully curating prompts and examples, and incorporating fairness constraints into the ICL process can help mitigate these biases \cite{ref327,ref328}.

Additionally, the development of robust evaluation frameworks and benchmarks for assessing the fairness and bias in ICL systems is crucial \cite{ref329,ref330}. These tools can help researchers and practitioners identify and address biases, as well as provide a means for transparent and accountable development of ICL-based applications.

Finally, the ethical and societal implications of biases in ICL must be carefully considered. \cite{ref58} highlights the need to incorporate ethical and legal perspectives when developing ICL systems for high-stakes applications. By acknowledging the susceptibility of ICL to biases and proactively addressing these challenges, the research community can work towards building more trustworthy and equitable in-context learning systems that serve the needs of all members of society.

\subsection{Need for Scalable and Efficient Algorithms}

The computational and memory challenges associated with in-context learning (ICL) have become increasingly apparent as the field has progressed. As the number of in-context examples and the complexity of the tasks increase, the resource requirements for effective ICL can become prohibitive, posing a significant obstacle to its practical deployment in real-world applications.

One of the primary challenges lies in the scalability of ICL algorithms. As models are required to process longer prompt sequences, the computational complexity grows rapidly, often following a quadratic or even exponential relationship with the input length \cite{ref331}. This can lead to significant performance bottlenecks, particularly for resource-constrained devices or applications that require real-time responses.

Moreover, the memory requirements of ICL can also pose a significant challenge. The need to store and process a large number of in-context examples, along with the model's internal representations, can quickly consume available memory resources, especially for models with high parameter counts \cite{ref332}. This issue is exacerbated when deploying ICL-based systems in environments with limited memory, such as edge devices or mobile platforms.

The rapid scaling of in-context examples can also lead to increased computational costs during the inference process. As the number of examples grows, the time required to process the entire prompt and generate the desired output can become prohibitively long, making it impractical for many real-world applications that require timely responses \cite{ref331}.

To address these challenges, there is a pressing need for the development of scalable and efficient ICL algorithms. Researchers have explored various strategies to mitigate the computational and memory requirements of ICL, such as the use of structured prompting \cite{ref331}, which can reduce the quadratic complexity of attention-based models, and the integration of processing-in-memory (PIM) architectures \cite{ref332}, which can alleviate the data movement bottleneck inherent in traditional processor-centric systems.

Another promising direction is the exploration of more efficient neural network architectures and optimization techniques specifically tailored for ICL. For example, the development of novel attention mechanisms \cite{ref333} or the incorporation of meta-learning principles \cite{ref334} could lead to more scalable and efficient ICL algorithms.

Furthermore, the integration of ICL with other learning paradigms, such as transfer learning and multi-task learning, could also contribute to the development of more efficient and generalizable ICL systems. By leveraging knowledge gained from related tasks or domains, ICL models may become more data-efficient and better equipped to handle a wider range of applications \cite{ref335}.

In addition, the exploration of compression techniques, such as model quantization and weight pruning, could also help to reduce the memory footprint of ICL models, enabling their deployment on resource-constrained platforms \cite{ref336}.

Overall, the need for scalable and efficient ICL algorithms is a pressing challenge that must be addressed to unlock the full potential of in-context learning in real-world applications. The development of innovative approaches that can handle increasing task complexity and scale while maintaining high performance and low resource requirements will be crucial for the widespread adoption and practical deployment of ICL-based systems.

\subsection{Ethical Implications}

The widespread deployment of in-context learning models, particularly in high-stakes domains, raises significant ethical concerns that require urgent attention. As these models become increasingly prevalent, the potential for unintended consequences and the need for accountability and transparency in their development and deployment have become particularly pressing issues.

One of the primary ethical concerns with in-context learning is the potential for amplifying and perpetuating societal biases \cite{ref337}. In-context learning models are often trained on large datasets that reflect historical and systemic biases, which can then be encoded and propagated through the models' outputs. This can lead to unfair and discriminatory decisions, with disproportionate impacts on marginalized communities \cite{ref338}. The sensitivity of in-context learning models to prompt design further exacerbates this issue, as biases can be introduced through the selection and framing of the examples provided to the model \cite{ref51}.

Additionally, the lack of transparency and interpretability in many in-context learning models poses a significant challenge to ensuring accountability. These models often operate as ``black boxes,'' making it difficult to understand the reasoning behind their decisions and the factors that influence their outputs \cite{ref339}. This lack of transparency can hinder the ability to identify and address biases, as well as to hold developers and deployers accountable for the consequences of their systems \cite{ref340}.

Furthermore, the potential for unintended consequences in the deployment of in-context learning models is a serious ethical concern. These models may be applied in complex, high-stakes domains, such as healthcare, criminal justice, and education, where their decisions can have profound impacts on people's lives \cite{ref341}. The unpredictable and context-dependent nature of in-context learning can lead to unforeseen outcomes that may contradict societal values and human rights \cite{ref56}.

The ethical implications of in-context learning also extend to issues of privacy and data rights. These models often rely on large, diverse datasets that may contain sensitive personal information. The potential for data breaches, unauthorized use, or misuse of this information raises concerns about individual privacy and the right to self-determination \cite{ref342}.

To address these ethical concerns, it is essential to develop in-context learning systems that adhere to robust ethical principles and consider the societal impact of their decisions \cite{ref343}. This may involve incorporating ethical considerations into the entire lifecycle of in-context learning, from data collection and model design to deployment and monitoring \cite{ref344}. Specific strategies may include implementing transparent and interpretable model architectures, establishing clear accountability mechanisms, and engaging with diverse stakeholders to understand and mitigate the potential harms \cite{ref345}.

Additionally, the development of comprehensive evaluation frameworks that assess the ethical implications of in-context learning models, beyond just their technical performance, is crucial \cite{ref346}. Furthermore, the integration of in-context learning with other ethical AI frameworks, such as fairness, transparency, and privacy-preserving techniques, can help ensure that these models are developed and deployed in a responsible and socially conscious manner \cite{ref347}.

Ultimately, the ethical implications of in-context learning are complex and multifaceted, requiring a holistic approach that balances the potential benefits of these models with the need to protect individual rights, promote social justice, and ensure the responsible development and deployment of AI systems \cite{ref348}.

\subsection{Mitigation Strategies}

The limitations and ethical challenges associated with in-context learning (ICL) are multifaceted and demand a comprehensive set of mitigation strategies to address them effectively. Recognizing the sensitivity of ICL to prompt design \cite{ref51}, the susceptibility to biases \cite{ref349}, and the need for scalable and efficient algorithms \cite{ref350}, researchers have proposed several promising approaches to overcome these obstacles.

One key mitigation strategy involves the development of novel prompt engineering techniques that can enhance the robustness and reliability of in-context learning. Researchers have explored methods to systematically analyze the sensitivity of different prompts and their impact on model performance \cite{ref51}. By identifying the crucial factors that contribute to prompt sensitivity, such as the choice of examples, their ordering, and formatting, researchers can devise techniques to design prompts that are less susceptible to subtle variations and more likely to elicit the desired model behavior. Furthermore, the incorporation of human feedback and expert knowledge into the prompt engineering process can help ensure that the prompts align with societal values and expectations \cite{ref56}.

Another promising approach to mitigate the limitations of in-context learning involves the development of bias-aware in-context learning algorithms. Researchers have proposed techniques that can actively identify and address various types of biases, such as data biases, social biases, and algorithmic biases, that can get amplified through the in-context learning process \cite{ref349}. By integrating fairness-aware mechanisms into the in-context learning algorithms, the models can be trained to be more cognizant of these biases and take proactive steps to mitigate their adverse impacts. This can involve the incorporation of debiasing techniques, the use of adversarial training, or the development of more sophisticated model architectures that are inherently less susceptible to biases.

In addition to technical solutions, the mitigation of in-context learning limitations and ethical challenges also requires the incorporation of human oversight and governance frameworks. Researchers have advocated for the involvement of diverse stakeholders, including domain experts, policymakers, and affected communities, in the development and deployment of in-context learning systems \cite{ref351}. By engaging these stakeholders, the in-context learning systems can be designed and evaluated with a more holistic understanding of the societal context and the potential impacts on different groups. Furthermore, the establishment of robust governance frameworks, such as model auditing protocols, transparency and accountability measures, and ethical review processes, can help ensure that the in-context learning systems adhere to relevant ethical principles and societal values \cite{ref352}.

The mitigation of in-context learning limitations and ethical challenges also necessitates a shift in the research and development culture towards more responsible and inclusive practices. Researchers have highlighted the importance of diversifying the teams involved in the development of in-context learning systems, as the lack of diverse perspectives can lead to the perpetuation of existing biases and the overlooking of critical societal considerations \cite{ref353}. By fostering a culture of collaboration and interdisciplinary engagement, the in-context learning community can better address the complex sociotechnical issues that arise from the deployment of these systems \cite{ref354}.

In conclusion, the mitigation of in-context learning limitations and ethical challenges requires a multifaceted approach that combines technical innovations, human-centered design, and responsible governance frameworks. By developing novel prompt engineering techniques, bias-aware in-context learning algorithms, and incorporating human oversight and diverse stakeholder engagement, the research community can work towards the development of in-context learning systems that are more reliable, fair, and aligned with societal values. This collective effort is crucial to realizing the full potential of in-context learning while ensuring its responsible and ethical deployment in various domains.

\section{Future Directions and Conclusion}

\subsection{Integration with Other Learning Paradigms}

The integration of in-context learning with other prominent learning paradigms, such as transfer learning, meta-learning, and multi-task learning, holds immense potential for enhancing the adaptability, sample efficiency, and generalization capabilities of intelligent systems. By leveraging the synergies between these approaches, researchers aim to unlock new frontiers in the development of versatile and robust learning algorithms.

Transfer learning has long been recognized as a powerful technique for leveraging knowledge acquired from one domain or task to improve performance on a related, but distinct, target task \cite{ref68}. In the context of in-context learning, transfer learning can play a crucial role in enabling the rapid adaptation of language models and other intelligent agents to new tasks and domains. By pretraining on a diverse set of tasks and then fine-tuning the model's parameters on a few in-context examples, the model can rapidly adapt its in-context learning capabilities to the new task at hand \cite{ref334}.

Furthermore, the integration of meta-learning and in-context learning has emerged as a particularly promising direction. Meta-learning, or learning-to-learn, focuses on developing models that can quickly adapt to new tasks by leveraging their previous experiences \cite{ref355,ref356}. In the context of in-context learning, meta-learning can enable language models and other intelligent agents to more efficiently learn new tasks from a few examples by extracting and leveraging relevant meta-knowledge from their previous in-context learning experiences \cite{ref357}.

The synergies between meta-learning and in-context learning are particularly compelling, as both approaches aim to enhance the adaptability and generalization capabilities of learning systems. By combining the meta-learning paradigm with the in-context learning framework, researchers can develop models that can quickly adapt to new tasks while still maintaining the benefits of in-context learning, such as the ability to leverage contextual information and demonstrate few-shot learning abilities \cite{ref2}.

In addition to the integration of transfer learning and meta-learning, the combination of in-context learning with multi-task learning also holds significant promise. Multi-task learning, which involves the simultaneous learning of multiple related tasks, can help language models and other intelligent agents develop more robust and generalizable representations \cite{ref358,ref359}. By exposing in-context learning models to a diverse set of tasks during training, they can learn to better leverage the shared structure and relationships between tasks, leading to enhanced adaptability and generalization \cite{ref360}.

Moreover, the integration of in-context learning with other emerging learning paradigms, such as self-supervised learning and continual learning, can further expand the capabilities of intelligent systems. Self-supervised learning, which involves learning representations from unlabeled data, can help in-context learning models develop more robust and versatile feature representations that can be effectively leveraged in new tasks \cite{ref361}. Similarly, the integration of in-context learning with continual learning, which focuses on the ability to learn new skills without catastrophically forgetting previously acquired knowledge, can lead to highly adaptable and resilient learning agents \cite{ref362}.

Overall, the synergistic integration of in-context learning with other prominent learning paradigms, such as transfer learning, meta-learning, and multi-task learning, holds immense promise for advancing the field of artificial intelligence. By leveraging the strengths of these complementary approaches, researchers can develop intelligent systems that can rapidly adapt to new tasks and domains, maintain high performance on a diverse set of challenges, and continuously expand their capabilities over time. As the field of in-context learning continues to evolve, the exploration of these integrative approaches will undoubtedly play a critical role in shaping the future of adaptive and versatile artificial intelligence.

\subsection{Novel Architectural Designs}

One promising direction for improving the in-context learning abilities of language models (LMs) and vision-language models (VLMs) is the development of novel neural network architectures and attention mechanisms. Recent studies have highlighted the potential of architectural innovations to enhance models' capacity to effectively leverage contextual information.

A key focus of this line of research has been on attention mechanisms, which have become a fundamental component of modern transformer-based models. The self-attention mechanism in transformers enables the models to capture dependencies between different parts of the input sequence, which has been crucial for their success in various natural language processing tasks \cite{ref363}. However, as transformer-based LMs and VLMs are scaled up to handle longer input sequences and more complex tasks, the quadratic complexity of the attention mechanism can become a significant bottleneck, limiting their ability to effectively incorporate contextual information.

To address this issue, researchers have proposed novel attention architectures that aim to improve the efficiency and effectiveness of contextual modeling. One approach is to introduce structured attention mechanisms that constrain the attention patterns to better match the inherent structure of the data. For example, \cite{ref364} introduces a structure-regularized attention mechanism that formalizes feature interaction as structural factorization, enabling the model to capture the structural dependencies in the data in an efficient manner. Similarly, \cite{ref365} presents a context outlooker attention mechanism that encodes local syntactic context by considering word proximity and other pair-wise constraints, which can be more effective than standard self-attention for certain natural language processing tasks.

Another direction for architectural innovation is the integration of memory modules within the LM and VLM architectures. \cite{ref366} proposes a memory-augmented transformer model that incorporates trainable memory tokens to selectively store local and global representations of the input sequence, which can improve the model's ability to process contextual information. By introducing a memory bottleneck and a dedicated layer for memory update, the Memory Transformer aims to better capture the global context and reduce the negative interference between different aspects of the input.

In the domain of computer vision, researchers have also explored novel architectural designs to enhance the in-context learning capabilities of VLMs. \cite{ref241} investigates the hypothesis that the in-context learning ability of LMs can be transferred to VLMs. The authors propose a method to meta-train an LM for in-context learning on NLP tasks and then transfer the learned abilities to a VLM by attaching a visual encoder. Their experiments demonstrate that this cross-modal transfer can significantly improve the in-context learning performance of VLMs on vision-and-language tasks, even with a smaller model size.

Moreover, researchers have explored the integration of different modalities, such as vision and language, within a unified architectural framework to better leverage contextual information. \cite{ref367} presents an architecture that combines state-of-the-art transformer and ResNeXt modules with a novel attentive multimodal module to produce a combined model trained on various vision and language tasks. The authors show that this approach can obtain good results on a wide range of vision and language tasks, demonstrating the potential of using a single, versatile model that can effectively handle multimodal contextual information.

Beyond attention mechanisms and memory modules, researchers have also investigated other architectural innovations to enhance in-context learning. \cite{ref368} introduces a novel architecture called Zebra that efficiently manages the quadratic time and memory complexity of standard transformer attention by employing a grouped local-global attention mechanism. This approach aims to balance the trade-off between local and global attention, allowing the model to process longer input sequences without compromising its ability to effectively leverage contextual information.

Overall, the development of novel neural network architectures and attention mechanisms is a promising avenue for improving the in-context learning abilities of LMs and VLMs. By designing architectures that can more effectively capture and leverage contextual information, researchers hope to overcome the limitations of current models and enable more versatile and adaptive intelligent systems that can rapidly adapt to new tasks and situations \cite{ref369,ref370,ref371}. As the field of in-context learning continues to evolve, we can expect to see more exciting advancements in this area, with potential applications ranging from natural language processing to multimodal reasoning and beyond.

\subsection{Cross-Modal and Cross-Lingual Transfer}

The potential for transferring in-context learning abilities across different modalities and languages holds significant promise in enhancing the versatility and accessibility of these powerful models. Recent advancements in the field have demonstrated the ability of large language models (LLMs) to excel in various tasks through in-context learning \cite{ref372,ref121}.

One key area of exploration is the transfer of in-context learning from language to vision. The work by \cite{ref241} has shown that the in-context learning ability developed in language models can be effectively transferred to vision-language models, leading to significant improvements in performance on various vision-language tasks, even with a smaller model size.

Building upon this, researchers have also explored the potential of incorporating multi-modal context to further boost in-context learning performance. The \cite{ref373} framework demonstrates the importance of modeling and leveraging various types of contextual information, including multi-modal, visual-semantic, and temporal context, to enhance the in-context learning abilities of vision-language models. Similarly, the work by \cite{ref132} provides valuable insights into the significance of both visual and language information in the in-context learning process.

The challenge of enabling cross-lingual in-context learning is another critical area of research. With the growing importance of supporting diverse languages in various applications, the ability to transfer in-context learning capabilities across languages is crucial. The work by \cite{ref77} explores the use of in-context learning prompts that leverage both source and target languages, demonstrating significant improvements in cross-lingual transfer performance on various multilingual benchmarks.

Extending the cross-modal and cross-lingual capabilities of in-context learning models can unlock a wide range of opportunities. For instance, the ability to seamlessly transfer in-context learning skills from language to vision can enable more efficient and adaptable vision-language systems, capable of rapidly adjusting to new tasks and contexts through in-context demonstrations \cite{ref241}. Similarly, the cross-lingual transfer of in-context learning can greatly enhance the accessibility and inclusivity of these models, making them more widely applicable across different languages and cultural contexts \cite{ref77}.

However, the realization of this potential is not without its challenges. Researchers need to address issues such as the sensitivity of in-context learning to prompt design, the susceptibility of these models to biases, and the need for scalable and efficient algorithms to enable practical deployment \cite{ref1}. Strategies like advanced prompt engineering, bias-aware in-context learning algorithms, and the incorporation of human oversight and governance frameworks will be crucial in mitigating these limitations \cite{ref1}.

Overall, the exploration of cross-modal and cross-lingual in-context learning represents a promising frontier in the field of intelligent systems. By leveraging the adaptability and versatility of in-context learning, researchers can work towards developing models that can seamlessly bridge the gap between different modalities and languages, ultimately enhancing the accessibility, inclusivity, and utility of these powerful technologies.

\subsection{Concluding Perspectives}

In-context learning has emerged as a transformative capability in modern artificial intelligence, reshaping the landscape of how we approach the development of intelligent systems. As we look to the future, the importance of in-context learning and its profound implications for the future of AI cannot be overstated.

At the heart of in-context learning lies the remarkable ability of large language models (LLMs) to adapt and perform new tasks by simply observing a few relevant examples, without the need for extensive retraining or fine-tuning \cite{ref310,ref350}. This remarkable feat represents a significant departure from the traditional machine learning paradigm, which has often struggled with the challenges of poor generalization, data-hungry training, and the inability to effectively leverage contextual information \cite{ref374}. In-context learning has the potential to address these limitations, paving the way for more adaptable, efficient, and human-centric AI systems.

One of the key contributions of in-context learning is its capacity to bridge the gap between the narrow specialization of current AI applications and the human-like flexibility and adaptability that we aspire to in artificial intelligence \cite{ref375,ref376}. By enabling AI systems to rapidly acquire new skills and knowledge from limited examples, in-context learning can unlock a new era of AI that is better suited to navigate the complex, ever-changing, and often ambiguous real-world environments \cite{ref377}.

Moreover, the insights gleaned from the study of in-context learning have the potential to inform and catalyze advancements in other areas of AI research. For instance, the mechanisms underlying in-context learning can provide valuable clues about the nature of human intelligence and cognition, potentially informing the development of more biologically plausible and psychologically realistic models of intelligence \cite{ref378}. Similarly, the exploration of in-context learning can inspire new architectures, training paradigms, and learning algorithms that can enhance the performance and capabilities of AI systems in diverse domains, from natural language processing to computer vision and beyond \cite{ref379,ref380}.

As we look to the future, the continued advancement and integration of in-context learning into the broader field of AI is poised to play a pivotal role in shaping the next generation of intelligent systems. By harnessing the power of in-context learning, researchers and practitioners can work towards the development of AI that is more adaptable, efficient, and deeply attuned to the complexities of the human experience \cite{ref381,ref382}. This, in turn, can unlock new frontiers in human-AI collaboration, where the unique strengths of both artificial and human intelligence are seamlessly combined to tackle some of the most complex challenges facing our world \cite{ref383,ref377}.


\begin{thebibliography}{383}
\addcontentsline{toc}{section}{References}
\bibitem{ref1} The Learnability of In-Context Learning. arXiv:2303.07895v1, 2023. \url{https://arxiv.org/abs/2303.07895v1}
\bibitem{ref2} Meta-in-context learning in large language models. arXiv:2305.12907v1, 2023. \url{https://arxiv.org/abs/2305.12907v1}
\bibitem{ref3} An Explanation of In-context Learning as Implicit Bayesian Inference. arXiv:2111.02080v6, 2021. \url{https://arxiv.org/abs/2111.02080v6}
\bibitem{ref4} Are Emergent Abilities in Large Language Models just In-Context Learning. arXiv:2309.01809v1, 2023. \url{https://arxiv.org/abs/2309.01809v1}
\bibitem{ref5} Data Distributional Properties Drive Emergent In-Context Learning in Transformers. arXiv:2205.05055v6, 2022. \url{https://arxiv.org/abs/2205.05055v6}
\bibitem{ref6} The Mystery of In-Context Learning A Comprehensive Survey on Interpretation and Analysis. arXiv:2311.00237v2, 2023. \url{https://arxiv.org/abs/2311.00237v2}
\bibitem{ref7} Explaining Emergent In-Context Learning as Kernel Regression. arXiv:2305.12766v2, 2023. \url{https://arxiv.org/abs/2305.12766v2}
\bibitem{ref8} Supervised Knowledge Makes Large Language Models Better In-context Learners. arXiv:2312.15918v2, 2023. \url{https://arxiv.org/abs/2312.15918v2}
\bibitem{ref9} Instruct Me More! Random Prompting for Visual In-Context Learning. arXiv:2311.03648v1, 2023. \url{https://arxiv.org/abs/2311.03648v1}
\bibitem{ref10} Fine-tune Language Models to Approximate Unbiased In-context Learning. arXiv:2310.03331v1, 2023. \url{https://arxiv.org/abs/2310.03331v1}
\bibitem{ref11} Towards More Unified In-context Visual Understanding. arXiv:2312.02520v2, 2023. \url{https://arxiv.org/abs/2312.02520v2}
\bibitem{ref12} PRODIGY Enabling In-context Learning Over Graphs. arXiv:2305.12600v1, 2023. \url{https://arxiv.org/abs/2305.12600v1}
\bibitem{ref13} Rethinking the Role of Scale for In-Context Learning An Interpretability-based Case Study at 66 Billion Scale. arXiv:2212.09095v2, 2022. \url{https://arxiv.org/abs/2212.09095v2}
\bibitem{ref14} Learning To Retrieve Prompts for In-Context Learning. arXiv:2112.08633v2, 2021. \url{https://arxiv.org/abs/2112.08633v2}
\bibitem{ref15} Revisiting Demonstration Selection Strategies in In-Context Learning. arXiv:2401.12087v1, 2024. \url{https://arxiv.org/abs/2401.12087v1}
\bibitem{ref16} The Strong Pull of Prior Knowledge in Large Language Models and Its Impact on Emotion Recognition. arXiv:2403.17125v1, 2024. \url{https://arxiv.org/abs/2403.17125v1}
\bibitem{ref17} Improving Input-label Mapping with Demonstration Replay for In-context Learning. arXiv:2310.19572v1, 2023. \url{https://arxiv.org/abs/2310.19572v1}
\bibitem{ref18} Unified Demonstration Retriever for In-Context Learning. arXiv:2305.04320v2, 2023. \url{https://arxiv.org/abs/2305.04320v2}
\bibitem{ref19} Iterative Forward Tuning Boosts In-context Learning in Language Models. arXiv:2305.13016v2, 2023. \url{https://arxiv.org/abs/2305.13016v2}
\bibitem{ref20} An Exploration of In-Context Learning for Speech Language Model. arXiv:2310.12477v1, 2023. \url{https://arxiv.org/abs/2310.12477v1}
\bibitem{ref21} Interactive Natural Language Processing. arXiv:2305.13246v1, 2023. \url{https://arxiv.org/abs/2305.13246v1}
\bibitem{ref22} In-context Learning Distillation Transferring Few-shot Learning Ability of Pre-trained Language Models. arXiv:2212.10670v1, 2022. \url{https://arxiv.org/abs/2212.10670v1}
\bibitem{ref23} Learning an Explicit Hyperparameter Prediction Function Conditioned on Tasks. arXiv:2107.02378v3, 2021. \url{https://arxiv.org/abs/2107.02378v3}
\bibitem{ref24} Exploring Effective Factors for Improving Visual In-Context Learning. arXiv:2304.04748v1, 2023. \url{https://arxiv.org/abs/2304.04748v1}
\bibitem{ref25} Meta-Learning without Memorization. arXiv:1912.03820v3, 2019. \url{https://arxiv.org/abs/1912.03820v3}
\bibitem{ref26} Meta-Learning with Context-Agnostic Initialisations. arXiv:2007.14658v2, 2020. \url{https://arxiv.org/abs/2007.14658v2}
\bibitem{ref27} Contextualizing Enhances Gradient Based Meta Learning. arXiv:2007.10143v1, 2020. \url{https://arxiv.org/abs/2007.10143v1}
\bibitem{ref28} Meta-Learning in Neural Networks A Survey. arXiv:2004.05439v2, 2020. \url{https://arxiv.org/abs/2004.05439v2}
\bibitem{ref29} In-Context Learning for Few-Shot Molecular Property Prediction. arXiv:2310.08863v1, 2023. \url{https://arxiv.org/abs/2310.08863v1}
\bibitem{ref30} Boosting Few-Shot Learning With Adaptive Margin Loss. arXiv:2005.13826v1, 2020. \url{https://arxiv.org/abs/2005.13826v1}
\bibitem{ref31} Few-shot Multimodal Multitask Multilingual Learning. arXiv:2303.12489v1, 2023. \url{https://arxiv.org/abs/2303.12489v1}
\bibitem{ref32} Exploiting the Potential of Seq2Seq Models as Robust Few-Shot Learners. arXiv:2307.14856v1, 2023. \url{https://arxiv.org/abs/2307.14856v1}
\bibitem{ref33} Towards Contextual Learning in Few-shot Object Classification. arXiv:1912.06679v3, 2019. \url{https://arxiv.org/abs/1912.06679v3}
\bibitem{ref34} Wandering Within a World Online Contextualized Few-Shot Learning. arXiv:2007.04546v3, 2020. \url{https://arxiv.org/abs/2007.04546v3}
\bibitem{ref35} Adaptive Cross-Modal Few-Shot Learning. arXiv:1902.07104v3, 2019. \url{https://arxiv.org/abs/1902.07104v3}
\bibitem{ref36} MERLOT Multimodal Neural Script Knowledge Models. arXiv:2106.02636v3, 2021. \url{https://arxiv.org/abs/2106.02636v3}
\bibitem{ref37} Optimization of Prompt Learning via Multi-Knowledge Representation for Vision-Language Models. arXiv:2404.10357v2, 2024. \url{https://arxiv.org/abs/2404.10357v2}
\bibitem{ref38} With a Little Help from my Temporal Context Multimodal Egocentric Action Recognition. arXiv:2111.01024v1, 2021. \url{https://arxiv.org/abs/2111.01024v1}
\bibitem{ref39} MMICT Boosting Multi-Modal Fine-Tuning with In-Context Examples. arXiv:2312.06363v2, 2023. \url{https://arxiv.org/abs/2312.06363v2}
\bibitem{ref40} Dynamic Fusion with Intra- and Inter- Modality Attention Flow for Visual Question Answering. arXiv:1812.05252v4, 2018. \url{https://arxiv.org/abs/1812.05252v4}
\bibitem{ref41} DataPerf Benchmarks for Data-Centric AI Development. arXiv:2207.10062v4, 2022. \url{https://arxiv.org/abs/2207.10062v4}
\bibitem{ref42} GEMv2 Multilingual NLG Benchmarking in a Single Line of Code. arXiv:2206.11249v3, 2022. \url{https://arxiv.org/abs/2206.11249v3}
\bibitem{ref43} InstructEval Systematic Evaluation of Instruction Selection Methods. arXiv:2307.00259v2, 2023. \url{https://arxiv.org/abs/2307.00259v2}
\bibitem{ref44} Personalized Benchmarking with the Ludwig Benchmarking Toolkit. arXiv:2111.04260v1, 2021. \url{https://arxiv.org/abs/2111.04260v1}
\bibitem{ref45} NICE To Optimize In-Context Examples or Not. arXiv:2402.06733v2, 2024. \url{https://arxiv.org/abs/2402.06733v2}
\bibitem{ref46} Measuring Pointwise \$\textbackslash\{\}mathcal\{V\}\$-Usable Information In-Context-ly. arXiv:2310.12300v2, 2023. \url{https://arxiv.org/abs/2310.12300v2}
\bibitem{ref47} What are the best systems New perspectives on NLP Benchmarking. arXiv:2202.03799v4, 2022. \url{https://arxiv.org/abs/2202.03799v4}
\bibitem{ref48} Mind Your Format Towards Consistent Evaluation of In-Context Learning Improvements. arXiv:2401.06766v2, 2024. \url{https://arxiv.org/abs/2401.06766v2}
\bibitem{ref49} Evaluation Gaps in Machine Learning Practice. arXiv:2205.05256v1, 2022. \url{https://arxiv.org/abs/2205.05256v1}
\bibitem{ref50} Stronger Random Baselines for In-Context Learning. arXiv:2404.13020v1, 2024. \url{https://arxiv.org/abs/2404.13020v1}
\bibitem{ref51} How are Prompts Different in Terms of Sensitivity. arXiv:2311.07230v1, 2023. \url{https://arxiv.org/abs/2311.07230v1}
\bibitem{ref52} Should ChatGPT be Biased Challenges and Risks of Bias in Large Language Models. arXiv:2304.03738v3, 2023. \url{https://arxiv.org/abs/2304.03738v3}
\bibitem{ref53} Towards Understanding and Mitigating Social Biases in Language Models. arXiv:2106.13219v1, 2021. \url{https://arxiv.org/abs/2106.13219v1}
\bibitem{ref54} Tiny, always-on and fragile Bias propagation through design choices in on-device machine learning workflows. arXiv:2201.07677v4, 2022. \url{https://arxiv.org/abs/2201.07677v4}
\bibitem{ref55} Base rate neglect in computer science education. arXiv:2209.08312v2, 2022. \url{https://arxiv.org/abs/2209.08312v2}
\bibitem{ref56} Designing for Human Rights in AI. arXiv:2005.04949v2, 2020. \url{https://arxiv.org/abs/2005.04949v2}
\bibitem{ref57} A Broader View on Bias in Automated Decision-Making Reflecting on Epistemology and Dynamics. arXiv:1807.00553v2, 2018. \url{https://arxiv.org/abs/1807.00553v2}
\bibitem{ref58} Fairness and Bias in Robot Learning. arXiv:2207.03444v2, 2022. \url{https://arxiv.org/abs/2207.03444v2}
\bibitem{ref59} Coping with Mistreatment in Fair Algorithms. arXiv:2102.10750v1, 2021. \url{https://arxiv.org/abs/2102.10750v1}
\bibitem{ref60} Developing a Philosophical Framework for Fair Machine Learning Lessons From The Case of Algorithmic Collusion. arXiv:2208.06308v2, 2022. \url{https://arxiv.org/abs/2208.06308v2}
\bibitem{ref61} Queering the ethics of AI. arXiv:2308.13591v1, 2023. \url{https://arxiv.org/abs/2308.13591v1}
\bibitem{ref62} From Algorithmic Black Boxes to Adaptive White Boxes Declarative Decision-Theoretic Ethical Programs as Codes of Ethics. arXiv:1711.06035v1, 2017. \url{https://arxiv.org/abs/1711.06035v1}
\bibitem{ref63} Fairway A Way to Build Fair ML Software. arXiv:2003.10354v6, 2020. \url{https://arxiv.org/abs/2003.10354v6}
\bibitem{ref64} Adaptive Data Debiasing through Bounded Exploration. arXiv:2110.13054v2, 2021. \url{https://arxiv.org/abs/2110.13054v2}
\bibitem{ref65} Responsible and Representative Multimodal Data Acquisition and Analysis On Auditability, Benchmarking, Confidence, Data-Reliance \& Explainability. arXiv:1903.07171v1, 2019. \url{https://arxiv.org/abs/1903.07171v1}
\bibitem{ref66} Ethical behavior in humans and machines -- Evaluating training data quality for beneficial machine learning. arXiv:2008.11463v1, 2020. \url{https://arxiv.org/abs/2008.11463v1}
\bibitem{ref67} An AI-based Learning Companion Promoting Lifelong Learning Opportunities for All. arXiv:2112.01242v1, 2021. \url{https://arxiv.org/abs/2112.01242v1}
\bibitem{ref68} Sharing to learn and learning to share -- Fitting together Meta-Learning, Multi-Task Learning, and Transfer Learning A meta review. arXiv:2111.12146v6, 2021. \url{https://arxiv.org/abs/2111.12146v6}
\bibitem{ref69} Advances and Challenges in Meta-Learning A Technical Review. arXiv:2307.04722v1, 2023. \url{https://arxiv.org/abs/2307.04722v1}
\bibitem{ref70} Fine-Tune Language Models as Multi-Modal Differential Equation Solvers. arXiv:2308.05061v4, 2023. \url{https://arxiv.org/abs/2308.05061v4}
\bibitem{ref71} Perspectives and Prospects on Transformer Architecture for Cross-Modal Tasks with Language and Vision. arXiv:2103.04037v2, 2021. \url{https://arxiv.org/abs/2103.04037v2}
\bibitem{ref72} Switch-BERT Learning to Model Multimodal Interactions by Switching Attention and Input. arXiv:2306.14182v1, 2023. \url{https://arxiv.org/abs/2306.14182v1}
\bibitem{ref73} Multimodality of AI for Education Towards Artificial General Intelligence. arXiv:2312.06037v2, 2023. \url{https://arxiv.org/abs/2312.06037v2}
\bibitem{ref74} Multimodal Intelligence Representation Learning, Information Fusion, and Applications. arXiv:1911.03977v3, 2019. \url{https://arxiv.org/abs/1911.03977v3}
\bibitem{ref75} Multilingual Multimodality A Taxonomical Survey of Datasets, Techniques, Challenges and Opportunities. arXiv:2210.16960v1, 2022. \url{https://arxiv.org/abs/2210.16960v1}
\bibitem{ref76} Enhancing Cross-lingual Transfer via Phonemic Transcription Integration. arXiv:2307.04361v1, 2023. \url{https://arxiv.org/abs/2307.04361v1}
\bibitem{ref77} Boosting Cross-lingual Transferability in Multilingual Models via In-Context Learning. arXiv:2305.15233v1, 2023. \url{https://arxiv.org/abs/2305.15233v1}
\bibitem{ref78} Building Human-like Communicative Intelligence A Grounded Perspective. arXiv:2201.02734v1, 2022. \url{https://arxiv.org/abs/2201.02734v1}
\bibitem{ref79} A Survey of Vision-Language Pre-training from the Lens of Multimodal Machine Translation. arXiv:2306.07198v1, 2023. \url{https://arxiv.org/abs/2306.07198v1}
\bibitem{ref80} Unleashing the potential of prompt engineering in Large Language Models a comprehensive review. arXiv:2310.14735v2, 2023. \url{https://arxiv.org/abs/2310.14735v2}
\bibitem{ref81} Prompt Engineering and Calibration for Zero-Shot Commonsense Reasoning. arXiv:2304.06962v1, 2023. \url{https://arxiv.org/abs/2304.06962v1}
\bibitem{ref82} Prompt Design and Engineering Introduction and Advanced Methods. arXiv:2401.14423v3, 2024. \url{https://arxiv.org/abs/2401.14423v3}
\bibitem{ref83} Unnatural language processing How do language models handle machine-generated prompts. arXiv:2310.15829v1, 2023. \url{https://arxiv.org/abs/2310.15829v1}
\bibitem{ref84} Do prompt positions really matter. arXiv:2305.14493v3, 2023. \url{https://arxiv.org/abs/2305.14493v3}
\bibitem{ref85} Improving Knowledge Extraction from LLMs for Task Learning through Agent Analysis. arXiv:2306.06770v4, 2023. \url{https://arxiv.org/abs/2306.06770v4}
\bibitem{ref86} Prompt Engineering for Healthcare Methodologies and Applications. arXiv:2304.14670v2, 2023. \url{https://arxiv.org/abs/2304.14670v2}
\bibitem{ref87} Language Model Prompt Selection via Simulation Optimization. arXiv:2404.08164v1, 2024. \url{https://arxiv.org/abs/2404.08164v1}
\bibitem{ref88} Promptbreeder Self-Referential Self-Improvement Via Prompt Evolution. arXiv:2309.16797v1, 2023. \url{https://arxiv.org/abs/2309.16797v1}
\bibitem{ref89} Multitask Prompt Tuning Enables Parameter-Efficient Transfer Learning. arXiv:2303.02861v1, 2023. \url{https://arxiv.org/abs/2303.02861v1}
\bibitem{ref90} To be or not to be an exploration of continuously controllable prompt engineering. arXiv:2311.09773v1, 2023. \url{https://arxiv.org/abs/2311.09773v1}
\bibitem{ref91} Multi-Modal Subjective Context Modelling and Recognition. arXiv:2011.09671v1, 2020. \url{https://arxiv.org/abs/2011.09671v1}
\bibitem{ref92} Learning and Leveraging World Models in Visual Representation Learning. arXiv:2403.00504v1, 2024. \url{https://arxiv.org/abs/2403.00504v1}
\bibitem{ref93} Modeling and Leveraging Prerequisite Context in Recommendation. arXiv:2209.11471v1, 2022. \url{https://arxiv.org/abs/2209.11471v1}
\bibitem{ref94} Electoral Susceptibility. arXiv:1211.0656v1, 2012. \url{https://arxiv.org/abs/1211.0656v1}
\bibitem{ref95} An Information-theoretic Approach to Prompt Engineering Without Ground Truth Labels. arXiv:2203.11364v1, 2022. \url{https://arxiv.org/abs/2203.11364v1}
\bibitem{ref96} Model-Agnostic Meta-Learning for Fast Adaptation of Deep Networks. arXiv:1703.03400v3, 2017. \url{https://arxiv.org/abs/1703.03400v3}
\bibitem{ref97} Recasting Gradient-Based Meta-Learning as Hierarchical Bayes. arXiv:1801.08930v1, 2018. \url{https://arxiv.org/abs/1801.08930v1}
\bibitem{ref98} Model-Agnostic Learning to Meta-Learn. arXiv:2012.02684v2, 2020. \url{https://arxiv.org/abs/2012.02684v2}
\bibitem{ref99} When MAML Can Adapt Fast and How to Assist When It Cannot. arXiv:1910.13603v3, 2019. \url{https://arxiv.org/abs/1910.13603v3}
\bibitem{ref100} Multimodal Prototypical Networks for Few-shot Learning. arXiv:2011.08899v1, 2020. \url{https://arxiv.org/abs/2011.08899v1}
\bibitem{ref101} How to train your MAML. arXiv:1810.09502v3, 2018. \url{https://arxiv.org/abs/1810.09502v3}
\bibitem{ref102} Alpha MAML Adaptive Model-Agnostic Meta-Learning. arXiv:1905.07435v1, 2019. \url{https://arxiv.org/abs/1905.07435v1}
\bibitem{ref103} Contextual Gradient Scaling for Few-Shot Learning. arXiv:2110.10353v1, 2021. \url{https://arxiv.org/abs/2110.10353v1}
\bibitem{ref104} How Does the Task Landscape Affect MAML Performance. arXiv:2010.14672v5, 2020. \url{https://arxiv.org/abs/2010.14672v5}
\bibitem{ref105} Task-Agnostic Meta-Learning for Few-shot Learning. arXiv:1805.07722v1, 2018. \url{https://arxiv.org/abs/1805.07722v1}
\bibitem{ref106} Multimodal Model-Agnostic Meta-Learning via Task-Aware Modulation. arXiv:1910.13616v1, 2019. \url{https://arxiv.org/abs/1910.13616v1}
\bibitem{ref107} Rapid Learning or Feature Reuse Towards Understanding the Effectiveness of MAML. arXiv:1909.09157v2, 2019. \url{https://arxiv.org/abs/1909.09157v2}
\bibitem{ref108} BOIL Towards Representation Change for Few-shot Learning. arXiv:2008.08882v2, 2020. \url{https://arxiv.org/abs/2008.08882v2}
\bibitem{ref109} Learning from Few Examples A Summary of Approaches to Few-Shot Learning. arXiv:2203.04291v1, 2022. \url{https://arxiv.org/abs/2203.04291v1}
\bibitem{ref110} Meta-learning for Few-shot Natural Language Processing A Survey. arXiv:2007.09604v1, 2020. \url{https://arxiv.org/abs/2007.09604v1}
\bibitem{ref111} Finding Task-Relevant Features for Few-Shot Learning by Category Traversal. arXiv:1905.11116v1, 2019. \url{https://arxiv.org/abs/1905.11116v1}
\bibitem{ref112} Prototype Rectification for Few-Shot Learning. arXiv:1911.10713v4, 2019. \url{https://arxiv.org/abs/1911.10713v4}
\bibitem{ref113} Few-Shot Text Classification with Triplet Networks, Data Augmentation, and Curriculum Learning. arXiv:2103.07552v1, 2021. \url{https://arxiv.org/abs/2103.07552v1}
\bibitem{ref114} Revisiting Fine-tuning for Few-shot Learning. arXiv:1910.00216v2, 2019. \url{https://arxiv.org/abs/1910.00216v2}
\bibitem{ref115} A Study on Representation Transfer for Few-Shot Learning. arXiv:2209.02073v1, 2022. \url{https://arxiv.org/abs/2209.02073v1}
\bibitem{ref116} Partial Is Better Than All Revisiting Fine-tuning Strategy for Few-shot Learning. arXiv:2102.03983v1, 2021. \url{https://arxiv.org/abs/2102.03983v1}
\bibitem{ref117} Boosting Supervision with Self-Supervision for Few-shot Learning. arXiv:1906.07079v1, 2019. \url{https://arxiv.org/abs/1906.07079v1}
\bibitem{ref118} Self-Supervised Prototypical Transfer Learning for Few-Shot Classification. arXiv:2006.11325v1, 2020. \url{https://arxiv.org/abs/2006.11325v1}
\bibitem{ref119} Shot in the Dark Few-Shot Learning with No Base-Class Labels. arXiv:2010.02430v2, 2020. \url{https://arxiv.org/abs/2010.02430v2}
\bibitem{ref120} Scalability in Computing and Robotics. arXiv:2006.04969v2, 2020. \url{https://arxiv.org/abs/2006.04969v2}
\bibitem{ref121} Like a Baby Visually Situated Neural Language Acquisition. arXiv:1805.11546v2, 2018. \url{https://arxiv.org/abs/1805.11546v2}
\bibitem{ref122} Learning Multimodal Word Representation via Dynamic Fusion Methods. arXiv:1801.00532v1, 2018. \url{https://arxiv.org/abs/1801.00532v1}
\bibitem{ref123} Hybrid Contrastive Learning of Tri-Modal Representation for Multimodal Sentiment Analysis. arXiv:2109.01797v1, 2021. \url{https://arxiv.org/abs/2109.01797v1}
\bibitem{ref124} End-to-end Multi-modal Video Temporal Grounding. arXiv:2107.05624v2, 2021. \url{https://arxiv.org/abs/2107.05624v2}
\bibitem{ref125} Causal-CoG A Causal-Effect Look at Context Generation for Boosting Multi-modal Language Models. arXiv:2312.06685v1, 2023. \url{https://arxiv.org/abs/2312.06685v1}
\bibitem{ref126} Understanding Dark Scenes by Contrasting Multi-Modal Observations. arXiv:2308.12320v2, 2023. \url{https://arxiv.org/abs/2308.12320v2}
\bibitem{ref127} Semantic Role Aware Correlation Transformer for Text to Video Retrieval. arXiv:2206.12849v1, 2022. \url{https://arxiv.org/abs/2206.12849v1}
\bibitem{ref128} Training an adaptive dialogue policy for interactive learning of visually grounded word meanings. arXiv:1709.10426v1, 2017. \url{https://arxiv.org/abs/1709.10426v1}
\bibitem{ref129} Multi-modal Visual Understanding with Prompts for Semantic Information Disentanglement of Image. arXiv:2305.09333v1, 2023. \url{https://arxiv.org/abs/2305.09333v1}
\bibitem{ref130} Visual In-Context Learning for Large Vision-Language Models. arXiv:2402.11574v1, 2024. \url{https://arxiv.org/abs/2402.11574v1}
\bibitem{ref131} Multimodal Contrastive Learning via Uni-Modal Coding and Cross-Modal Prediction for Multimodal Sentiment Analysis. arXiv:2210.14556v1, 2022. \url{https://arxiv.org/abs/2210.14556v1}
\bibitem{ref132} Understanding and Improving In-Context Learning on Vision-language Models. arXiv:2311.18021v1, 2023. \url{https://arxiv.org/abs/2311.18021v1}
\bibitem{ref133} Can Text-to-image Model Assist Multi-modal Learning for Visual Recognition with Visual Modality Missing. arXiv:2402.09036v1, 2024. \url{https://arxiv.org/abs/2402.09036v1}
\bibitem{ref134} Temporal Cross-Media Retrieval with Soft-Smoothing. arXiv:1810.04547v1, 2018. \url{https://arxiv.org/abs/1810.04547v1}
\bibitem{ref135} HyperLearn A Distributed Approach for Representation Learning in Datasets With Many Modalities. arXiv:1909.09252v1, 2019. \url{https://arxiv.org/abs/1909.09252v1}
\bibitem{ref136} UNIMO-3 Multi-granularity Interaction for Vision-Language Representation Learning. arXiv:2305.13697v1, 2023. \url{https://arxiv.org/abs/2305.13697v1}
\bibitem{ref137} One-Versus-Others Attention Scalable Multimodal Integration for Clinical Data. arXiv:2307.05435v3, 2023. \url{https://arxiv.org/abs/2307.05435v3}
\bibitem{ref138} Understanding the Multi-modal Prompts of the Pre-trained Vision-Language Model. arXiv:2312.11570v3, 2023. \url{https://arxiv.org/abs/2312.11570v3}
\bibitem{ref139} Putting visual object recognition in context. arXiv:1911.07349v3, 2019. \url{https://arxiv.org/abs/1911.07349v3}
\bibitem{ref140} Text-centric Alignment for Multi-Modality Learning. arXiv:2402.08086v1, 2024. \url{https://arxiv.org/abs/2402.08086v1}
\bibitem{ref141} Context Understanding in Computer Vision A Survey. arXiv:2302.05011v1, 2023. \url{https://arxiv.org/abs/2302.05011v1}
\bibitem{ref142} Modeling Text-visual Mutual Dependency for Multi-modal Dialog Generation. arXiv:2105.14445v1, 2021. \url{https://arxiv.org/abs/2105.14445v1}
\bibitem{ref143} FSMR A Feature Swapping Multi-modal Reasoning Approach with Joint Textual and Visual Clues. arXiv:2403.20026v1, 2024. \url{https://arxiv.org/abs/2403.20026v1}
\bibitem{ref144} REX Reasoning-aware and Grounded Explanation. arXiv:2203.06107v1, 2022. \url{https://arxiv.org/abs/2203.06107v1}
\bibitem{ref145} Missing Modality Robustness in Semi-Supervised Multi-Modal Semantic Segmentation. arXiv:2304.10756v1, 2023. \url{https://arxiv.org/abs/2304.10756v1}
\bibitem{ref146} Probing Multimodal Embeddings for Linguistic Properties the Visual-Semantic Case. arXiv:2102.11115v1, 2021. \url{https://arxiv.org/abs/2102.11115v1}
\bibitem{ref147} Compositional Mixture Representations for Vision and Text. arXiv:2206.06404v1, 2022. \url{https://arxiv.org/abs/2206.06404v1}
\bibitem{ref148} IMProv Inpainting-based Multimodal Prompting for Computer Vision Tasks. arXiv:2312.01771v1, 2023. \url{https://arxiv.org/abs/2312.01771v1}
\bibitem{ref149} Investigating the Role of Attribute Context in Vision-Language Models for Object Recognition and Detection. arXiv:2303.10093v2, 2023. \url{https://arxiv.org/abs/2303.10093v2}
\bibitem{ref150} Efficient Object-Level Visual Context Modeling for Multimodal Machine Translation Masking Irrelevant Objects Helps Grounding. arXiv:2101.05208v1, 2020. \url{https://arxiv.org/abs/2101.05208v1}
\bibitem{ref151} Moments in Time Dataset one million videos for event understanding. arXiv:1801.03150v3, 2018. \url{https://arxiv.org/abs/1801.03150v3}
\bibitem{ref152} Technical Report Temporal Aggregate Representations. arXiv:2106.03152v2, 2021. \url{https://arxiv.org/abs/2106.03152v2}
\bibitem{ref153} Only Time Can Tell Discovering Temporal Data for Temporal Modeling. arXiv:1907.08340v2, 2019. \url{https://arxiv.org/abs/1907.08340v2}
\bibitem{ref154} Generic Temporal Reasoning with Differential Analysis and Explanation. arXiv:2212.10467v2, 2022. \url{https://arxiv.org/abs/2212.10467v2}
\bibitem{ref155} ATM Action Temporality Modeling for Video Question Answering. arXiv:2309.02290v1, 2023. \url{https://arxiv.org/abs/2309.02290v1}
\bibitem{ref156} Learning Fine-Grained Visual Understanding for Video Question Answering via Decoupling Spatial-Temporal Modeling. arXiv:2210.03941v1, 2022. \url{https://arxiv.org/abs/2210.03941v1}
\bibitem{ref157} Can Temporal Information Help with Contrastive Self-Supervised Learning. arXiv:2011.13046v1, 2020. \url{https://arxiv.org/abs/2011.13046v1}
\bibitem{ref158} Time Is MattEr Temporal Self-supervision for Video Transformers. arXiv:2207.09067v1, 2022. \url{https://arxiv.org/abs/2207.09067v1}
\bibitem{ref159} Test of Time Instilling Video-Language Models with a Sense of Time. arXiv:2301.02074v2, 2023. \url{https://arxiv.org/abs/2301.02074v2}
\bibitem{ref160} Revisiting the Video in Video-Language Understanding. arXiv:2206.01720v1, 2022. \url{https://arxiv.org/abs/2206.01720v1}
\bibitem{ref161} Unified Graph Structured Models for Video Understanding. arXiv:2103.15662v1, 2021. \url{https://arxiv.org/abs/2103.15662v1}
\bibitem{ref162} Videos as Space-Time Region Graphs. arXiv:1806.01810v2, 2018. \url{https://arxiv.org/abs/1806.01810v2}
\bibitem{ref163} iReason Multimodal Commonsense Reasoning using Videos and Natural Language with Interpretability. arXiv:2107.10300v1, 2021. \url{https://arxiv.org/abs/2107.10300v1}
\bibitem{ref164} Neuro-Symbolic Video Search. arXiv:2403.11021v1, 2024. \url{https://arxiv.org/abs/2403.11021v1}
\bibitem{ref165} Spatio-temporal Tendency Reasoning for Human Body Pose and Shape Estimation from Videos. arXiv:2210.03659v2, 2022. \url{https://arxiv.org/abs/2210.03659v2}
\bibitem{ref166} CLEVRER CoLlision Events for Video REpresentation and Reasoning. arXiv:1910.01442v2, 2019. \url{https://arxiv.org/abs/1910.01442v2}
\bibitem{ref167} Cross-Modality Attentive Feature Fusion for Object Detection in Multispectral Remote Sensing Imagery. arXiv:2112.02991v1, 2021. \url{https://arxiv.org/abs/2112.02991v1}
\bibitem{ref168} Multi-modal Semantic Understanding with Contrastive Cross-modal Feature Alignment. arXiv:2403.06355v1, 2024. \url{https://arxiv.org/abs/2403.06355v1}
\bibitem{ref169} MSAF Multimodal Split Attention Fusion. arXiv:2012.07175v2, 2020. \url{https://arxiv.org/abs/2012.07175v2}
\bibitem{ref170} Cross-Attention is Not Enough Incongruity-Aware Dynamic Hierarchical Fusion for Multimodal Affect Recognition. arXiv:2305.13583v4, 2023. \url{https://arxiv.org/abs/2305.13583v4}
\bibitem{ref171} A Joint Cross-Attention Model for Audio-Visual Fusion in Dimensional Emotion Recognition. arXiv:2203.14779v3, 2022. \url{https://arxiv.org/abs/2203.14779v3}
\bibitem{ref172} Context-Based Multimodal Fusion. arXiv:2403.04650v2, 2024. \url{https://arxiv.org/abs/2403.04650v2}
\bibitem{ref173} Multimodal Hyperspectral Image Classification via Interconnected Fusion. arXiv:2304.00495v1, 2023. \url{https://arxiv.org/abs/2304.00495v1}
\bibitem{ref174} A Self-Adjusting Fusion Representation Learning Model for Unaligned Text-Audio Sequences. arXiv:2212.11772v1, 2022. \url{https://arxiv.org/abs/2212.11772v1}
\bibitem{ref175} Adaptive Fusion Techniques for Multimodal Data. arXiv:1911.03821v2, 2019. \url{https://arxiv.org/abs/1911.03821v2}
\bibitem{ref176} Multi-channel Attentive Graph Convolutional Network With Sentiment Fusion For Multimodal Sentiment Analysis. arXiv:2201.10274v1, 2022. \url{https://arxiv.org/abs/2201.10274v1}
\bibitem{ref177} Sequential Cross Attention Based Multi-task Learning. arXiv:2209.02518v1, 2022. \url{https://arxiv.org/abs/2209.02518v1}
\bibitem{ref178} Towards Robust Multimodal Prompting With Missing Modalities. arXiv:2312.15890v2, 2023. \url{https://arxiv.org/abs/2312.15890v2}
\bibitem{ref179} Multimodal Prompting with Missing Modalities for Visual Recognition. arXiv:2303.03369v2, 2023. \url{https://arxiv.org/abs/2303.03369v2}
\bibitem{ref180} Learning Domain Invariant Prompt for Vision-Language Models. arXiv:2212.04196v2, 2022. \url{https://arxiv.org/abs/2212.04196v2}
\bibitem{ref181} Instruction-ViT Multi-Modal Prompts for Instruction Learning in ViT. arXiv:2305.00201v1, 2023. \url{https://arxiv.org/abs/2305.00201v1}
\bibitem{ref182} ContextVis Envision Contextual Learning and Interaction with Generative Models. arXiv:2403.12768v1, 2024. \url{https://arxiv.org/abs/2403.12768v1}
\bibitem{ref183} S3Eval A Synthetic, Scalable, Systematic Evaluation Suite for Large Language Models. arXiv:2310.15147v2, 2023. \url{https://arxiv.org/abs/2310.15147v2}
\bibitem{ref184} ConTextual Evaluating Context-Sensitive Text-Rich Visual Reasoning in Large Multimodal Models. arXiv:2401.13311v1, 2024. \url{https://arxiv.org/abs/2401.13311v1}
\bibitem{ref185} RADDLE An Evaluation Benchmark and Analysis Platform for Robust Task-oriented Dialog Systems. arXiv:2012.14666v1, 2020. \url{https://arxiv.org/abs/2012.14666v1}
\bibitem{ref186} Benchmark Environments for Multitask Learning in Continuous Domains. arXiv:1708.04352v1, 2017. \url{https://arxiv.org/abs/1708.04352v1}
\bibitem{ref187} Evaluation Methods and Measures for Causal Learning Algorithms. arXiv:2202.02896v1, 2022. \url{https://arxiv.org/abs/2202.02896v1}
\bibitem{ref188} Assessing Generalization for Subpopulation Representative Modeling via In-Context Learning. arXiv:2402.07368v1, 2024. \url{https://arxiv.org/abs/2402.07368v1}
\bibitem{ref189} Identifying the Context Shift between Test Benchmarks and Production Data. arXiv:2207.01059v2, 2022. \url{https://arxiv.org/abs/2207.01059v2}
\bibitem{ref190} Data-SUITE Data-centric identification of in-distribution incongruous examples. arXiv:2202.08836v3, 2022. \url{https://arxiv.org/abs/2202.08836v3}
\bibitem{ref191} Pre-trained Vision and Language Transformers Are Few-Shot Incremental Learners. arXiv:2404.02117v1, 2024. \url{https://arxiv.org/abs/2404.02117v1}
\bibitem{ref192} Language Models are Few-Shot Learners. arXiv:2005.14165v4, 2020. \url{https://arxiv.org/abs/2005.14165v4}
\bibitem{ref193} PaLM Scaling Language Modeling with Pathways. arXiv:2204.02311v5, 2022. \url{https://arxiv.org/abs/2204.02311v5}
\bibitem{ref194} In-Context Learning Dynamics with Random Binary Sequences. arXiv:2310.17639v3, 2023. \url{https://arxiv.org/abs/2310.17639v3}
\bibitem{ref195} Post Hoc Explanations of Language Models Can Improve Language Models. arXiv:2305.11426v3, 2023. \url{https://arxiv.org/abs/2305.11426v3}
\bibitem{ref196} Complementary Explanations for Effective In-Context Learning. arXiv:2211.13892v2, 2022. \url{https://arxiv.org/abs/2211.13892v2}
\bibitem{ref197} Siren's Song in the AI Ocean A Survey on Hallucination in Large Language Models. arXiv:2309.01219v2, 2023. \url{https://arxiv.org/abs/2309.01219v2}
\bibitem{ref198} Fine-tuning Language Models for Factuality. arXiv:2311.08401v1, 2023. \url{https://arxiv.org/abs/2311.08401v1}
\bibitem{ref199} Unveiling the Generalization Power of Fine-Tuned Large Language Models. arXiv:2403.09162v1, 2024. \url{https://arxiv.org/abs/2403.09162v1}
\bibitem{ref200} Reinforcement Learning Fine-tuning of Language Models is Biased Towards More Extractable Features. arXiv:2311.04046v1, 2023. \url{https://arxiv.org/abs/2311.04046v1}
\bibitem{ref201} In-context learning agents are asymmetric belief updaters. arXiv:2402.03969v1, 2024. \url{https://arxiv.org/abs/2402.03969v1}
\bibitem{ref202} Can Large Language Models Explain Themselves A Study of LLM-Generated Self-Explanations. arXiv:2310.11207v1, 2023. \url{https://arxiv.org/abs/2310.11207v1}
\bibitem{ref203} The Unreliability of Explanations in Few-shot Prompting for Textual Reasoning. arXiv:2205.03401v2, 2022. \url{https://arxiv.org/abs/2205.03401v2}
\bibitem{ref204} Self-AMPLIFY Improving Small Language Models with Self Post Hoc Explanations. arXiv:2402.12038v2, 2024. \url{https://arxiv.org/abs/2402.12038v2}
\bibitem{ref205} Explanations from Large Language Models Make Small Reasoners Better. arXiv:2210.06726v1, 2022. \url{https://arxiv.org/abs/2210.06726v1}
\bibitem{ref206} ExaRanker Explanation-Augmented Neural Ranker. arXiv:2301.10521v2, 2023. \url{https://arxiv.org/abs/2301.10521v2}
\bibitem{ref207} Towards Benchmarking the Utility of Explanations for Model Debugging. arXiv:2105.04505v1, 2021. \url{https://arxiv.org/abs/2105.04505v1}
\bibitem{ref208} Are Large Language Models Post Hoc Explainers. arXiv:2310.05797v3, 2023. \url{https://arxiv.org/abs/2310.05797v3}
\bibitem{ref209} Large Language Models As Faithful Explainers. arXiv:2402.04678v1, 2024. \url{https://arxiv.org/abs/2402.04678v1}
\bibitem{ref210} A Hypothesis-Driven Framework for the Analysis of Self-Rationalising Models. arXiv:2402.04787v1, 2024. \url{https://arxiv.org/abs/2402.04787v1}
\bibitem{ref211} Do Large Language Models Understand Logic or Just Mimick Context. arXiv:2402.12091v1, 2024. \url{https://arxiv.org/abs/2402.12091v1}
\bibitem{ref212} CLOMO Counterfactual Logical Modification with Large Language Models. arXiv:2311.17438v3, 2023. \url{https://arxiv.org/abs/2311.17438v3}
\bibitem{ref213} LLM Inference Unveiled Survey and Roofline Model Insights. arXiv:2402.16363v5, 2024. \url{https://arxiv.org/abs/2402.16363v5}
\bibitem{ref214} On the Compressibility of Quantized Large Language Models. arXiv:2403.01384v1, 2024. \url{https://arxiv.org/abs/2403.01384v1}
\bibitem{ref215} Compresso Structured Pruning with Collaborative Prompting Learns Compact Large Language Models. arXiv:2310.05015v2, 2023. \url{https://arxiv.org/abs/2310.05015v2}
\bibitem{ref216} Synthetic Data Generation in Low-Resource Settings via Fine-Tuning of Large Language Models. arXiv:2310.01119v2, 2023. \url{https://arxiv.org/abs/2310.01119v2}
\bibitem{ref217} Retrieval-based Knowledge Transfer An Effective Approach for Extreme Large Language Model Compression. arXiv:2310.15594v1, 2023. \url{https://arxiv.org/abs/2310.15594v1}
\bibitem{ref218} Scaling In-Context Demonstrations with Structured Attention. arXiv:2307.02690v1, 2023. \url{https://arxiv.org/abs/2307.02690v1}
\bibitem{ref219} Context Compression for Auto-regressive Transformers with Sentinel Tokens. arXiv:2310.08152v2, 2023. \url{https://arxiv.org/abs/2310.08152v2}
\bibitem{ref220} Can language models learn from explanations in context. arXiv:2204.02329v4, 2022. \url{https://arxiv.org/abs/2204.02329v4}
\bibitem{ref221} Exploring Automatically Perturbed Natural Language Explanations in Relation Extraction. arXiv:2305.15520v1, 2023. \url{https://arxiv.org/abs/2305.15520v1}
\bibitem{ref222} Explanation Selection Using Unlabeled Data for Chain-of-Thought Prompting. arXiv:2302.04813v3, 2023. \url{https://arxiv.org/abs/2302.04813v3}
\bibitem{ref223} Rationale-Augmented Ensembles in Language Models. arXiv:2207.00747v1, 2022. \url{https://arxiv.org/abs/2207.00747v1}
\bibitem{ref224} Curated LLM Synergy of LLMs and Data Curation for tabular augmentation in ultra low-data regimes. arXiv:2312.12112v2, 2023. \url{https://arxiv.org/abs/2312.12112v2}
\bibitem{ref225} LLM2LLM Boosting LLMs with Novel Iterative Data Enhancement. arXiv:2403.15042v1, 2024. \url{https://arxiv.org/abs/2403.15042v1}
\bibitem{ref226} Generative AI for Synthetic Data Generation Methods, Challenges and the Future. arXiv:2403.04190v1, 2024. \url{https://arxiv.org/abs/2403.04190v1}
\bibitem{ref227} Exploring Self-supervised Logic-enhanced Training for Large Language Models. arXiv:2305.13718v6, 2023. \url{https://arxiv.org/abs/2305.13718v6}
\bibitem{ref228} Many-Shot In-Context Learning. arXiv:2404.11018v1, 2024. \url{https://arxiv.org/abs/2404.11018v1}
\bibitem{ref229} Beyond the Limits A Survey of Techniques to Extend the Context Length in Large Language Models. arXiv:2402.02244v1, 2024. \url{https://arxiv.org/abs/2402.02244v1}
\bibitem{ref230} Attention-Driven Reasoning Unlocking the Potential of Large Language Models. arXiv:2403.14932v2, 2024. \url{https://arxiv.org/abs/2403.14932v2}
\bibitem{ref231} TriBERT Full-body Human-centric Audio-visual Representation Learning for Visual Sound Separation. arXiv:2110.13412v1, 2021. \url{https://arxiv.org/abs/2110.13412v1}
\bibitem{ref232} Foundational Models Defining a New Era in Vision A Survey and Outlook. arXiv:2307.13721v1, 2023. \url{https://arxiv.org/abs/2307.13721v1}
\bibitem{ref233} LoGoPrompt Synthetic Text Images Can Be Good Visual Prompts for Vision-Language Models. arXiv:2309.01155v2, 2023. \url{https://arxiv.org/abs/2309.01155v2}
\bibitem{ref234} A Cooperative Group Optimization System. arXiv:1808.01342v1, 2018. \url{https://arxiv.org/abs/1808.01342v1}
\bibitem{ref235} Prompting through Prototype A Prototype-based Prompt Learning on Pretrained Vision-Language Models. arXiv:2210.10841v1, 2022. \url{https://arxiv.org/abs/2210.10841v1}
\bibitem{ref236} Mastering Robot Manipulation with Multimodal Prompts through Pretraining and Multi-task Fine-tuning. arXiv:2310.09676v1, 2023. \url{https://arxiv.org/abs/2310.09676v1}
\bibitem{ref237} Retrieval-Enhanced Visual Prompt Learning for Few-shot Classification. arXiv:2306.02243v1, 2023. \url{https://arxiv.org/abs/2306.02243v1}
\bibitem{ref238} Learning Expressive Prompting With Residuals for Vision Transformers. arXiv:2303.15591v1, 2023. \url{https://arxiv.org/abs/2303.15591v1}
\bibitem{ref239} Why Can GPT Learn In-Context Language Models Implicitly Perform Gradient Descent as Meta-Optimizers. arXiv:2212.10559v3, 2022. \url{https://arxiv.org/abs/2212.10559v3}
\bibitem{ref240} Scratching Visual Transformer's Back with Uniform Attention. arXiv:2210.08457v1, 2022. \url{https://arxiv.org/abs/2210.08457v1}
\bibitem{ref241} MetaVL Transferring In-Context Learning Ability From Language Models to Vision-Language Models. arXiv:2306.01311v1, 2023. \url{https://arxiv.org/abs/2306.01311v1}
\bibitem{ref242} Attentive Feature Reuse for Multi Task Meta learning. arXiv:2006.07438v1, 2020. \url{https://arxiv.org/abs/2006.07438v1}
\bibitem{ref243} Is Cross-Attention Preferable to Self-Attention for Multi-Modal Emotion Recognition. arXiv:2202.09263v1, 2022. \url{https://arxiv.org/abs/2202.09263v1}
\bibitem{ref244} Deeply Coupled Cross-Modal Prompt Learning. arXiv:2305.17903v3, 2023. \url{https://arxiv.org/abs/2305.17903v3}
\bibitem{ref245} IGLUE A Benchmark for Transfer Learning across Modalities, Tasks, and Languages. arXiv:2201.11732v2, 2022. \url{https://arxiv.org/abs/2201.11732v2}
\bibitem{ref246} MMBench Is Your Multi-modal Model an All-around Player. arXiv:2307.06281v3, 2023. \url{https://arxiv.org/abs/2307.06281v3}
\bibitem{ref247} Images Speak in Images A Generalist Painter for In-Context Visual Learning. arXiv:2212.02499v2, 2022. \url{https://arxiv.org/abs/2212.02499v2}
\bibitem{ref248} BLINK Multimodal Large Language Models Can See but Not Perceive. arXiv:2404.12390v2, 2024. \url{https://arxiv.org/abs/2404.12390v2}
\bibitem{ref249} VALUED -- Vision and Logical Understanding Evaluation Dataset. arXiv:2311.12610v2, 2023. \url{https://arxiv.org/abs/2311.12610v2}
\bibitem{ref250} ContextRef Evaluating Referenceless Metrics For Image Description Generation. arXiv:2309.11710v1, 2023. \url{https://arxiv.org/abs/2309.11710v1}
\bibitem{ref251} Towards Reliable Assessments of Demographic Disparities in Multi-Label Image Classifiers. arXiv:2302.08572v1, 2023. \url{https://arxiv.org/abs/2302.08572v1}
\bibitem{ref252} Active Testing An Efficient and Robust Framework for Estimating Accuracy. arXiv:1807.00493v1, 2018. \url{https://arxiv.org/abs/1807.00493v1}
\bibitem{ref253} In-context learning enables multimodal large language models to classify cancer pathology images. arXiv:2403.07407v1, 2024. \url{https://arxiv.org/abs/2403.07407v1}
\bibitem{ref254} Domain Generalization by Learning from Privileged Medical Imaging Information. arXiv:2311.05861v1, 2023. \url{https://arxiv.org/abs/2311.05861v1}
\bibitem{ref255} On the Road with GPT-4V(ision) Early Explorations of Visual-Language Model on Autonomous Driving. arXiv:2311.05332v2, 2023. \url{https://arxiv.org/abs/2311.05332v2}
\bibitem{ref256} Light the Night A Multi-Condition Diffusion Framework for Unpaired Low-Light Enhancement in Autonomous Driving. arXiv:2404.04804v1, 2024. \url{https://arxiv.org/abs/2404.04804v1}
\bibitem{ref257} Causal Scene BERT Improving object detection by searching for challenging groups of data. arXiv:2202.03651v2, 2022. \url{https://arxiv.org/abs/2202.03651v2}
\bibitem{ref258} Multiperspective Teaching of Unknown Objects via Shared-gaze-based Multimodal Human-Robot Interaction. arXiv:2303.00423v1, 2023. \url{https://arxiv.org/abs/2303.00423v1}
\bibitem{ref259} Vision-based Robot Manipulation Learning via Human Demonstrations. arXiv:2003.00385v1, 2020. \url{https://arxiv.org/abs/2003.00385v1}
\bibitem{ref260} Survey of Social Bias in Vision-Language Models. arXiv:2309.14381v1, 2023. \url{https://arxiv.org/abs/2309.14381v1}
\bibitem{ref261} Exploring Relational Context for Multi-Task Dense Prediction. arXiv:2104.13874v2, 2021. \url{https://arxiv.org/abs/2104.13874v2}
\bibitem{ref262} Neural Operator Learning Maps Between Function Spaces. arXiv:2108.08481v5, 2021. \url{https://arxiv.org/abs/2108.08481v5}
\bibitem{ref263} FMint Bridging Human Designed and Data Pretrained Models for Differential Equation Foundation Model. arXiv:2404.14688v1, 2024. \url{https://arxiv.org/abs/2404.14688v1}
\bibitem{ref264} Encoding physics to learn reaction-diffusion processes. arXiv:2106.04781v2, 2021. \url{https://arxiv.org/abs/2106.04781v2}
\bibitem{ref265} Large-scale Neural Solvers for Partial Differential Equations. arXiv:2009.03730v1, 2020. \url{https://arxiv.org/abs/2009.03730v1}
\bibitem{ref266} Deep-learning enhancement of large scale numerical simulations. arXiv:2004.03454v1, 2020. \url{https://arxiv.org/abs/2004.03454v1}
\bibitem{ref267} Accelerating Physics Simulations with TPUs An Inundation Modeling Example. arXiv:2204.10323v1, 2022. \url{https://arxiv.org/abs/2204.10323v1}
\bibitem{ref268} Pretrained Semantic Speech Embeddings for End-to-End Spoken Language Understanding via Cross-Modal Teacher-Student Learning. arXiv:2007.01836v2, 2020. \url{https://arxiv.org/abs/2007.01836v2}
\bibitem{ref269} BERT Meets CTC New Formulation of End-to-End Speech Recognition with Pre-trained Masked Language Model. arXiv:2210.16663v2, 2022. \url{https://arxiv.org/abs/2210.16663v2}
\bibitem{ref270} Learning Speech Representation From Contrastive Token-Acoustic Pretraining. arXiv:2309.00424v5, 2023. \url{https://arxiv.org/abs/2309.00424v5}
\bibitem{ref271} Bridging Speech and Textual Pre-trained Models with Unsupervised ASR. arXiv:2211.03025v1, 2022. \url{https://arxiv.org/abs/2211.03025v1}
\bibitem{ref272} Contrastive Learning for Improving ASR Robustness in Spoken Language Understanding. arXiv:2205.00693v2, 2022. \url{https://arxiv.org/abs/2205.00693v2}
\bibitem{ref273} Long-span language modeling for speech recognition. arXiv:1911.04571v1, 2019. \url{https://arxiv.org/abs/1911.04571v1}
\bibitem{ref274} Speech-Text Dialog Pre-training for Spoken Dialog Understanding with Explicit Cross-Modal Alignment. arXiv:2305.11579v2, 2023. \url{https://arxiv.org/abs/2305.11579v2}
\bibitem{ref275} Contextual Audio-Visual Switching For Speech Enhancement in Real-World Environments. arXiv:1808.09825v1, 2018. \url{https://arxiv.org/abs/1808.09825v1}
\bibitem{ref276} Universal Paralinguistic Speech Representations Using Self-Supervised Conformers. arXiv:2110.04621v4, 2021. \url{https://arxiv.org/abs/2110.04621v4}
\bibitem{ref277} Improving RNN-T ASR Performance with Date-Time and Location Awareness. arXiv:2106.06183v2, 2021. \url{https://arxiv.org/abs/2106.06183v2}
\bibitem{ref278} Can Large Language Models Understand Context. arXiv:2402.00858v1, 2024. \url{https://arxiv.org/abs/2402.00858v1}
\bibitem{ref279} A Pragmatics-Centered Evaluation Framework for Natural Language Understanding. arXiv:1907.08672v2, 2019. \url{https://arxiv.org/abs/1907.08672v2}
\bibitem{ref280} KOBEST Korean Balanced Evaluation of Significant Tasks. arXiv:2204.04541v1, 2022. \url{https://arxiv.org/abs/2204.04541v1}
\bibitem{ref281} Resources and Few-shot Learners for In-context Learning in Slavic Languages. arXiv:2304.01922v1, 2023. \url{https://arxiv.org/abs/2304.01922v1}
\bibitem{ref282} On the Effect of Pretraining Corpora on In-context Learning by a Large-scale Language Model. arXiv:2204.13509v2, 2022. \url{https://arxiv.org/abs/2204.13509v2}
\bibitem{ref283} The ICL Consistency Test. arXiv:2312.04945v1, 2023. \url{https://arxiv.org/abs/2312.04945v1}
\bibitem{ref284} CausalBench A Comprehensive Benchmark for Causal Learning Capability of Large Language Models. arXiv:2404.06349v1, 2024. \url{https://arxiv.org/abs/2404.06349v1}
\bibitem{ref285} Few-shot Fine-tuning vs. In-context Learning A Fair Comparison and Evaluation. arXiv:2305.16938v2, 2023. \url{https://arxiv.org/abs/2305.16938v2}
\bibitem{ref286} Evaluating the Evaluators Are Current Few-Shot Learning Benchmarks Fit for Purpose. arXiv:2307.02732v1, 2023. \url{https://arxiv.org/abs/2307.02732v1}
\bibitem{ref287} On (Normalised) Discounted Cumulative Gain as an Off-Policy Evaluation Metric for Top-\$n\$ Recommendation. arXiv:2307.15053v2, 2023. \url{https://arxiv.org/abs/2307.15053v2}
\bibitem{ref288} Multi-Dimensional Evaluation of Text Summarization with In-Context Learning. arXiv:2306.01200v1, 2023. \url{https://arxiv.org/abs/2306.01200v1}
\bibitem{ref289} Identifying and Benchmarking Natural Out-of-Context Prediction Problems. arXiv:2110.13223v1, 2021. \url{https://arxiv.org/abs/2110.13223v1}
\bibitem{ref290} Self-Improving-Leaderboard(SIL) A Call for Real-World Centric Natural Language Processing Leaderboards. arXiv:2303.10888v1, 2023. \url{https://arxiv.org/abs/2303.10888v1}
\bibitem{ref291} Fairness-guided Few-shot Prompting for Large Language Models. arXiv:2303.13217v3, 2023. \url{https://arxiv.org/abs/2303.13217v3}
\bibitem{ref292} Addressing Order Sensitivity of In-Context Demonstration Examples in Causal Language Models. arXiv:2402.15637v1, 2024. \url{https://arxiv.org/abs/2402.15637v1}
\bibitem{ref293} Significant Improvements over the State of the Art A Case Study of the MS MARCO Document Ranking Leaderboard. arXiv:2102.12887v1, 2021. \url{https://arxiv.org/abs/2102.12887v1}
\bibitem{ref294} Evaluation of Synthetic Datasets for Conversational Recommender Systems. arXiv:2212.08167v1, 2022. \url{https://arxiv.org/abs/2212.08167v1}
\bibitem{ref295} Utility is in the Eye of the User A Critique of NLP Leaderboards. arXiv:2009.13888v4, 2020. \url{https://arxiv.org/abs/2009.13888v4}
\bibitem{ref296} Emergent Abilities in Reduced-Scale Generative Language Models. arXiv:2404.02204v1, 2024. \url{https://arxiv.org/abs/2404.02204v1}
\bibitem{ref297} Parallel Structures in Pre-training Data Yield In-Context Learning. arXiv:2402.12530v1, 2024. \url{https://arxiv.org/abs/2402.12530v1}
\bibitem{ref298} Is attention required for ICL Exploring the Relationship Between Model Architecture and In-Context Learning Ability. arXiv:2310.08049v3, 2023. \url{https://arxiv.org/abs/2310.08049v3}
\bibitem{ref299} Understanding In-Context Learning via Supportive Pretraining Data. arXiv:2306.15091v1, 2023. \url{https://arxiv.org/abs/2306.15091v1}
\bibitem{ref300} Understanding the role of FFNs in driving multilingual behaviour in LLMs. arXiv:2404.13855v1, 2024. \url{https://arxiv.org/abs/2404.13855v1}
\bibitem{ref301} ClinicalGPT Large Language Models Finetuned with Diverse Medical Data and Comprehensive Evaluation. arXiv:2306.09968v1, 2023. \url{https://arxiv.org/abs/2306.09968v1}
\bibitem{ref302} MedInsight A Multi-Source Context Augmentation Framework for Generating Patient-Centric Medical Responses using Large Language Models. arXiv:2403.08607v1, 2024. \url{https://arxiv.org/abs/2403.08607v1}
\bibitem{ref303} Precedent-Enhanced Legal Judgment Prediction with LLM and Domain-Model Collaboration. arXiv:2310.09241v1, 2023. \url{https://arxiv.org/abs/2310.09241v1}
\bibitem{ref304} Towards Automatic Evaluation for LLMs' Clinical Capabilities Metric, Data, and Algorithm. arXiv:2403.16446v1, 2024. \url{https://arxiv.org/abs/2403.16446v1}
\bibitem{ref305} SynBench Task-Agnostic Benchmarking of Pretrained Representations using Synthetic Data. arXiv:2210.02989v2, 2022. \url{https://arxiv.org/abs/2210.02989v2}
\bibitem{ref306} Evaluation of Categorical Generative Models -- Bridging the Gap Between Real and Synthetic Data. arXiv:2210.16405v1, 2022. \url{https://arxiv.org/abs/2210.16405v1}
\bibitem{ref307} Synthesising Multi-Modal Minority Samples for Tabular Data. arXiv:2105.08204v1, 2021. \url{https://arxiv.org/abs/2105.08204v1}
\bibitem{ref308} Towards Task Sampler Learning for Meta-Learning. arXiv:2307.08924v3, 2023. \url{https://arxiv.org/abs/2307.08924v3}
\bibitem{ref309} DERAIL Diagnostic Environments for Reward And Imitation Learning. arXiv:2012.01365v1, 2020. \url{https://arxiv.org/abs/2012.01365v1}
\bibitem{ref310} A Closer Look at In-Context Learning under Distribution Shifts. arXiv:2305.16704v1, 2023. \url{https://arxiv.org/abs/2305.16704v1}
\bibitem{ref311} A Fine-Grained Analysis on Distribution Shift. arXiv:2110.11328v2, 2021. \url{https://arxiv.org/abs/2110.11328v2}
\bibitem{ref312} Revisiting Out-of-distribution Robustness in NLP Benchmark, Analysis, and LLMs Evaluations. arXiv:2306.04618v2, 2023. \url{https://arxiv.org/abs/2306.04618v2}
\bibitem{ref313} A Human-Centered Data-Driven Planner-Actor-Critic Architecture via Logic Programming. arXiv:1909.09209v1, 2019. \url{https://arxiv.org/abs/1909.09209v1}
\bibitem{ref314} A Human-in-the-loop Framework to Construct Context-aware Mathematical Notions of Outcome Fairness. arXiv:1911.03020v2, 2019. \url{https://arxiv.org/abs/1911.03020v2}
\bibitem{ref315} A Rationale-Centric Framework for Human-in-the-loop Machine Learning. arXiv:2203.12918v1, 2022. \url{https://arxiv.org/abs/2203.12918v1}
\bibitem{ref316} Identification of Tasks, Datasets, Evaluation Metrics, and Numeric Scores for Scientific Leaderboards Construction. arXiv:1906.09317v1, 2019. \url{https://arxiv.org/abs/1906.09317v1}
\bibitem{ref317} Quantifying Language Models' Sensitivity to Spurious Features in Prompt Design or How I learned to start worrying about prompt formatting. arXiv:2310.11324v1, 2023. \url{https://arxiv.org/abs/2310.11324v1}
\bibitem{ref318} Coverage-based Example Selection for In-Context Learning. arXiv:2305.14907v3, 2023. \url{https://arxiv.org/abs/2305.14907v3}
\bibitem{ref319} Identifying and Analyzing Task-Encoding Tokens in Large Language Models. arXiv:2401.11323v2, 2024. \url{https://arxiv.org/abs/2401.11323v2}
\bibitem{ref320} Bias in Machine Learning Software Why How What to do. arXiv:2105.12195v3, 2021. \url{https://arxiv.org/abs/2105.12195v3}
\bibitem{ref321} Learning what and where to attend. arXiv:1805.08819v4, 2018. \url{https://arxiv.org/abs/1805.08819v4}
\bibitem{ref322} Fairness Evaluation in Text Classification Machine Learning Practitioner Perspectives of Individual and Group Fairness. arXiv:2303.00673v1, 2023. \url{https://arxiv.org/abs/2303.00673v1}
\bibitem{ref323} Bias and Fairness on Multimodal Emotion Detection Algorithms. arXiv:2205.08383v1, 2022. \url{https://arxiv.org/abs/2205.08383v1}
\bibitem{ref324} Dynamic fairness - Breaking vicious cycles in automatic decision making. arXiv:1902.00375v2, 2019. \url{https://arxiv.org/abs/1902.00375v2}
\bibitem{ref325} Prisoners of Their Own Devices How Models Induce Data Bias in Performative Prediction. arXiv:2206.13183v1, 2022. \url{https://arxiv.org/abs/2206.13183v1}
\bibitem{ref326} Factors Influencing Perceived Fairness in Algorithmic Decision-Making Algorithm Outcomes, Development Procedures, and Individual Differences. arXiv:2001.09604v1, 2020. \url{https://arxiv.org/abs/2001.09604v1}
\bibitem{ref327} Fairness Explainability using Optimal Transport with Applications in Image Classification. arXiv:2308.11090v2, 2023. \url{https://arxiv.org/abs/2308.11090v2}
\bibitem{ref328} FITNESS A Causal De-correlation Approach for Mitigating Bias in Machine Learning Software. arXiv:2305.14396v1, 2023. \url{https://arxiv.org/abs/2305.14396v1}
\bibitem{ref329} Evaluating Fairness Metrics in the Presence of Dataset Bias. arXiv:1809.09245v1, 2018. \url{https://arxiv.org/abs/1809.09245v1}
\bibitem{ref330} Monitoring Algorithmic Fairness. arXiv:2305.15979v1, 2023. \url{https://arxiv.org/abs/2305.15979v1}
\bibitem{ref331} Structured Prompting Scaling In-Context Learning to 1,000 Examples. arXiv:2212.06713v1, 2022. \url{https://arxiv.org/abs/2212.06713v1}
\bibitem{ref332} An Experimental Evaluation of Machine Learning Training on a Real Processing-in-Memory System. arXiv:2207.07886v5, 2022. \url{https://arxiv.org/abs/2207.07886v5}
\bibitem{ref333} Scaling Transformer to 1M tokens and beyond with RMT. arXiv:2304.11062v2, 2023. \url{https://arxiv.org/abs/2304.11062v2}
\bibitem{ref334} General-Purpose In-Context Learning by Meta-Learning Transformers. arXiv:2212.04458v2, 2022. \url{https://arxiv.org/abs/2212.04458v2}
\bibitem{ref335} How does Multi-Task Training Affect Transformer In-Context Capabilities Investigations with Function Classes. arXiv:2404.03558v1, 2024. \url{https://arxiv.org/abs/2404.03558v1}
\bibitem{ref336} Exploring Parameter-Efficient Fine-Tuning Techniques for Code Generation with Large Language Models. arXiv:2308.10462v2, 2023. \url{https://arxiv.org/abs/2308.10462v2}
\bibitem{ref337} Beyond Fairness Alternative Moral Dimensions for Assessing Algorithms and Designing Systems. arXiv:2312.12559v1, 2023. \url{https://arxiv.org/abs/2312.12559v1}
\bibitem{ref338} Harms from Increasingly Agentic Algorithmic Systems. arXiv:2302.10329v2, 2023. \url{https://arxiv.org/abs/2302.10329v2}
\bibitem{ref339} What Are You Hiding Algorithmic Transparency and User Perceptions. arXiv:1812.03220v1, 2018. \url{https://arxiv.org/abs/1812.03220v1}
\bibitem{ref340} Logics and practices of transparency and opacity in real-world applications of public sector machine learning. arXiv:1706.09249v3, 2017. \url{https://arxiv.org/abs/1706.09249v3}
\bibitem{ref341} Machine Decisions and Human Consequences. arXiv:1811.06747v2, 2018. \url{https://arxiv.org/abs/1811.06747v2}
\bibitem{ref342} Towards an Understanding of Developers' Perceptions of Transparency in Software Development A Preliminary Study. arXiv:2309.06161v1, 2023. \url{https://arxiv.org/abs/2309.06161v1}
\bibitem{ref343} A Human Rights-Based Approach to Responsible AI. arXiv:2210.02667v1, 2022. \url{https://arxiv.org/abs/2210.02667v1}
\bibitem{ref344} Embedded EthiCS Integrating Ethics Broadly Across Computer Science Education. arXiv:1808.05686v1, 2018. \url{https://arxiv.org/abs/1808.05686v1}
\bibitem{ref345} Responsible Autonomy. arXiv:1706.02513v1, 2017. \url{https://arxiv.org/abs/1706.02513v1}
\bibitem{ref346} A Study on the Calibration of In-context Learning. arXiv:2312.04021v4, 2023. \url{https://arxiv.org/abs/2312.04021v4}
\bibitem{ref347} Beyond Bias and Compliance Towards Individual Agency and Plurality of Ethics in AI. arXiv:2302.12149v1, 2023. \url{https://arxiv.org/abs/2302.12149v1}
\bibitem{ref348} Why we need biased AI -- How including cognitive and ethical machine biases can enhance AI systems. arXiv:2203.09911v1, 2022. \url{https://arxiv.org/abs/2203.09911v1}
\bibitem{ref349} FAIRO Fairness-aware Adaptation in Sequential-Decision Making for Human-in-the-Loop Systems. arXiv:2307.05857v2, 2023. \url{https://arxiv.org/abs/2307.05857v2}
\bibitem{ref350} Breaking through the learning plateaus of in-context learning in Transformer. arXiv:2309.06054v2, 2023. \url{https://arxiv.org/abs/2309.06054v2}
\bibitem{ref351} Towards Equitable Agile Research and Development of AI and Robotics. arXiv:2402.08242v1, 2024. \url{https://arxiv.org/abs/2402.08242v1}
\bibitem{ref352} AI Fairness in Practice. arXiv:2403.14636v1, 2024. \url{https://arxiv.org/abs/2403.14636v1}
\bibitem{ref353} No computation without representation Avoiding data and algorithm biases through diversity. arXiv:2002.11836v1, 2020. \url{https://arxiv.org/abs/2002.11836v1}
\bibitem{ref354} Extending the Machine Learning Abstraction Boundary A Complex Systems Approach to Incorporate Societal Context. arXiv:2006.09663v1, 2020. \url{https://arxiv.org/abs/2006.09663v1}
\bibitem{ref355} Meta-learning the Learning Trends Shared Across Tasks. arXiv:2010.09291v1, 2020. \url{https://arxiv.org/abs/2010.09291v1}
\bibitem{ref356} Meta-Learning with Fewer Tasks through Task Interpolation. arXiv:2106.02695v2, 2021. \url{https://arxiv.org/abs/2106.02695v2}
\bibitem{ref357} MetaICL Learning to Learn In Context. arXiv:2110.15943v2, 2021. \url{https://arxiv.org/abs/2110.15943v2}
\bibitem{ref358} Measuring and Harnessing Transference in Multi-Task Learning. arXiv:2010.15413v3, 2020. \url{https://arxiv.org/abs/2010.15413v3}
\bibitem{ref359} Multi-Task Learning for Argumentation Mining. arXiv:1904.10162v1, 2019. \url{https://arxiv.org/abs/1904.10162v1}
\bibitem{ref360} Bridging Multi-Task Learning and Meta-Learning Towards Efficient Training and Effective Adaptation. arXiv:2106.09017v1, 2021. \url{https://arxiv.org/abs/2106.09017v1}
\bibitem{ref361} Diverse Distributions of Self-Supervised Tasks for Meta-Learning in NLP. arXiv:2111.01322v1, 2021. \url{https://arxiv.org/abs/2111.01322v1}
\bibitem{ref362} Reconciling meta-learning and continual learning with online mixtures of tasks. arXiv:1812.06080v3, 2018. \url{https://arxiv.org/abs/1812.06080v3}
\bibitem{ref363} Attention Is All You Need. arXiv:1706.03762v7, 2017. \url{https://arxiv.org/abs/1706.03762v7}
\bibitem{ref364} Structure-Regularized Attention for Deformable Object Representation. arXiv:2106.06672v1, 2021. \url{https://arxiv.org/abs/2106.06672v1}
\bibitem{ref365} COOL, a Context Outlooker, and its Application to Question Answering and other Natural Language Processing Tasks. arXiv:2204.09593v2, 2022. \url{https://arxiv.org/abs/2204.09593v2}
\bibitem{ref366} Memory Transformer. arXiv:2006.11527v2, 2020. \url{https://arxiv.org/abs/2006.11527v2}
\bibitem{ref367} All-in-One Image-Grounded Conversational Agents. arXiv:1912.12394v2, 2019. \url{https://arxiv.org/abs/1912.12394v2}
\bibitem{ref368} Zebra Extending Context Window with Layerwise Grouped Local-Global Attention. arXiv:2312.08618v1, 2023. \url{https://arxiv.org/abs/2312.08618v1}
\bibitem{ref369} Exploring Different Dimensions of Attention for Uncertainty Detection. arXiv:1612.06549v2, 2016. \url{https://arxiv.org/abs/1612.06549v2}
\bibitem{ref370} Translating Natural Language Instructions for Behavioral Robot Navigation with a Multi-Head Attention Mechanism. arXiv:2006.00697v3, 2020. \url{https://arxiv.org/abs/2006.00697v3}
\bibitem{ref371} Constrained Convolutional-Recurrent Networks to Improve Speech Quality with Low Impact on Recognition Accuracy. arXiv:1802.05874v1, 2018. \url{https://arxiv.org/abs/1802.05874v1}
\bibitem{ref372} FLAVA A Foundational Language And Vision Alignment Model. arXiv:2112.04482v3, 2021. \url{https://arxiv.org/abs/2112.04482v3}
\bibitem{ref373} MMICL Empowering Vision-language Model with Multi-Modal In-Context Learning. arXiv:2309.07915v3, 2023. \url{https://arxiv.org/abs/2309.07915v3}
\bibitem{ref374} Thinking Fast and Slow in AI. arXiv:2010.06002v2, 2020. \url{https://arxiv.org/abs/2010.06002v2}
\bibitem{ref375} Artificial Intelligence and its Role in Near Future. arXiv:1804.01396v1, 2018. \url{https://arxiv.org/abs/1804.01396v1}
\bibitem{ref376} Automating Ambiguity Challenges and Pitfalls of Artificial Intelligence. arXiv:2206.04179v1, 2022. \url{https://arxiv.org/abs/2206.04179v1}
\bibitem{ref377} The Interplay of Learning, Analytics, and Artificial Intelligence in Education. arXiv:2403.16081v2, 2024. \url{https://arxiv.org/abs/2403.16081v2}
\bibitem{ref378} What Artificial Neural Networks Can Tell Us About Human Language Acquisition. arXiv:2208.07998v2, 2022. \url{https://arxiv.org/abs/2208.07998v2}
\bibitem{ref379} Extending Machine Language Models toward Human-Level Language Understanding. arXiv:1912.05877v2, 2019. \url{https://arxiv.org/abs/1912.05877v2}
\bibitem{ref380} Beyond Interpretable Benchmarks Contextual Learning through Cognitive and Multimodal Perception. arXiv:2304.00002v2, 2022. \url{https://arxiv.org/abs/2304.00002v2}
\bibitem{ref381} Intelligent behavior depends on the ecological niche Scaling up AI to human-like intelligence in socio-cultural environments. arXiv:2103.06769v1, 2021. \url{https://arxiv.org/abs/2103.06769v1}
\bibitem{ref382} Lifelong Personal Context Recognition. arXiv:2205.10123v1, 2022. \url{https://arxiv.org/abs/2205.10123v1}
\bibitem{ref383} Learning to Complement Humans. arXiv:2005.00582v1, 2020. \url{https://arxiv.org/abs/2005.00582v1}
\end{thebibliography}

\end{document}
